% Dosages acido-basiques et solutions tampons — Chimie Tle TI — Cours signé OnBuch+
% Programme officiel MINESEC, Terminale TI. Compiler avec : tectonic L08-dosages-acido-basiques-solutions-tampons.tex
\newcommand{\DOCMATIERE}{Chimie}
\newcommand{\DOCNIVEAU}{Tle TI}
\newcommand{\DOCMODULE}{Module 2 · Acides et bases}
\newcommand{\DOCLECON}{08}
\newcommand{\DOCTITRE}{Dosages acido-basiques et solutions tampons}
% OnBuch+ — préambule « cours signés » ENRICHI (compilé avec tectonic / XeLaTeX)
% Système visuel « Pop » d'OnBuch+ : fond crème (jamais blanc), bordures encre
% épaisses, ombres « dures » décalées, titres Archivo Black en capitales.
% Version enrichie pour les cours de chimie : chimie (mhchem, chemfig),
% graphes (pgfplots), boîtes pédagogiques supplémentaires (définition,
% propriété, méthode, attention, le savais-tu, exercice résolu, exercice,
% TP, bilan), page de garde refaite et signature OnBuch+ ancrée en bas.
%
% Macros attendues dans main.tex AVANT \input{preamble} :
%   \DOCMATIERE  \DOCNIVEAU  \DOCMODULE  \DOCLECON (numéro)  \DOCTITRE
\documentclass[11pt]{article}

\usepackage{fontspec}
\usepackage[french]{babel}
\usepackage[a4paper,margin=2cm,top=3.4cm,bottom=2.8cm,headheight=1.4cm,footskip=1.3cm]{geometry}
\usepackage{amsmath,amssymb,amsfonts}
\usepackage[table]{xcolor} % [table] : fournit aussi \rowcolors (alternance de lignes)
\usepackage{pagecolor}
\usepackage{tikz}
\usetikzlibrary{arrows.meta,positioning,calc,shapes.geometric,shapes.misc,shapes.arrows,decorations.pathmorphing,decorations.markings,patterns,shadows,mindmap,trees,fit,backgrounds,matrix,intersections}
\usepackage{pgfplots}
\pgfplotsset{compat=1.18}
\usepackage[version=4]{mhchem}
\usepackage{chemfig}
\usepackage{enumitem}
\usepackage{fancyhdr}
% [most] charge toutes les bibliothèques tcolorbox (listings, minted, raster,
% documentation...) et gonfle inutilement la mémoire TeX sur un document long
% (~300 boîtes) : on ne charge que ce qui est réellement utilisé.
\usepackage[skins,breakable]{tcolorbox}
\usepackage{graphicx}
\usepackage{colortbl}
\usepackage{booktabs}
\usepackage{tabularx}
\usepackage{array}
\usepackage{multicol}
\usepackage{siunitx}
% parse-numbers=false : accepte aussi \SI{2,5 \times 10^{-3}}{...}
\sisetup{locale=FR,output-decimal-marker={,},group-minimum-digits=4,parse-numbers=false}
\DeclareSIUnit\ohm{\ensuremath{\symup{\Omega}}} % U+2126 absent de la police
\usepackage{qrcode}
% \degree : commande fréquemment utilisée par les modèles pour un angle ou une
% température en dehors de \SI{}{} (non fournie par les paquets chargés).
\providecommand{\degree}{^{\circ}}
\usepackage{lastpage}
\usepackage{hyperref}
\hypersetup{hidelinks}

% ============================================================
% Palette « Pop » OnBuch+ — mêmes hex que la classe P (plus_ui.dart)
% ============================================================
\definecolor{poporange}{HTML}{F59321}
\definecolor{popdark}{HTML}{E07A0C}
\definecolor{popcream}{HTML}{FFF6E8}
\definecolor{popink}{HTML}{1C1714}
\definecolor{poppurple}{HTML}{7A5AE0}
\definecolor{popgreen}{HTML}{2E9E6A}
\definecolor{poppink}{HTML}{E0457B}
\definecolor{popblue}{HTML}{2F7DE1}
\definecolor{popgold}{HTML}{C98A12}
\definecolor{popmuted}{HTML}{8A7F76}
\definecolor{popline}{HTML}{EADFD2}
\definecolor{popgray}{HTML}{8A7F76}
\colorlet{popgrey}{popgray}
% Alias tolérants (le modèle nomme parfois les couleurs différemment)
\colorlet{popwhite}{white}
\colorlet{popblack}{popink}
\colorlet{popred}{poppink}
\colorlet{popyellow}{popgold}
% Teintes claires pour fonds de tableaux / zones
\colorlet{popblueL}{popblue!10!white}
\colorlet{poporangeL}{poporange!14!white}
\colorlet{popgreenL}{popgreen!12!white}
\colorlet{poppurpleL}{poppurple!12!white}
\colorlet{poppinkL}{poppink!12!white}
\colorlet{popgoldL}{popgold!14!white}

\pagecolor{popcream}

% ============================================================
% Typographie
% ============================================================
\defaultfontfeatures{Ligatures=TeX}
\setmainfont{PlusJakartaSans-Medium.ttf}[Path=../../../fonts/,BoldFont=PlusJakartaSans-Bold.ttf]
% Maths en sans-serif (Fira Math) avec chiffres et lettres droites de Plus
% Jakarta, pour que les formules se fondent dans le texte.
\usepackage[math-style=ISO,bold-style=ISO]{unicode-math}
\setmathfont{FiraMath-Regular.otf}[Path=../../../fonts/]
\setmathfont{PlusJakartaSans-Medium.ttf}[Path=../../../fonts/,range={up/num,up/{Latin,latin},\mathup}]
\usepackage{icomma} % virgule décimale sans espace en mode math
% Caractères absents de la police de texte (sinon carré vide) : rendus par
% leur équivalent mathématique, en texte comme en formule.
\usepackage{newunicodechar}
\newunicodechar{α}{\ensuremath{\alpha}}
\newunicodechar{Α}{\ensuremath{\alpha}}
\newunicodechar{β}{\ensuremath{\beta}}
\newunicodechar{δ}{\ensuremath{\delta}}
\newunicodechar{✓}{\popcheck{popgreen}}
\newfontfamily\popdisplay{ArchivoBlack-Regular.ttf}[Path=../../../fonts/]
\newfontfamily\poplogo{SpaceGrotesk-Bold.ttf}[Path=../../../fonts/]
\newfontfamily\popsemi{PlusJakartaSans-SemiBold.ttf}[Path=../../../fonts/]
\newfontfamily\popxbold{PlusJakartaSans-ExtraBold.ttf}[Path=../../../fonts/]
\newfontfamily\popxboldls{PlusJakartaSans-ExtraBold.ttf}[Path=../../../fonts/,LetterSpace=3.0]

\setlist[enumerate]{leftmargin=1.7em,itemsep=5pt,topsep=4pt}
\setlist[itemize]{leftmargin=1.7em,itemsep=4pt,topsep=4pt,label={\color{poporange}\textbullet}}
\setlength{\parindent}{0pt}
\setlength{\parskip}{5pt}
\renewcommand{\arraystretch}{1.25}

% chemfig : liaisons un peu plus épaisses, encre Pop
\setchemfig{atom sep=2em,bond style={line width=0.9pt,popink},cram width=3pt}

% pgfplots : style Pop par défaut
\pgfplotsset{
  every axis/.append style={axis line style={popink,line width=1pt},tick style={popink},
    grid style={popline,line width=0.5pt},label style={font=\small},tick label style={font=\footnotesize}},
  popcurve/.style={poporange,line width=1.8pt,no markers,smooth},
}

% ============================================================
% Pill / squiggle
% ============================================================
\newcommand{\poppill}[3]{%
  \tikz[baseline=-0.55ex]\node[draw=popink,line width=1.2pt,rounded corners=2.6pt,%
    fill=#2,inner xsep=6.5pt,inner ysep=3pt]{%
    \popxboldls\fontsize{7}{7}\selectfont\color{#3}\MakeUppercase{#1}};%
}
\newcommand{\popsquiggle}[1]{%
  \tikz[baseline=-0.4ex]{
    \draw[#1,line width=2.6pt,line cap=round]
      plot [smooth] coordinates {(0,0.09) (0.34,0.01) (0.68,0.21) (1.02,0.01) (1.36,0.10)};
  }%
}
% Mot-clé surligné (marqueur orange sous le texte)
\newcommand{\cle}[1]{{\popxbold\color{popdark}#1}}

% ============================================================
% En-tête / pied
% ============================================================
\pagestyle{fancy}
\fancyhf{}
\renewcommand{\headrulewidth}{0pt}
\renewcommand{\footrulewidth}{0pt}
\lhead{\raisebox{-0.15ex}{\poplogo\fontsize{13}{13}\selectfont\color{popink}OnBuch\color{poporange}+}}
\rhead{\poppill{\DOCMATIERE\ \textbullet\ \DOCNIVEAU\ \textbullet\ Leçon \DOCLECON}{white}{popink}}
\lfoot{\popxbold\fontsize{7.6}{7.6}\selectfont\color{popmuted}\MakeUppercase{Cours signé OnBuch+}\ \textbullet\ {\popsemi\DOCTITRE}}
\rfoot{\tikz[baseline=-0.6ex]\node[rounded corners=8.5pt,fill=popink,minimum height=17pt,minimum width=17pt,inner xsep=5pt,inner sep=0]{\popxbold\fontsize{8}{8}\selectfont\color{popcream}\thepage\,/\,\pageref*{LastPage}};}

% ============================================================
% Titres
% ============================================================
\newcommand{\chaptitre}[2]{%
  \begin{tcolorbox}[enhanced,colback=poporange,colframe=popink,boxrule=2.6pt,arc=11pt,
      drop shadow=popink,left=15pt,right=15pt,top=12pt,bottom=11pt,before skip=3pt,after skip=18pt]
    \poppill{#2}{popink}{popcream}\par\vspace{9pt}
    {\raggedright\hyphenpenalty=10000\exhyphenpenalty=10000\popdisplay\fontsize{22}{24}\selectfont\color{popcream}\MakeUppercase{#1}\par}
  \end{tcolorbox}
}
% Section numérotée : pastille orange avec le numéro + titre Archivo
\newcounter{popsec}
\newcommand{\coursec}[1]{%
  \stepcounter{popsec}\setcounter{popsub}{0}%
  \par\vspace{16pt}\noindent
  \begin{minipage}{\linewidth}
  \tikz[baseline=-0.7ex]\node[draw=popink,line width=1.6pt,fill=poporange,rounded corners=4pt,minimum size=22pt,inner sep=0]{\popdisplay\fontsize{12}{12}\selectfont\color{popcream}\thepopsec};%
  \hspace{8pt}{\popdisplay\fontsize{15}{17}\selectfont\color{popink}\MakeUppercase{#1}}\par
  \vspace{4pt}\noindent{\color{poporange}\rule{36pt}{3.6pt}}
  \end{minipage}\par\nopagebreak\vspace{8pt}
}
\newcounter{popsub}
\newcommand{\courssub}[1]{%
  \stepcounter{popsub}%
  \par\vspace{9pt}\noindent{\popxbold\fontsize{11.5}{13}\selectfont\color{popink}{\color{poporange}\thepopsec.\thepopsub}\hspace{6pt}#1}\par\nopagebreak\vspace{4pt}
}

% ============================================================
% Boîtes pédagogiques — même recette : fond blanc, bordure épaisse colorée,
% ombre dure encre, badge plein. Titre optionnel [#1].
% ============================================================
\tcbset{popbase/.style={breakable,enhanced,colback=white,boxrule=2.3pt,arc=9pt,
  shadow={3pt}{-3pt}{0pt}{popink},left=14pt,right=14pt,top=11pt,bottom=11pt,before skip=11pt,after skip=15pt,
  fontupper=\popsemi\fontsize{10.6}{14.5}\selectfont\color{popink},
  fontlower=\popsemi\fontsize{10.6}{14.5}\selectfont\color{popink}}}
% Coche dessinée (le glyphe ✓ n'existe pas dans les polices du document)
\newcommand{\popcheck}[1]{\tikz[baseline=-0.5ex]\draw[#1,line width=1.6pt,line cap=round,line join=round](-0.13,0.0)--(-0.04,-0.09)--(0.13,0.1);}
\newcommand{\popboxhead}[3]{\poppill{#1}{#2}{white}\ifx\relax#3\relax\else\hspace{6pt}{\popxbold\fontsize{10.6}{12}\selectfont\color{popink}#3}\fi\par\vspace{7pt}}

\newtcolorbox{aretenir}[1][]{popbase,colframe=popgreen,before upper={\popboxhead{À retenir}{popgreen}{#1}}}
\newtcolorbox{exemplebox}[1][]{popbase,colframe=poporange,before upper={\popboxhead{Exemple}{poporange}{#1}}}
\newtcolorbox{definition}[1][]{popbase,colframe=popblue,before upper={\popboxhead{Définition}{popblue}{#1}}}
\newtcolorbox{propriete}[1][]{popbase,colframe=poppurple,before upper={\popboxhead{Propriété}{poppurple}{#1}}}
\newtcolorbox{methode}[1][]{popbase,colframe=popgold,before upper={\popboxhead{Méthode}{popgold}{#1}}}
\newtcolorbox{attention}[1][]{popbase,colframe=poppink,before upper={\popboxhead{Attention}{poppink}{#1}}}
\newtcolorbox{experience}[1][]{popbase,colframe=popblue,colback=popblueL,before upper={\popboxhead{Expérience}{popblue}{#1}}}
% « Le savais-tu ? » : carte violette pleine (contexte, culture, Cameroun)
\newtcolorbox{savaistu}[1][]{popbase,colback=poppurple,colframe=popink,
  fontupper=\popsemi\fontsize{10.4}{14.2}\selectfont\color{white},
  before upper={\poppill{Le savais-tu\,?}{white}{poppurple}\ifx\relax#1\relax\else\hspace{6pt}{\popxbold\color{white}#1}\fi\par\vspace{7pt}}}
% Exercice résolu : énoncé en haut, \tcblower puis la solution sur fond crème
\newcounter{popexo}\newcounter{popexr}
\newtcolorbox{exoresolu}[1][]{popbase,colframe=poporange,colbacklower=popcream,
  segmentation style={popink,line width=1.2pt,dashed},
  before upper={\stepcounter{popexr}\popboxhead{Exercice résolu \thepopexr}{poporange}{#1}},
  before lower={\poppill{Solution}{popink}{popcream}\par\vspace{6pt}}}
% Exercice (non résolu en place) — la correction est dans la partie Corrigés
\newtcolorbox{exercice}[1][]{popbase,colframe=popink,
  before upper={\stepcounter{popexo}\popboxhead{Exercice \thepopexo}{popink}{#1}}}
% Corrigé
\newtcolorbox{corrige}[1][]{popbase,colframe=popgreen,colback=popgreenL,
  before upper={\popboxhead{Corrigé}{popgreen}{#1}}}
% Objectifs / prérequis / bilan
\newtcolorbox{objectifs}{popbase,colframe=popink,colback=poporangeL,
  before upper={\popboxhead{Objectifs de la leçon}{popink}{}}}
\newtcolorbox{prerequis}{popbase,colframe=popink,colback=popblueL,
  before upper={\popboxhead{Prérequis}{popblue}{}}}
\newtcolorbox{bilan}{popbase,colframe=popink,colback=white,boxrule=2.8pt,
  before upper={\popboxhead{Fiche bilan}{poporange}{L'essentiel à connaître pour l'examen}}}
% Figure encadrée avec légende
\newtcolorbox{popfigure}[1][]{enhanced,breakable=false,colback=white,colframe=popline,boxrule=1.4pt,arc=8pt,
  left=10pt,right=10pt,top=10pt,bottom=8pt,before skip=10pt,after skip=12pt,halign=center,
  lower separated=false,halign lower=center,
  fontlower=\popsemi\fontsize{9}{11.5}\selectfont\color{popmuted},
  #1}
\newcommand{\legende}[1]{\par\vspace{6pt}{\popsemi\fontsize{9}{11.5}\selectfont\color{popmuted}#1}}

% ============================================================
% Signature OnBuch+
% ============================================================
\newcommand{\poplogoinline}[1]{{\poplogo\fontsize{#1}{#1}\selectfont\color{popink}OnBuch\color{poporange}+}}
\newcommand{\signatureonbuch}{%
  \begin{tcolorbox}[enhanced,colback=popink,colframe=popink,boxrule=2.2pt,arc=10pt,
      drop shadow=poporange,left=14pt,right=14pt,top=10pt,bottom=10pt,before skip=0pt,after skip=0pt]
    \tikz[baseline=-0.7ex]\node[circle,fill=popgreen,draw=popcream,line width=1.3pt,minimum size=22pt,inner sep=0]{\popcheck{white}};%
    \hspace{9pt}\begin{minipage}[c]{0.62\linewidth}
      {\popxbold\fontsize{12}{13}\selectfont\color{popcream}Signé\ \textbullet\ {\poplogo OnBuch\color{poporange}+}}\par\vspace{2pt}
      {\raggedright\popsemi\fontsize{8}{10.5}\selectfont\color{popline}\MakeUppercase{Équipe pédagogique OnBuch+}\\\MakeUppercase{Conforme au programme officiel MINESEC}\par}
    \end{minipage}\hfill
    {\poplogo\fontsize{22}{22}\selectfont\color{popcream}OnBuch\color{poporange}+}
  \end{tcolorbox}%
}

% ============================================================
% Page de garde
% #1 titre · #2 matière · #3 niveau · #4 module · #5 n° leçon · #6 durée ·
% #7 description / objectifs (texte ou itemize)
% ============================================================
\newcommand{\pagegarde}[7]{%
  \thispagestyle{empty}
  \newgeometry{a4paper,margin=1.7cm,top=1.6cm,bottom=1.4cm}%
  \begin{tikzpicture}[remember picture,overlay]
    \fill[poporange!18] ([xshift=-3.2cm,yshift=-3.6cm]current page.north east) circle (4.6cm);
    \fill[poppurple!14] ([xshift=2.4cm,yshift=7.6cm]current page.south west) circle (3.8cm);
  \end{tikzpicture}%
  {\poplogo\fontsize{30}{30}\selectfont\color{popink}OnBuch\color{poporange}+}\hfill
  \raisebox{0.6ex}{\poppill{Cours signé \textbullet\ Premium}{popink}{popcream}}\par
  \vspace{1pt}\popsquiggle{poppurple}\par
  \vspace{8pt}
  \begin{tcolorbox}[enhanced,colback=poporange,colframe=popink,boxrule=2.8pt,arc=13pt,
      drop shadow=popink,left=18pt,right=18pt,top=12pt,bottom=13pt,after skip=10pt]
    \poppill{#2 \textbullet\ #3}{popink}{popcream}\hspace{5pt}\poppill{#4}{white}{popink}\par\vspace{8pt}
    {\popdisplay\fontsize{14}{14}\selectfont\color{popink}LEÇON #5}\par\vspace{4pt}
    {\raggedright\hyphenpenalty=10000\exhyphenpenalty=10000\popdisplay\fontsize{25}{27}\selectfont\color{popcream}\MakeUppercase{#1}\par}
    \vspace{8pt}
    {\popxbold\fontsize{9.5}{9.5}\selectfont\color{popink}\MakeUppercase{Durée conseillée : #6}}
  \end{tcolorbox}
  \begin{tcolorbox}[enhanced,colback=white,colframe=popink,boxrule=2.2pt,arc=10pt,
      drop shadow=popink,left=16pt,right=16pt,top=10pt,bottom=10pt,after skip=8pt]
    \poppill{Dans cette leçon}{poppurple}{white}\par\vspace{5pt}
    \popsemi\fontsize{9.3}{12.6}\selectfont\color{popink}#7
  \end{tcolorbox}
  \noindent\begin{minipage}[t]{0.487\linewidth}
    \begin{tcolorbox}[enhanced,colback=popgreen,colframe=popink,boxrule=2.2pt,arc=10pt,
        drop shadow=popink,left=11pt,right=11pt,top=10pt,bottom=10pt,height=3.1cm,valign=center]
      \tcbox[colback=white,colframe=popink,boxrule=1.3pt,arc=5pt,boxsep=0pt,nobeforeafter,
        left=3pt,right=3pt,top=3pt,bottom=3pt,valign=center]{\qrcode[height=1.9cm]{https://play.google.com/store/apps/details?id=com.luvvix.onbuch&hl=fr}}%
      \hspace{8pt}\begin{minipage}[c]{0.5\linewidth}
        {\popdisplay\fontsize{11}{12.5}\selectfont\color{white}\MakeUppercase{Télécharge OnBuch}}\par\vspace{4pt}
        {\raggedright\popsemi\fontsize{8.4}{11}\selectfont\color{white}Gratuit \textbullet\ Google Play\\Tuteur IA, annales, concours, résultats\par}
      \end{minipage}
    \end{tcolorbox}
  \end{minipage}\hfill
  \begin{minipage}[t]{0.487\linewidth}
    \begin{tcolorbox}[enhanced,colback=poppurple,colframe=popink,boxrule=2.2pt,arc=10pt,
        drop shadow=popink,left=11pt,right=11pt,top=10pt,bottom=10pt,height=3.1cm,valign=center]
      \tikz[baseline=-0.7ex]\node[circle,fill=white,draw=popink,line width=1.3pt,minimum size=1.6cm,inner sep=0]{\popdisplay\fontsize{24}{24}\selectfont\color{poppurple}+};%
      \hspace{8pt}\begin{minipage}[c]{0.6\linewidth}
        {\popdisplay\fontsize{11}{12.5}\selectfont\color{white}\MakeUppercase{Passe à OnBuch+}}\par\vspace{4pt}
        {\raggedright\popsemi\fontsize{8.4}{11}\selectfont\color{white}Tous les cours signés, Tuteur IA illimité, fiches PDF — onglet OnBuch+ de l'app\par}
      \end{minipage}
    \end{tcolorbox}
  \end{minipage}
  \vspace{8pt}
  \popcontenu
  \vfill
  \signatureonbuch
  \restoregeometry
  \newpage
}

% Rangée « Ce cours contient » (page de garde)
\newcommand{\poptile}[3]{% #1 couleur #2 gros texte #3 légende
  \begin{tcolorbox}[enhanced,colback=#1,colframe=popink,boxrule=2pt,arc=9pt,shadow={2.5pt}{-2.5pt}{0pt}{popink},
      left=6pt,right=6pt,top=9pt,bottom=9pt,halign=center,valign=center,height=1.75cm,width=\linewidth,nobeforeafter]
    {\popdisplay\fontsize{12.5}{13}\selectfont\color{white}\MakeUppercase{#2}}\par\vspace{3pt}
    {\centering\popsemi\fontsize{7.8}{9.5}\selectfont\color{white}#3\par}
  \end{tcolorbox}}
\newcommand{\popcontenu}{%
  \par\noindent{\popxboldls\fontsize{7.5}{7.5}\selectfont\color{popmuted}CE COURS SIGNÉ CONTIENT}\par\vspace{7pt}
  \noindent\begin{minipage}[t]{0.235\linewidth}\poptile{popblue}{Cours}{complet, illustré, pas à pas}\end{minipage}\hfill
  \begin{minipage}[t]{0.235\linewidth}\poptile{poporange}{Méthodes}{et exercices résolus}\end{minipage}\hfill
  \begin{minipage}[t]{0.235\linewidth}\poptile{poppink}{Exercices}{type examen + corrigés}\end{minipage}\hfill
  \begin{minipage}[t]{0.235\linewidth}\poptile{popgreen}{Fiche bilan}{l'essentiel à retenir}\end{minipage}\par}

% Sommaire
\newtcolorbox{sommaire}{enhanced,colback=white,colframe=popink,boxrule=2.2pt,arc=10pt,
  drop shadow=popink,left=18pt,right=18pt,top=13pt,bottom=13pt,before skip=6pt,after skip=20pt,
  before upper={\poppill{Sommaire}{poppurple}{white}\par\vspace{9pt}\popsemi\fontsize{10.6}{14.5}\selectfont\color{popink}}}

% Signature de fin de cours — poussée tout en bas de la dernière page
\newcommand{\findecours}{%
  \par\vfill\nopagebreak
  \begin{center}\popsquiggle{poporange}\end{center}\vspace{4pt}
  \signatureonbuch
}
% Glyphes absents de Fira Math : redéfinis à partir de glyphes disponibles
\AtBeginDocument{%
  \def\setminus{\mathbin{\backslash}}\let\smallsetminus\setminus
  \def\vdots{\mathord{\vbox{\baselineskip3.2pt\lineskiplimit0pt\kern4pt\hbox{.}\hbox{.}\hbox{.}}}}%
  \def\ddots{\mathinner{\mkern1mu\raise6.5pt\hbox{.}\mkern2mu\raise3.6pt\hbox{.}\mkern2mu\raise0.7pt\hbox{.}\mkern1mu}}%
  \def\triangle{\mathord{\scalebox{0.9}{$\symup{\Delta}$}}}%
  \def\nabla{\mathord{\raisebox{\depth}{\scalebox{1}[-1]{$\symup{\Delta}$}}}}%
  \def\checkmark{\popcheck{popdark}}%
  \def\star{\mathbin{\ast}}\def\bigstar{\mathord{\ast}}%
  \def\diamond{\mathbin{\scalebox{0.6}{$\Diamond$}}}%
  \def\bigcup{\mathop{\mathchoice{\raisebox{-0.3em}{\scalebox{1.7}{$\cup$}}}{\scalebox{1.25}{$\cup$}}{\cup}{\cup}}\displaylimits}%
  \def\bigcap{\mathop{\mathchoice{\raisebox{-0.3em}{\scalebox{1.7}{$\cap$}}}{\scalebox{1.25}{$\cap$}}{\cap}{\cap}}\displaylimits}%
  \def\lesseqqgtr{\mathrel{\lessgtr}}\def\bigoplus{\mathop{\oplus}\displaylimits}%
}
\providecommand{\ssi}{si et seulement si} % abréviation fréquente

%% FIGFIX-BEGIN v4 — correctifs globaux des figures (docs/onbuch-generation/tools/figfix)
\usepackage{adjustbox}\usepackage{etoolbox}
% toute figure TikZ/pgfplots plus large que la ligne est réduite à la largeur disponible
\BeforeBeginEnvironment{tikzpicture}{\begin{adjustbox}{max width=\linewidth}}
\AfterEndEnvironment{tikzpicture}{\end{adjustbox}}
% légendes pgfplots lisibles par-dessus les courbes
\pgfplotsset{every axis/.append style={legend style={fill=white,fill opacity=0.92,text opacity=1,draw=black!15,inner sep=3pt}}}
% étiquettes d'arêtes (syntaxe edge["..."]) avec fond blanc
\tikzset{every edge quotes/.append style={fill=white,inner sep=1.2pt}}
%% FIGFIX-END
\begin{document}

\pagegarde{Dosages acido-basiques et solutions tampons}{Chimie}{Tle TI}{Module 2 · Acides et bases}{08}{4 h (cours + TP)}{Cette leçon étudie les réactions de dosage acido-basique (acide fort/base forte, acide faible/base forte et réciproquement), l'exploitation des courbes pH = f(V) et le choix des indicateurs colorés. Elle introduit enfin les solutions tampons, leur préparation à pH = pKa et leur importance en biologie et en industrie.
\vspace{4pt}
\begin{multicols}{2}\small
\begin{itemize}
\item À la fin de cette leçon, je sais écrire l'équation-bilan de la réaction de dosage entre un acide et une base (forts ou faibles).
\item À la fin de cette leçon, je sais définir l'équivalence acido-basique et exploiter la relation entre les quantités de matière à l'équivalence.
\item À la fin de cette leçon, je sais tracer et interpréter la courbe pH = f(V) d'un dosage pH-métrique.
\item À la fin de cette leçon, je sais déterminer graphiquement le point équivalent par la méthode des tangentes et le point de demi-équivalence.
\item À la fin de cette leçon, je sais calculer une concentration inconnue à partir du volume équivalent.
\item À la fin de cette leçon, je sais choisir judicieusement un indicateur coloré adapté à un dosage donné.
\end{itemize}\end{multicols}}

\begin{sommaire}
\begin{multicols}{2}
\begin{enumerate}
\item Réactions de dosage et équivalence acido-basique
\item Dosage pH-métrique d'un acide fort par une base forte
\item Dosage d'un acide faible par une base forte (et réciproquement)
\item Dosage colorimétrique et choix de l'indicateur coloré
\item Solutions tampons : définition, propriétés et préparation
\item Travaux pratiques
\item Méthodes et exercices résolus
\item Exercices
\item Corrigés des exercices
\item Fiche bilan
\end{enumerate}
\end{multicols}
\end{sommaire}

\begin{objectifs}
\begin{itemize}
\item À la fin de cette leçon, je sais écrire l'équation-bilan de la réaction de dosage entre un acide et une base (forts ou faibles).
\item À la fin de cette leçon, je sais définir l'équivalence acido-basique et exploiter la relation entre les quantités de matière à l'équivalence.
\item À la fin de cette leçon, je sais tracer et interpréter la courbe pH = f(V) d'un dosage pH-métrique.
\item À la fin de cette leçon, je sais déterminer graphiquement le point équivalent par la méthode des tangentes et le point de demi-équivalence.
\item À la fin de cette leçon, je sais calculer une concentration inconnue à partir du volume équivalent.
\item À la fin de cette leçon, je sais choisir judicieusement un indicateur coloré adapté à un dosage donné.
\item À la fin de cette leçon, je sais définir une solution tampon, citer ses propriétés et son importance.
\item À la fin de cette leçon, je sais préparer une solution tampon de pH = pKa.
\end{itemize}
\end{objectifs}

\begin{prerequis}
\begin{itemize}
\item Couples acide/base, Ka et pKa, force des acides et des bases (leçons précédentes du module 2)
\item pH d'une solution, relation pH = -log[H3O+] et produit ionique de l'eau Ke
\item Réactions totales et réactions limitées, avancement maximal et tableau d'avancement (Première)
\item Calculs de quantités de matière : n = C × V, dilutions (Première)
\item Lecture et exploitation de courbes expérimentales
\end{itemize}
\end{prerequis}

\newpage

\coursec{Réactions de dosage et équivalence acido-basique}

\courssub{Pourquoi doser ? Une situation concrète}

À la pharmacie du quartier à Yaoundé, un flacon de vinaigre commercial affiche « 8° d'acidité ». Comment vérifier cette information sans matériel sophistiqué ? Au laboratoire du lycée, on dispose d'une solution de soude (hydroxyde de sodium), d'un pH-mètre et de quelques indicateurs colorés. En versant progressivement la soude dans le vinaigre dilué et en suivant l'évolution du pH, on peut remonter à la concentration exacte en acide éthanoïque (acide acétique) du vinaigre. C'est le principe du \cle{dosage acido-basique}, technique incontournable du contrôle de qualité : eaux minérales, lait, médicaments, jus de bissap, bière... Et pourquoi le pH de notre sang reste-t-il stable autour de 7,4 malgré ce que nous mangeons ? La réponse viendra des \cle{solutions tampons}, étudiées à la fin de cette leçon.

\textbf{Problématique :} comment déterminer avec précision la concentration inconnue d'une solution acide ou basique, et quelles réactions chimiques rendent cette détermination possible ?

\courssub{La notion de dosage (titrage)}

\begin{definition}[Dosage ou titrage]
Un \cle{dosage} (ou \cle{titrage}) est une opération qui consiste à déterminer la \textbf{concentration inconnue} d'une espèce chimique (l'espèce \emph{titrée}) en la faisant réagir avec une espèce de concentration \textbf{exactement connue} (l'espèce \emph{titrante}), au moyen d'une réaction chimique.
\end{definition}

En dosage acido-basique, l'espèce titrée est un acide (ou une base) et l'espèce titrante est une base (ou un acide). La solution de concentration connue est appelée \cle{solution étalon} ; la détermination de sa concentration est parfois appelée \emph{étalonnage}.

\begin{propriete}[Conditions d'une réaction de dosage]
Pour servir de support à un dosage, une réaction chimique doit être :
\begin{itemize}
\item \textbf{totale} (constante d'équilibre $K$ très grande, en pratique $K > 10^4$) : les réactifs sont consommés intégralement ;
\item \textbf{rapide} : l'équilibre est atteint quasi instantanément à chaque ajout de réactif ;
\item \textbf{unique} : aucune réaction parasite ne doit consommer les réactifs.
\end{itemize}
On retient le triptyque : \textbf{T-R-U} (Totale, Rapide, Unique).
\end{propriete}

\begin{attention}[Solution titrante de qualité]
La solution titrante doit être de concentration \emph{précisément} connue : on la prépare par pesée d'un composé très pur (étalon primaire) ou on l'étalonne au préalable. La soude, par exemple, absorbe le dioxyde de carbone de l'air : une solution de soude ancienne n'est plus fiable et doit être réétalonnée avant tout dosage sérieux.
\end{attention}

\courssub{Le dispositif expérimental}

Le dispositif classique d'un dosage suivi au pH-mètre (dosage \cle{pH-métrique}) comprend une burette graduée contenant la solution titrante, un erlenmeyer (ou bécher) contenant un volume précis de solution à doser prélevé à la pipette jaugée, un agitateur magnétique pour homogénéiser le mélange, et une électrode de pH-mètre plongeant dans la solution.

\begin{popfigure}
\begin{center}
\begin{tikzpicture}[scale=0.9, every node/.style={font=\small}]
% Agitateur magnétique (socle)
\draw[fill=popmuted!40, rounded corners=3pt] (-2.6,-0.5) rectangle (2.6,-1.1);
\node at (0,-0.8) {\textbf{agitateur magnétique}};
% Barreau aimanté
\draw[fill=poporange, rounded corners=2pt] (-0.7,-0.32) rectangle (0.7,-0.12);
% Erlenmeyer
\draw[thick, popblue] (-0.55,2.1) -- (-0.55,1.55) -- (-1.7,-0.3) -- (1.7,-0.3) -- (0.55,1.55) -- (0.55,2.1);
% Liquide
\fill[popblueL] (-1.52,-0.28) -- (1.52,-0.28) -- (1.13,0.75) -- (-1.13,0.75) -- cycle;
\node[popdark] at (0,0.25) {solution à doser};
% Burette
\draw[thick, poppurple] (3.4,5.6) -- (3.4,1.4);
\draw[thick, poppurple] (3.9,5.6) -- (3.9,1.4);
\fill[poppurpleL] (3.4,5.55) rectangle (3.9,2.6);
% graduations
\foreach \y in {2.2,2.6,3.0,3.4,3.8,4.2,4.6,5.0,5.4} \draw[poppurple] (3.4,\y) -- (3.55,\y);
% robinet
\draw[thick, poppurple] (3.4,1.4) -- (3.65,1.15) -- (3.9,1.4);
\draw[thick, poppurple, fill=popgold] (3.25,1.15) rectangle (4.05,0.95);
\draw[thick, poppurple] (3.6,0.95) -- (3.6,0.6) -- (3.7,0.6) -- (3.7,0.95);
\node[poppurple, align=center] at (5.6,3.6) {burette graduée\\ (solution titrante)};
% goutte
\fill[poppurple] (3.65,0.1) circle (0.07);
\draw[->, poppurple, thick] (3.65,0.55) -- (3.65,0.2);
% pH-mètre : électrode
\draw[thick, popgreen!60!black, fill=popgreenL] (-1.15,3.4) rectangle (-0.85,0.4);
\draw[thick, popgreen!60!black] (-1.0,3.4) -- (-1.0,4.4) -- (-2.6,4.4);
% boîtier pH-mètre
\draw[thick, popgreen!60!black, fill=popcream, rounded corners=4pt] (-4.6,3.9) rectangle (-2.6,4.9);
\node[popdark] at (-3.6,4.6) {\textbf{pH-mètre}};
\node[popgreen!60!black] at (-3.6,4.2) {pH = 4,8};
\node[popgreen!60!black, align=center] at (-2.6,2.0) {électrode\\ de verre};
% annotations erlenmeyer
\node[popblue, align=center] at (-3.4,0.6) {erlenmeyer : volume $V_a$\\ prélevé à la pipette jaugée};
\draw[->, popblue] (-1.9,0.6) -- (-1.25,0.55);
\node[poporange!80!black] at (2.4,-0.2) {barreau aimanté};
\draw[->, poporange!80!black] (1.6,-0.22) -- (0.75,-0.22);
\end{tikzpicture}
\end{center}
\legende{Dispositif d'un dosage acido-basique suivi au pH-mètre : la burette délivre la solution titrante goutte à goutte ; l'agitation magnétique homogénéise le mélange ; le pH-mètre enregistre le pH après chaque ajout.}
\end{popfigure}

\begin{methode}[Réaliser un dosage en toute sécurité]
\begin{enumerate}
\item Rincer la burette avec la solution titrante, puis la remplir et ajuster le zéro.
\item Prélever à la \textbf{pipette jaugée} un volume exact $V_a$ de la solution à doser ; l'introduire dans l'erlenmeyer.
\item Ajouter éventuellement de l'eau distillée (sans influence sur la quantité de matière titrée) pour que l'électrode plonge correctement.
\item Placer le barreau aimanté, mettre en route l'agitation modérée, introduire l'électrode.
\item Verser la solution titrante par fractions décroissantes (mL, puis goutte à goutte près de l'équivalence) en notant le pH après chaque ajout.
\end{enumerate}
\end{methode}

\courssub{Les quatre grandes réactions de dosage}

\courssub{Acide fort et base forte : la réaction de référence}

Dans l'eau, un acide fort (comme \ce{HCl}) est intégralement transformé en ions oxonium \ce{H3O+}, et une base forte (comme la soude \ce{NaOH}) est intégralement dissociée en ions hydroxyde \ce{HO-}. Les ions sodium \ce{Na+} et chlorure \ce{Cl-} sont \emph{spectateurs} : ils ne participent pas à la réaction. La réaction de dosage s'écrit donc :

\[ \ce{H3O+ + HO- -> 2 H2O} \]

Cette réaction est \textbf{totale} (sa constante vaut $K = 10^{14}$, inverse du produit ionique de l'eau), \textbf{rapide} et \textbf{unique} : elle remplit parfaitement les trois conditions d'un dosage. C'est la réaction de référence des dosages acido-basiques.

\begin{exemplebox}[Dosage de l'acide chlorhydrique par la soude]
On dose \SI{20,0}{mL} d'acide chlorhydrique de concentration inconnue $C_a$ par de la soude à $C_b = \SI{0,10}{mol.L^{-1}}$. À chaque ajout de soude, les ions \ce{HO-} versés consomment instantanément une quantité égale d'ions \ce{H3O+}. Tant qu'il reste des ions \ce{H3O+}, le pH reste acide ; quand tout l'acide est consommé, la moindre goutte de soude en excès fait brutalement monter le pH.
\end{exemplebox}

\courssub{Acide faible et base forte}

Considérons un acide faible \ce{AH} (couple \ce{AH}/\ce{A-}, de constante d'acidité $K_a$) dosé par une base forte. Les ions \ce{HO-} arrachent le proton de l'acide faible :

\[ \ce{AH + HO- -> A- + H2O} \]

La constante de cette réaction vaut $K = \dfrac{K_a}{K_e} = 10^{14 - pK_a}$. Elle est donc \textbf{totale dès que $pK_a < 14$} (en pratique dès que $pK_a \lesssim 10$), ce qui est le cas de tous les acides faibles usuels : acide éthanoïque ($pK_a = 4,8$, $K \approx 10^{9,2}$), acides carboxyliques, ion ammonium ($pK_a = 9,2$)...

\begin{exemplebox}[Dosage du vinaigre]
Le vinaigre contient de l'acide éthanoïque \ce{CH3COOH}. Sa réaction avec la soude s'écrit :
\[ \ce{CH3COOH + HO- -> CH3COO- + H2O} \quad K = 10^{14-4{,}8} \approx 1{,}6 \times 10^{9} \]
Réaction totale, rapide, unique : le dosage est possible. C'est exactement le dosage de la situation d'accroche !
\end{exemplebox}

\courssub{Acide fort et base faible}

De manière symétrique, une base faible \ce{B} (couple \ce{BH+}/\ce{B}) peut être dosée par un acide fort, c'est-à-dire par les ions oxonium :

\[ \ce{B + H3O+ -> BH+ + H2O} \]

La constante de cette réaction vaut $K = \dfrac{1}{K_a} = 10^{pK_a}$ où $pK_a$ est celui du couple \ce{BH+}/\ce{B}. Elle est totale dès que $pK_a > 0$, ce qui est toujours vérifié pour les bases faibles usuelles. Par exemple, pour l'ammoniac \ce{NH3} ($pK_a(\ce{NH4+}/\ce{NH3}) = 9,2$) dosé par l'acide chlorhydrique :

\[ \ce{NH3 + H3O+ -> NH4+ + H2O} \quad K = 10^{9{,}2} \]

\courssub{Acide faible et base faible : le cas défavorable}

Et si l'on dosait un acide faible \ce{AH} par une base faible \ce{B} ? La réaction serait :

\[ \ce{AH + B <=> A- + BH+} \quad K = 10^{pK_a(\ce{BH+}/\ce{B}) - pK_a(\ce{AH}/\ce{A-})} \]

Pour que la réaction soit totale, il faudrait que le $pK_a$ du couple de la base dépasse largement celui du couple de l'acide, situation rare et jamais franche. En pratique, la réaction est \textbf{peu avancée} et l'équilibre s'établit \emph{avant} que les réactifs ne soient consommés.

\begin{savaistu}[Pourquoi ce dosage est-il inutilisable ?]
Même lorsque la constante $K$ est acceptable, le dosage acide faible / base faible est \textbf{déconseillé en pratique} : la courbe $pH = f(V)$ ne présente \textbf{aucun saut de pH marqué} à l'équivalence. Le pH varie lentement et continûment, si bien qu'aucun indicateur coloré ne vire nettement et qu'aucune méthode graphique ne permet de repérer le point équivalent avec précision. Au laboratoire, on contourne toujours le problème : pour doser une base faible, on utilise un \emph{acide fort} ; pour doser un acide faible, une \emph{base forte}.
\end{savaistu}

\begin{center}
\begin{tabularx}{\linewidth}{|l|X|X|c|c|}
\hline
\rowcolor{popblueL} \textbf{Réactifs} & \textbf{Équation-bilan} & \textbf{Constante $K$} & \textbf{Totale ?} & \textbf{Dosage ?} \\
\hline
Acide fort + base forte & \ce{H3O+ + HO- -> 2 H2O} & $10^{14}$ & Oui & Excellent \\
\hline
Acide faible + base forte & \ce{AH + HO- -> A- + H2O} & $10^{14 - pK_a}$ & Oui si $pK_a < 14$ & Bon \\
\hline
Acide fort + base faible & \ce{B + H3O+ -> BH+ + H2O} & $10^{pK_a}$ & Oui si $pK_a > 0$ & Bon \\
\hline
Acide faible + base faible & \ce{AH + B <=> A- + BH+} & $10^{\Delta pK_a}$ & Rarement & Déconseillé \\
\hline
\end{tabularx}
\end{center}

\begin{methode}[Écrire l'équation-bilan d'une réaction de dosage]
\begin{enumerate}
\item \textbf{Identifier} la nature de chaque réactif : acide fort (\ce{H3O+}), base forte (\ce{HO-}), acide faible \ce{AH} ou base faible \ce{B}.
\item \textbf{Repérer les couples} mis en jeu : l'acide d'un couple réagit avec la base de l'autre couple.
\item \textbf{Écrire le transfert de proton} : le donneur de \ce{H+} (l'acide) cède son proton à l'accepteur (la base). Éliminer les ions spectateurs (\ce{Na+}, \ce{Cl-}...).
\item \textbf{Vérifier le caractère total} en calculant $K$ à partir des $pK_a$ : si $K > 10^4$, la réaction convient pour un dosage.
\end{enumerate}
\end{methode}

\courssub{L'équivalence acido-basique}

Revenons à notre erlenmeyer. Au début, il contient $n_a = C_a \times V_a$ moles d'acide. On verse la base goutte à goutte : chaque mole de base versée détruit une mole d'acide. Il existe un \textbf{instant précis} où la quantité de base versée est \emph{exactement} égale à la quantité d'acide initialement présente : c'est \cle{l'équivalence}.

\begin{definition}[Équivalence acido-basique]
L'\cle{équivalence} est atteinte lorsque les réactifs ont été mélangés dans les \textbf{proportions stœchiométriques} de la réaction de dosage : l'acide et la base ont alors été \textbf{totalement consommés}. Pour un monoacide et une monobase (coefficients 1 et 1) :
\[ n(\text{acide apporté}) = n(\text{base apportée}) \quad \Longleftrightarrow \quad \boxed{C_a \times V_a = C_b \times V_{bE}} \]
où $V_{bE}$ est le volume de solution titrante versé à l'équivalence.
\end{definition}

\begin{attention}[Équivalence $\neq$ neutralité !]
Erreur classique du baccalauréat : confondre \textbf{équivalence} et \textbf{neutralité}. L'équivalence est une \emph{relation entre quantités de matière} ($C_a V_a = C_b V_{bE}$) ; la neutralité est une \emph{valeur de pH} ($pH = 7$ à \SI{25}{\celsius}). Les deux ne coïncident que pour le dosage d'un acide fort par une base forte. À l'équivalence du dosage d'un acide faible par la soude, la solution contient la base faible \ce{A-} : le pH est \textbf{supérieur à 7}. Inversement, à l'équivalence du dosage de l'ammoniac par \ce{HCl}, le pH est \textbf{inférieur à 7} à cause de l'ion \ce{NH4+} formé. Retiens : \emph{on peut être à l'équivalence sans être neutre}.
\end{attention}

\courssub{Tableau d'avancement à l'équivalence}

Le tableau d'avancement rend la relation d'équivalence limpide. Pour le dosage d'un volume $V_a$ d'acide chlorhydrique ($n_a = C_a V_a$ moles de \ce{H3O+}) par un volume $V_b$ de soude ($n_b = C_b V_b$ moles de \ce{HO-}) :

\begin{center}
\begin{tabularx}{\linewidth}{|l|c|X|X|X|}
\hline
\rowcolor{popblueL} \textbf{État} & \textbf{Avancement} & \ce{H3O+} & \ce{HO-} & \ce{H2O} \\
\hline
Initial & $x = 0$ & $C_a V_a$ & $C_b V_b$ & solvant (excès) \\
\hline
Intermédiaire & $x$ & $C_a V_a - x$ & $C_b V_b - x$ & excès \\
\hline
Équivalence & $x_E = x_{\max}$ & $C_a V_a - x_E = 0$ & $C_b V_{bE} - x_E = 0$ & excès \\
\hline
\end{tabularx}
\end{center}

À l'équivalence, \textbf{les deux réactifs sont simultanément et totalement consommés} (réactifs en proportions stœchiométriques). Les deux lignes « $= 0$ » donnent $x_E = C_a V_a = C_b V_{bE}$ : on retrouve la relation fondamentale.

\begin{attention}[Avant et après l'équivalence]
\begin{itemize}
\item \textbf{Avant l'équivalence} ($V_b < V_{bE}$) : la base versée est le réactif limitant ; elle disparaît entièrement et il reste de l'acide en excès. Le pH est fixé par l'acide restant.
\item \textbf{Après l'équivalence} ($V_b > V_{bE}$) : l'acide a disparu ; c'est la base versée qui s'accumule et impose le pH.
\end{itemize}
L'équivalence est donc le \emph{moment de bascule} où l'espèce qui fixe le pH change : c'est ce basculement qui produit le saut de pH exploité dans les dosages.
\end{attention}

\begin{exoresolu}[Vérification d'une concentration d'acide chlorhydrique]
\textbf{Énoncé.} Au laboratoire du lycée, on prélève $V_a = \SI{20,0}{mL}$ d'une solution d'acide chlorhydrique de concentration $C_a$ inconnue, que l'on dose par une solution de soude étalonnée à $C_b = \SI{0,10}{mol.L^{-1}}$. Le suivi pH-métrique montre que le saut de pH se produit pour un volume versé $V_{bE} = \SI{15,0}{mL}$. Déterminer $C_a$.
\tcblower
\textbf{Solution détaillée.}

\textbf{Étape 1 : équation de la réaction de dosage.} Acide fort + base forte :
\[ \ce{H3O+ + HO- -> 2 H2O} \]
Réaction totale ($K = 10^{14}$), rapide, unique : le dosage est légitime.

\textbf{Étape 2 : relation à l'équivalence.} Les coefficients stœchiométriques valent 1 et 1, donc :
\[ C_a \times V_a = C_b \times V_{bE} \]

\textbf{Étape 3 : application numérique} (les volumes peuvent rester en mL puisqu'ils sont dans la même unité des deux côtés) :
\[ C_a = \frac{C_b \times V_{bE}}{V_a} = \frac{0{,}10 \times 15{,}0}{20{,}0} = \SI{0,075}{mol.L^{-1}} \]

\textbf{Conclusion :} la solution d'acide chlorhydrique a pour concentration $C_a = \SI{7,5e-2}{mol.L^{-1}}$.

\textbf{Vérification de cohérence :} le volume équivalent (\SI{15,0}{mL}) est inférieur au volume d'acide (\SI{20,0}{mL}), donc l'acide doit être \emph{moins concentré} que la soude : $0{,}075 < 0{,}10$. Le résultat est cohérent.
\end{exoresolu}

\begin{exercice}[Dosage du vinaigre — retour à la situation d'accroche]
On dilue 10 fois un vinaigre commercial, puis on dose $V_a = \SI{20,0}{mL}$ de la solution diluée par de la soude à $C_b = \SI{0,100}{mol.L^{-1}}$. L'équivalence est obtenue pour $V_{bE} = \SI{13,3}{mL}$.
\begin{enumerate}
\item Écrire l'équation-bilan de la réaction de dosage et justifier qu'elle convient.
\item Calculer la concentration en acide éthanoïque de la solution diluée, puis du vinaigre commercial.
\item Sachant que le degré d'acidité est la masse d'acide éthanoïque ($M = \SI{60}{g.mol^{-1}}$) en grammes pour \SI{100}{g} de vinaigre (densité $\approx 1$), vérifier l'indication « 8° » de l'étiquette.
\end{enumerate}
\end{exercice}

\begin{corrige}[Exercice : dosage du vinaigre]
\textbf{1.} Acide faible + base forte : $\ce{CH3COOH + HO- -> CH3COO- + H2O}$, avec $K = 10^{14-4{,}8} \approx 1{,}6\times10^{9} \gg 10^4$ : réaction totale, rapide, unique, donc adaptée au dosage.

\textbf{2.} À l'équivalence : $C_{a,\text{dil}} = \dfrac{C_b V_{bE}}{V_a} = \dfrac{0{,}100 \times 13{,}3}{20{,}0} = \SI{6,65e-2}{mol.L^{-1}}$. Le vinaigre commercial est 10 fois plus concentré : $C_a = \SI{0,665}{mol.L^{-1}}$.

\textbf{3.} Masse d'acide par litre : $m = C_a \times M = 0{,}665 \times 60 \approx \SI{40}{g.L^{-1}}$, soit environ \SI{4,0}{g} pour \SI{100}{mL} $ \approx $ \SI{100}{g} de vinaigre : le degré mesuré est $ \approx 4° $, et non 8°. L'étiquette du flacon est donc \emph{fausse} (ou le vinaigre s'est dégradé) : le dosage a permis le contrôle qualité annoncé dans la situation d'accroche.
\end{corrige}

\begin{aretenir}[L'essentiel de la section]
\begin{itemize}
\item Un \cle{dosage} détermine une concentration inconnue grâce à une réaction \textbf{totale, rapide et unique}.
\item Réactions de dosage : \ce{H3O+ + HO- -> 2 H2O} (fort/fort), \ce{AH + HO- -> A- + H2O} (acide faible/base forte), \ce{B + H3O+ -> BH+ + H2O} (base faible/acide fort). Le couple acide faible + base faible est à proscrire.
\item À \cle{l'équivalence}, les réactifs sont en proportions stœchiométriques : $C_a V_a = C_b V_{bE}$ (monoacide/monobase).
\item Équivalence $\neq$ neutralité : le pH à l'équivalence dépend des espèces formées.
\end{itemize}
\end{aretenir}

\coursec{Dosage pH-métrique d'un acide fort par une base forte}

Dans la section précédente, nous avons défini l'\cle{équivalence acido-basique} : c'est l'état du mélange où l'acide et la base ont été introduits dans les proportions stoechiométriques de la réaction de dosage. Toute la question est maintenant pratique : \emph{comment repérer expérimentalement cette équivalence ?} La première réponse, la plus précise, est le \cle{dosage pH-métrique} : on suit le pH du mélange au fur et à mesure de l'addition du réactif titrant. Nous allons l'étudier sur le cas le plus simple et le plus fondamental : le dosage d'un acide fort (l'acide chlorhydrique) par une base forte (l'hydroxyde de sodium, la soude).

\courssub{Dispositif expérimental et réaction de dosage}

\textbf{Le montage.}\par
Le dosage pH-métrique exige un matériel précis et un protocole rigoureux :

\begin{itemize}
\item une \cle{burette} graduée, fixée sur un support, contenant la solution titrante de concentration connue (ici la soude à $C_b$) ;
\item un \cle{bécher} contenant un volume $V_a$ exact de la solution à titrer (ici l'acide chlorhydrique), prélevé à la \emph{pipette jaugée} ;
\item un \cle{pH-mètre} préalablement \emph{étalonné} avec deux solutions étalons (pH $= 4$ et pH $= 7$), dont l'électrode de verre plonge dans le bécher ;
\item un \cle{agitateur magnétique} qui homogénéise le mélange en permanence : sans agitation, le pH mesuré près de l'électrode ne serait pas représentatif de l'ensemble de la solution.
\end{itemize}

\begin{experience}[Dosage pH-métrique de l'acide chlorhydrique par la soude]
\textbf{Protocole :} on introduit $V_a = \SI{10,0}{mL}$ d'acide chlorhydrique de concentration inconnue $C_a$ dans un bécher, on ajoute de l'eau distillée pour que l'électrode soit bien immergée, puis on verse la soude de concentration $C_b = \SI{0,050}{mol.L^{-1}}$ par fractions de \SI{1}{mL} (puis goutte à goutte au voisinage du saut de pH), en relevant le pH après chaque addition.

\textbf{Observations :} le pH démarre vers $1$, croît d'abord très lentement, puis augmente brutalement de plusieurs unités pour quelques gouttes versées, avant de se stabiliser vers $13$.

\textbf{Interprétation :} chaque ajout de soude consomme des ions \ce{H3O+} ; tant qu'il reste de l'acide en excès, le pH reste faible. Lorsque tout l'acide est consommé, la moindre goutte supplémentaire de soude fait basculer le milieu en zone basique : c'est le saut de pH.
\end{experience}

\textbf{La réaction de dosage.}\par
L'acide chlorhydrique est totalement dissocié : \ce{HCl + H2O -> H3O+ + Cl-}. La soude est totalement dissociée : \ce{NaOH -> Na+ + OH-}. La réaction qui se produit dans le bécher est donc la réaction entre les ions oxonium et les ions hydroxyde :

\[ \ce{H3O+ + OH- -> 2 H2O} \]

Cette réaction possède les trois qualités exigées d'une réaction de dosage : elle est \cle{totale} (constante d'équilibre $K = 10^{14}$, immense), \cle{rapide} (instantanée) et \cle{unique} (aucune réaction parasite). Les ions \ce{Na+} et \ce{Cl-} sont des ions \emph{spectateurs} : ils n'interviennent ni dans le bilan, ni dans le pH.

\begin{attention}[Bien étalonner le pH-mètre]
Un pH-mètre non étalonné donne des valeurs fausses de plusieurs dixièmes d'unité : le saut de pH sera décalé et la détermination de l'équivalence compromise. L'étalonnage avec deux solutions tampons étalons (pH $4$ et pH $7$) est \emph{obligatoire} avant toute mesure. Autre piège : oublier l'agitation, ce qui crée des gradients de concentration autour de l'électrode et des mesures instables.
\end{attention}

\courssub{Allure de la courbe pH = f(Vb) : les quatre zones}

En portant le pH mesuré en ordonnée et le volume $V_b$ de soude versé en abscisse, on obtient une courbe croissante en forme de « S » couché, appelée \cle{courbe de dosage}. Décortiquons ses quatre zones.

\begin{popfigure}
\begin{center}
\begin{tikzpicture}
\begin{axis}[
width=0.85\linewidth, height=7cm,
xlabel={$V_b$ (\si{mL})}, ylabel={pH},
xmin=0, xmax=25, ymin=0, ymax=14,
xtick={0,5,10,12.4,15,20,25}, ytick={0,2,4,6,7,8,10,12,14},
grid=both, grid style={popline!40},
axis lines=left]
\addplot[popcurve,domain=0:25,samples=300] {7 + 6*tanh(2.2*(x-12.4)/1.6)};
\addplot[popmuted, dashed] coordinates {(12.4,0) (12.4,7)};
\addplot[popmuted, dashed] coordinates {(0,7) (12.4,7)};
\addplot[only marks, mark=*, poporange, mark size=2pt] coordinates {(12.4,7)};
\node[poporange, anchor=west] at (axis cs:13,6.2) {$E\,(V_{bE}\,;\,\mathrm{pH}_E=7)$};
\node[popmuted, anchor=south] at (axis cs:12.4,0.2) {$V_{bE}$};
\node[popblue, anchor=south west] at (axis cs:0.5,1.3) {\small zone 1};
\node[popblue, anchor=south] at (axis cs:8,2.6) {\small zone 2};
\node[popblue] at (axis cs:14.5,10.5) {\small zone 3};
\node[popblue, anchor=west] at (axis cs:18,12.3) {\small zone 4};
\end{axis}
\end{tikzpicture}
\end{center}
\legende{Courbe pH $= f(V_b)$ du dosage de \SI{10,0}{mL} de \ce{HCl} par la soude à \SI{0,050}{mol.L^{-1}}. Le point équivalent $E$ est le point d'inflexion situé au milieu du saut de pH.}
\end{popfigure}

\begin{itemize}
\item \textbf{Zone 1 — le point de départ ($V_b = 0$).} Le bécher ne contient que l'acide fort : le pH initial est faible, donné par $\mathrm{pH} = -\log C_a$. Pour $C_a = \SI{6,2 \times 10^{-2}}{mol.L^{-1}}$, pH $\approx 1{,}2$.
\item \textbf{Zone 2 — la variation lente (avant l'équivalence).} La soude versée est entièrement consommée, mais il reste un excès d'ions \ce{H3O+} qui impose un pH acide. Le pH monte doucement : l'acide est encore majoritaire.
\item \textbf{Zone 3 — le saut de pH (autour de l'équivalence).} Au voisinage immédiat de l'équivalence, quelques gouttes de soude suffisent à faire passer le pH d'environ $4$ à environ $10$ : la courbe devient quasi verticale. C'est dans cette zone que se trouve le point équivalent.
\item \textbf{Zone 4 — le palier (après l'équivalence).} Tout l'acide est consommé ; la soude versée s'accumule et impose son pH. La courbe tend vers le pH de la soude diluée, autour de $13$.
\end{itemize}

\begin{definition}[Point équivalent E]
Le \cle{point équivalent} $E$ est le point de la courbe pH $= f(V_b)$ où le pH varie le plus rapidement : c'est le \emph{point d'inflexion} de la courbe, situé au cœur du saut de pH. Il est repéré par ses coordonnées $(V_{bE}\,;\,\mathrm{pH}_E)$. À ce point, les réactifs ont été mélangés dans les proportions stoechiométriques : $n_a(\ce{H3O+}) = n_b(\ce{OH-})$, soit $C_a V_a = C_b V_{bE}$.
\end{definition}

\begin{propriete}[pH à l'équivalence d'un dosage acide fort -- base forte]
À l'équivalence du dosage d'un acide fort par une base forte, le bécher ne contient que de l'eau et un sel d'acide fort et de base forte (ici \ce{NaCl}), dont les ions \ce{Na+} et \ce{Cl-} sont sans action sur l'eau. La solution est donc \emph{neutre} :
\[ \mathrm{pH}_E = 7{,}0 \quad \text{à } \SI{25}{\celsius} \]
Cette valeur est une signature du couple « acide fort / base forte » : nous verrons dans la section suivante qu'avec un acide faible, $\mathrm{pH}_E > 7$.
\end{propriete}

\courssub{Détermination graphique du point équivalent}

\textbf{Méthode des tangentes parallèles.}\par
Comment localiser précisément le point d'inflexion $E$ sur une courbe expérimentale tracée point par point ? La méthode classique, exigible au Baccalauréat TI, est celle des \cle{tangentes parallèles}.

\begin{methode}[Méthode des tangentes parallèles]
\begin{enumerate}
\item Tracer la tangente à la courbe dans la zone 2 (avant le saut), là où la courbe est régulière.
\item Tracer la tangente à la courbe dans la zone 4 (après le saut), \emph{parallèle} à la première.
\item Tracer une troisième droite parallèle aux deux précédentes et \emph{équidistante} des deux tangentes.
\item Cette droite médiane coupe la courbe au point équivalent $E$.
\item Projeter $E$ sur l'axe des abscisses pour lire $V_{bE}$, et sur l'axe des ordonnées pour lire $\mathrm{pH}_E$.
\end{enumerate}
\end{methode}

\begin{popfigure}
\begin{center}
\begin{tikzpicture}
\begin{axis}[
width=0.85\linewidth, height=7.5cm,
xlabel={$V_b$ (\si{mL})}, ylabel={pH},
xmin=8, xmax=17, ymin=2, ymax=12,
xtick={8,10,12,12.4,14,16}, ytick={2,4,6,7,8,10,12},
grid=both, grid style={popline!40},
axis lines=left]
\addplot[popcurve,domain=8:17,samples=300] {7 + 6*tanh(2.2*(x-12.4)/1.6)};
\addplot[popgreen, thick, domain=9.6:12.2] {2.55 + 0.85*(x-10.5)};
\addplot[popgreen, thick, domain=12.8:15.6] {11.45 + 0.85*(x-14.5)};
\addplot[poporange, thick, dashed, domain=10.6:14.6] {7 + 0.85*(x-12.4)};
\addplot[only marks, mark=*, poporange, mark size=2.5pt] coordinates {(12.4,7)};
\addplot[popmuted, dashed] coordinates {(12.4,2) (12.4,7)};
\node[poporange, anchor=west] at (axis cs:12.7,6.3) {$E$};
\node[popmuted, anchor=north] at (axis cs:12.4,2) {$V_{bE}$};
\node[popgreen, anchor=south east, font=\small] at (axis cs:10.2,3.1) {tangente 1};
\node[popgreen, anchor=north west, font=\small] at (axis cs:15.2,11.9) {tangente 2};
\node[poporange, font=\small, anchor=south west] at (axis cs:13.9,8.4) {droite médiane};
\end{axis}
\end{tikzpicture}
\end{center}
\legende{Construction des tangentes parallèles : les deux tangentes (vertes) sont tracées de part et d'autre du saut ; la droite médiane (orange, pointillés) coupe la courbe au point équivalent $E$.}
\end{popfigure}

\begin{attention}[Erreur fréquente au Baccalauréat]
Beaucoup d'élèves lisent $V_{bE}$ « au milieu du saut à vue d'œil », sans construction géométrique. C'est une erreur : le milieu visuel du segment vertical ne coïncide pas toujours avec le point d'inflexion, surtout si la courbe est dissymétrique ou mal espacée en points. Le correcteur attend \emph{explicitement} les deux tangentes parallèles et la droite médiane tracées sur la courbe. Une lecture sans construction perd des points, même si la valeur numérique est correcte.
\end{attention}

\textbf{Méthode de la courbe dérivée.}\par
Une seconde méthode, très appréciée comme complément au Baccalauréat, consiste à tracer la \cle{courbe dérivée} $\dfrac{\mathrm{d}\mathrm{pH}}{\mathrm{d}V_b} = f(V_b)$. Or le point équivalent est précisément le point où le pH varie le plus vite : la dérivée y est donc \emph{maximale}.

\begin{propriete}[Exploitation de la courbe dérivée]
La courbe dérivée $\dfrac{\mathrm{d}\mathrm{pH}}{\mathrm{d}V_b}$ présente un \emph{pic} (maximum très étroit) dont l'abscisse est exactement $V_{bE}$. En pratique, on calcule pour chaque intervalle entre deux mesures successives :
\[ \left(\frac{\mathrm{d}\mathrm{pH}}{\mathrm{d}V_b}\right)_i \approx \frac{\mathrm{pH}_{i+1}-\mathrm{pH}_i}{V_{b,i+1}-V_{b,i}} \]
et l'on porte ce rapport en fonction de $V_b$. Le sommet du pic donne $V_{bE}$ sans construction de tangentes.
\end{propriete}

\begin{popfigure}
\begin{center}
\begin{tikzpicture}
\begin{axis}[
width=0.85\linewidth, height=6cm,
xlabel={$V_b$ (\si{mL})}, ylabel={$\dfrac{\mathrm{d}\mathrm{pH}}{\mathrm{d}V_b}$ (\si{pH.mL^{-1}})},
xmin=8, xmax=17, ymin=0, ymax=9,
xtick={8,10,12,12.4,14,16},
grid=both, grid style={popline!40},
axis lines=left]
\addplot[popcurve,domain=8:17,samples=300] {8.25*(1-(tanh(2.2*(x-12.4)/1.6))^2)};
\addplot[popmuted, dashed] coordinates {(12.4,0) (12.4,8.25)};
\addplot[only marks, mark=*, poporange, mark size=2pt] coordinates {(12.4,8.25)};
\node[poporange, anchor=south west] at (axis cs:12.6,8.1) {maximum : $V_{bE} = \SI{12,4}{mL}$};
\end{axis}
\end{tikzpicture}
\end{center}
\legende{Courbe dérivée du dosage : le pic, très marqué, a pour abscisse le volume équivalent $V_{bE}$.}
\end{popfigure}

\begin{savaistu}[Les pH-mètres qui trouvent l'équivalence tout seuls]
Les pH-mètres de laboratoire modernes, couplés à une burette automatique, calculent en temps réel la dérivée $\mathrm{d}\mathrm{pH}/\mathrm{d}V_b$ et détectent \emph{automatiquement} le maximum du pic : l'appareil affiche directement $V_{bE}$ et arrête même le dosage. C'est le principe des titrateurs automatiques utilisés dans les laboratoires de contrôle qualité, par exemple à la SONARA pour vérifier l'acidité des produits pétroliers, ou dans les brasseries pour contrôler l'acidité des moûts. La méthode des tangentes reste néanmoins exigible à l'examen : elle prouve que tu comprends ce que la machine calcule !
\end{savaistu}

\courssub{Exploitation quantitative : calcul de la concentration inconnue}

Tout l'enjeu du dosage est là : connaissant $V_{bE}$, on remonte à la concentration inconnue $C_a$ de l'acide grâce à la relation d'équivalence établie dans la section précédente.

\begin{aretenir}[Relation d'équivalence acide fort -- base forte]
À l'équivalence du dosage d'un acide fort (monoacide) par une base forte (monobase) :
\[ C_a \times V_a = C_b \times V_{bE} \]
avec $C_a$ et $C_b$ en \si{mol.L^{-1}}, $V_a$ et $V_{bE}$ dans la \emph{même} unité de volume (les deux en \si{mL} conviennent).
\end{aretenir}

\begin{exoresolu}[Vérification d'une étiquette d'acide chlorhydrique]
\textbf{Énoncé.} Un laborantin d'un lycée technique de Douala veut vérifier la concentration d'une bouteille d'acide chlorhydrique. Il prélève $V_a = \SI{10,0}{mL}$ de la solution, qu'il dose par pH-métrie avec de la soude à $C_b = \SI{0,050}{mol.L^{-1}}$. La méthode des tangentes donne $V_{bE} = \SI{12,4}{mL}$ et $\mathrm{pH}_E = 7{,}0$. Déterminer la concentration $C_a$ de l'acide.
\tcblower
\textbf{Solution détaillée.}

\textbf{Étape 1 — Équation de la réaction de dosage :}
\[ \ce{H3O+ + OH- -> 2 H2O} \]
La réaction est totale ($K = 10^{14}$), rapide et unique : c'est bien une réaction de dosage.

\textbf{Étape 2 — Relation à l'équivalence :} à l'équivalence, les réactifs ont été introduits en proportions stoechiométriques :
\[ n(\ce{H3O+}) = n(\ce{OH-}) \quad\Longrightarrow\quad C_a \times V_a = C_b \times V_{bE} \]

\textbf{Étape 3 — Application numérique :}
\[ C_a = \frac{C_b \times V_{bE}}{V_a} = \frac{0{,}050 \times 12{,}4}{10{,}0} = \SI{6,2 \times 10^{-2}}{mol.L^{-1}} \]

\textbf{Étape 4 — Vérification de cohérence :} le pH initial mesuré était d'environ $1{,}2$ ; or $-\log(6{,}2\times 10^{-2}) \approx 1{,}21$ : parfaitement cohérent. De plus, $\mathrm{pH}_E = 7{,}0$ confirme qu'il s'agit bien d'un dosage acide fort -- base forte (sel neutre \ce{NaCl} à l'équivalence).

\textbf{Conclusion :} $C_a = \SI{6,2 \times 10^{-2}}{mol.L^{-1}}$, soit environ \SI{2,3}{g.L^{-1}} de \ce{HCl}.
\end{exoresolu}

\courssub{Influence de la dilution et des concentrations sur le saut de pH}

\textbf{Question d'intuition :} si l'on ajoute de l'eau distillée dans le bécher avant de doser (pour immerger l'électrode), le résultat est-il faussé ? Et si les solutions sont très diluées, le saut de pH est-il toujours aussi net ?

\begin{propriete}[Effet de la dilution et des concentrations]
\begin{itemize}
\item L'ajout d'eau distillée dans le bécher \emph{ne change pas} la quantité de matière d'acide $n_a = C_a V_a$ : le volume équivalent $V_{bE}$ est donc \emph{inchangé}. La relation $C_a V_a = C_b V_{bE}$ reste valable.
\item En revanche, la dilution \emph{réduit l'amplitude du saut de pH} : le pH initial est plus élevé (acide dilué) et le palier final plus bas (soude diluée). Le saut devient moins vertical et le point équivalent moins précis à repérer.
\item Plus les concentrations $C_a$ et $C_b$ sont élevées, plus le saut de pH est \emph{grand et brutal} (par exemple de pH $3$ à pH $11$ pour des solutions à \SI{0,1}{mol.L^{-1}}, contre pH $5$ à pH $9$ seulement pour des solutions à \SI{1e-3}{mol.L^{-1}}).
\end{itemize}
\end{propriete}

\begin{attention}[Précision et concentration]
Pour un dosage précis, on choisit des solutions \emph{suffisamment concentrées} (typiquement \SIrange{e-2}{e-1}{mol.L^{-1}}) afin d'obtenir un saut de pH ample et bien vertical. Un saut trop étalé rend la méthode des tangentes incertaine et le choix d'un indicateur coloré (section suivante) délicat. Mais attention : diluer le bécher avec de l'eau distillée n'est jamais une erreur de dosage, c'est même recommandé pour bien immerger l'électrode !
\end{attention}

\begin{aretenir}[L'essentiel de la section]
\begin{itemize}
\item Le dosage pH-métrique suit le pH en fonction du volume $V_b$ de base versée : burette, pH-mètre étalonné, agitation permanente.
\item La courbe pH $= f(V_b)$ présente quatre zones : pH initial faible, montée lente, saut brutal, palier basique.
\item Le point équivalent $E$ est le point d'inflexion de la courbe ; pour un dosage acide fort -- base forte, $\mathrm{pH}_E = 7$ à \SI{25}{\celsius} (sel neutre).
\item $E$ se détermine par la méthode des tangentes parallèles ou par le maximum de la courbe dérivée $\mathrm{d}\mathrm{pH}/\mathrm{d}V_b$.
\item L'exploitation repose sur la relation $C_a V_a = C_b V_{bE}$.
\item La dilution ne modifie pas $V_{bE}$ mais réduit l'amplitude du saut de pH.
\end{itemize}
\end{aretenir}

\begin{exercice}[À toi de jouer]
On dose $V_a = \SI{20,0}{mL}$ d'une solution d'acide chlorhydrique par de la soude à $C_b = \SI{0,10}{mol.L^{-1}}$. La courbe dérivée présente un maximum pour $V_{bE} = \SI{15,0}{mL}$.
\begin{enumerate}
\item Écrire l'équation de la réaction de dosage et justifier qu'elle convient à un dosage.
\item Quelle est la valeur du pH à l'équivalence ? Justifier.
\item Calculer la concentration $C_a$ de l'acide, puis le pH initial de la solution.
\end{enumerate}
\end{exercice}

\begin{corrige}[Exercice]
\begin{enumerate}
\item \ce{H3O+ + OH- -> 2 H2O} : réaction totale ($K = 10^{14}$), rapide et unique, donc utilisable pour un dosage.
\item $\mathrm{pH}_E = 7{,}0$ à \SI{25}{\celsius} : à l'équivalence, la solution ne contient que \ce{Na+} et \ce{Cl-}, ions spectateurs sans action sur l'eau ; le sel formé est neutre.
\item $C_a = \dfrac{C_b V_{bE}}{V_a} = \dfrac{0{,}10 \times 15{,}0}{20{,}0} = \SI{7,5 \times 10^{-2}}{mol.L^{-1}}$. pH initial : $-\log(7{,}5\times 10^{-2}) \approx 1{,}1$.
\end{enumerate}
\end{corrige}

Nous savons désormais doser avec précision un acide fort par une base forte. Mais que se passe-t-il lorsque l'acide est \emph{faible}, comme l'acide éthanoïque du vinaigre ? La courbe change d'allure, le pH à l'équivalence n'est plus $7$, et un phénomène nouveau apparaît à la demi-équivalence : c'est l'objet de la section suivante.

\coursec{Dosage d'un acide faible par une base forte (et réciproquement)}

Dans la section précédente, nous avons dosé l'acide chlorhydrique, un \cle{acide fort}, par la soude : la courbe $pH = f(V_b)$ partait d'un pH très bas, présentait un immense saut vertical, et l'équivalence se situait exactement à $pH_E = 7$. Mais la plupart des acides de la vie courante — le vinaigre, l'acide citrique du jus de bissap, l'acide lactique du lait caillé — sont des \cle{acides faibles} : ils ne se dissocient que partiellement dans l'eau. Que devient la courbe de dosage dans ce cas ? C'est toute la richesse de cette section : nous allons découvrir une courbe à l'allure nouvelle, un équivalent basique ($pH_E > 7$), et surtout un point remarquable, la \cle{demi-équivalence}, qui permet de mesurer expérimentalement le $pK_a$ d'un couple acide-base.

\courssub{Exemple type : dosage de l'acide éthanoïque par la soude}

\textbf{Situation concrète.}\par
Au laboratoire du lycée, on veut vérifier le titre d'un vinaigre commercial. On prélève $V_a = \SI{20,0}{mL}$ d'une solution d'acide éthanoïque \ce{CH3COOH} de concentration $C_a = \SI{0,10}{mol.L^{-1}}$, placée dans un bécher avec un pH-mètre, et l'on verse progressivement, à la burette, une solution d'hydroxyde de sodium \ce{NaOH} (soude) de concentration $C_b = \SI{0,10}{mol.L^{-1}}$.

\textbf{La réaction de dosage.}\par
L'acide éthanoïque est un acide faible ($pK_a(\ce{CH3COOH}/\ce{CH3COO-}) = 4,8$), mais la réaction avec les ions hydroxyde est \textbf{quasi-totale} car son avancement est poussé très loin : la constante d'équilibre vaut
\[
K = \frac{K_a}{K_e} = \frac{10^{-4,8}}{10^{-14}} = 10^{9,2} \gg 10^4 .
\]
L'équation-bilan s'écrit donc avec une flèche simple :
\[
\ce{CH3COOH + OH- -> CH3COO- + H2O}
\]
Cette réaction est \textbf{rapide}, \textbf{quasi-totale} et \textbf{unique} : elle possède toutes les qualités d'une réaction de dosage. À noter : on écrit \ce{CH3COOH} (forme majoritaire de l'acide faible en solution) et non \ce{H3O+}, contrairement au dosage de l'acide fort.

\begin{attention}[Ne pas confondre « faible » et « peu réactif »]
Un acide \emph{faible} se dissocie peu dans l'eau, mais il réagit \emph{totalement} avec une base forte. « Faible » décrit l'équilibre de dissociation dans l'eau pure ; « totale » décrit la réaction de dosage. Les deux notions sont indépendantes.
\end{attention}

\courssub{Allure de la courbe $pH = f(V_b)$}

Suivons le pH au cours du versement de soude, en raisonnant sur les espèces présentes :

\begin{enumerate}
\item \textbf{$V_b = 0$ (point de départ).} La solution ne contient que l'acide faible \ce{CH3COOH}. Comme il est peu dissocié, $[\ce{H3O+}]$ est plus petite que $C_a$ : le pH initial vaut environ $2,9$ (et non $1$ comme pour un acide fort de même concentration). \textbf{Le pH initial est plus élevé.}
\item \textbf{$0 < V_b < V_{bE}$ (avant l'équivalence).} Les ions \ce{OH-} versés consomment \ce{CH3COOH} et produisent \ce{CH3COO-}. Le mélange contient le couple \ce{CH3COOH}/\ce{CH3COO-} : c'est une \textbf{solution tampon}, dont le pH varie lentement. La courbe monte doucement.
\item \textbf{$V_b = V_{bE}$ (équivalence).} Tout l'acide a été consommé ; la solution contient l'ion éthanoate \ce{CH3COO-}, qui est une \textbf{base faible}. Le pH est donc \textbf{supérieur à 7} (environ $8,7$ ici). Le saut de pH, bien visible, est \textbf{plus court} que pour un acide fort : il s'étend environ de $7$ à $10$, au lieu de $4$ à $10$.
\item \textbf{$V_b > V_{bE}$ (après l'équivalence).} La soude est en excès ; c'est elle qui impose le pH. La courbe rejoint l'asymptote $pH \approx 12,5$, comme dans le cas de l'acide fort.
\end{enumerate}

\begin{popfigure}
\begin{center}
\begin{tikzpicture}
\begin{axis}[
  width=0.92\linewidth, height=7.2cm,
  xlabel={$V_b$ (mL)}, ylabel={pH},
  xmin=0, xmax=40, ymin=0, ymax=14,
  xtick={0,5,10,15,20,25,30,35,40},
  ytick={0,2,4,6,8,10,12,14},
  grid=both, grid style={popline!40},
  axis line style={popdark},
  clip=false
]
% Avant équivalence : zone tampon (Henderson-Hasselbalch)
\addplot[popcurve, domain=0.5:19.5, samples=200] {4.8 + log10(x/(20-x))};
% Saut d'équivalence (linéaire approximatif)
\addplot[popcurve, domain=19.5:20.5, samples=50] {8.72 + 4.5*(x-20)};
% Après équivalence : excès de soude
\addplot[popcurve, domain=20.5:40, samples=200] {14 + log10(0.1*(x-20)/(x+20))};
% Point initial (calculé : pH = 1/2(pKa - log Ca) = 2.87)
\addplot[only marks, mark=*, mark size=2pt, poporange] coordinates {(0,2.87)};
% Équivalence
\addplot[dashed, poporange, thick] coordinates {(20,0) (20,8.72)};
\addplot[dashed, poporange, thick] coordinates {(0,8.72) (20,8.72)};
\node[poporange, font=\small\bfseries] at (axis cs:23.5,7.6) {E};
\node[poporange, font=\footnotesize] at (axis cs:26,9.4) {$pH_E \approx 8,7 > 7$};
% Demi-équivalence
\addplot[dashed, popgreen, thick] coordinates {(10,0) (10,4.8)};
\addplot[dashed, popgreen, thick] coordinates {(0,4.8) (10,4.8)};
\node[popgreen, font=\small\bfseries] at (axis cs:12.6,3.9) {$E_{1/2}$};
\node[popgreen, font=\footnotesize] at (axis cs:15.5,5.6) {$pH = pK_a = 4,8$};
\addplot[only marks, mark=*, mark size=2pt, popgreen] coordinates {(10,4.8)};
\addplot[only marks, mark=*, mark size=2pt, poporange] coordinates {(20,8.72)};
\end{axis}
\end{tikzpicture}
\end{center}
\legende{Courbe de dosage de \SI{20,0}{mL} d'acide éthanoïque \SI{0,10}{mol.L^{-1}} par la soude \SI{0,10}{mol.L^{-1}}. On repère l'équivalence E ($V_{bE} = \SI{20,0}{mL}$, $pH_E > 7$) et la demi-équivalence $E_{1/2}$ ($V_b = \SI{10,0}{mL}$, $pH = pK_a = 4,8$).}
\end{popfigure}

\begin{aretenir}[Trois différences avec le dosage d'un acide fort]
Pour le dosage d'un acide faible par une base forte, par rapport au dosage d'un acide fort de même concentration par la même base :
\begin{itemize}
\item le \textbf{pH initial est plus élevé} (l'acide est peu dissocié) ;
\item le \textbf{saut de pH est plus court} et décalé vers les pH élevés ;
\item le \textbf{pH à l'équivalence est supérieur à 7} ($pH_E > 7$).
\end{itemize}
En revanche, le volume équivalent $V_{bE}$ est \textbf{le même} : il ne dépend que des quantités de matière, pas de la force de l'acide.
\end{aretenir}

\courssub{Pourquoi $pH_E > 7$ à l'équivalence ?}

C'est LE point que les élèves oublient le plus souvent. Raisonnons sur les espèces.

À l'équivalence, la réaction $\ce{CH3COOH + OH- -> CH3COO- + H2O}$ a consommé exactement tout l'acide et toute la base versée. Que reste-t-il dans le bécher ? Des ions \ce{Na+} (spectateurs, sans action sur l'eau) et des ions éthanoate \ce{CH3COO-}. Or \ce{CH3COO-} est la \textbf{base conjuguée} de l'acide éthanoïque : elle réagit avec l'eau selon l'équilibre
\[
\ce{CH3COO- + H2O <=> CH3COOH + OH-}
\]
qui produit des ions hydroxyde. La solution à l'équivalence est donc \textbf{basique} : $pH_E > 7$.

\begin{attention}[Erreur fréquente : écrire $pH_E = 7$]
Beaucoup d'élèves écrivent machinalement $pH_E = 7$ à toute équivalence. C'est \textbf{faux} dès qu'un des réactifs est faible !
\begin{itemize}
\item Acide fort + base forte : $pH_E = 7$ (seuls des ions spectateurs restent).
\item Acide faible + base forte : $pH_E > 7$ (la base conjuguée \ce{A-} formée rend la solution basique).
\item Acide fort + base faible : $pH_E < 7$ (l'acide conjugué \ce{BH+} formé rend la solution acide).
\end{itemize}
Retiens l'image : \emph{à l'équivalence, c'est le « fort » qui gagne}. La soude étant forte, elle impose un pH basique.
\end{attention}

\courssub{Le point de demi-équivalence : $pH = pK_a$}

\begin{definition}[Demi-équivalence]
La \cle{demi-équivalence} est l'état du mélange obtenu lorsque le volume de base versé vaut la moitié du volume équivalent :
\[
V_b = \frac{V_{bE}}{2}.
\]
À la demi-équivalence, la moitié de l'acide faible a été transformée en sa base conjuguée : le mélange contient des \textbf{quantités égales} d'acide et de base conjuguée :
\[
[\ce{AH}] = [\ce{A-}].
\]
\end{definition}

Démontrons la conséquence capitale. Pour le couple \ce{CH3COOH}/\ce{CH3COO-}, la relation d'Henderson-Hasselbalch s'écrit :
\[
pH = pK_a + \log\frac{[\ce{CH3COO-}]}{[\ce{CH3COOH}]}.
\]
À la demi-équivalence, $[\ce{CH3COO-}] = [\ce{CH3COOH}]$, donc le logarithme vaut $\log 1 = 0$, et il reste :
\[
\boxed{\;pH_{1/2} = pK_a\;}
\]
C'est un résultat magnifique : \textbf{la simple lecture du pH à la demi-équivalence donne directement le $pK_a$ du couple}, sans aucun calcul de concentration.

\begin{methode}[Déterminer expérimentalement le $pK_a$ d'un acide faible]
\begin{enumerate}
\item Tracer la courbe $pH = f(V_b)$ et repérer le point équivalent E (méthode des tangentes ou de la dérivée) ; noter $V_{bE}$.
\item Calculer $V_b = \dfrac{V_{bE}}{2}$ et reporter cette valeur sur l'axe des volumes.
\item Lire l'ordonnée du point de la courbe d'abscisse $V_{bE}/2$.
\item Ce pH est le $pK_a$ du couple $\ce{AH}/\ce{A-}$.
\end{enumerate}
\end{methode}

\begin{exemplebox}[Le $pK_a$ de l'acide éthanoïque lu sur la courbe]
Sur la courbe de dosage du vinaigre, la méthode des tangentes donne $V_{bE} = \SI{20,0}{mL}$. On se place donc à $V_b = \SI{10,0}{mL}$ et on lit le pH correspondant : $pH = 4,8$. Conclusion immédiate :
\[
pK_a(\ce{CH3COOH}/\ce{CH3COO-}) = 4,8.
\]
C'est exactement la valeur des tables. La courbe de dosage est ainsi un véritable « appareil de mesure » de $pK_a$.
\end{exemplebox}

\begin{savaistu}[La demi-équivalence, zone de pouvoir tampon maximal]
C'est précisément à la demi-équivalence que le mélange contient autant d'acide que de base conjuguée : il peut alors « absorber » aussi bien un ajout d'acide qu'un ajout de base. Le \cle{pouvoir tampon} y est \textbf{maximal}. C'est pourquoi la courbe y est presque plate : autour de $V_b = V_{bE}/2$, le pH varie très lentement malgré l'ajout de soude. Les chimistes qui préparent des solutions tampons exploitent exactement cette propriété (section 5 de cette leçon).
\end{savaistu}

\courssub{Cas réciproque : dosage d'une base faible par un acide fort}

\textbf{Situation.}\par
On dose maintenant $V_b = \SI{20,0}{mL}$ d'une solution d'ammoniac \ce{NH3} de concentration $C_b = \SI{0,10}{mol.L^{-1}}$ par une solution d'acide chlorhydrique \ce{HCl} ($\ce{H3O+} + \ce{Cl-}$) de concentration $C_a = \SI{0,10}{mol.L^{-1}}$. L'ammoniac est une base faible : $pK_a(\ce{NH4+}/\ce{NH3}) = 9,2$.

La réaction de dosage, quasi-totale ($K = \dfrac{1}{K_a} = 10^{9,2} \gg 10^4$), s'écrit :
\[
\ce{NH3 + H3O+ -> NH4+ + H2O}
\]

Le raisonnement est symétrique du précédent :
\begin{itemize}
\item \textbf{pH initial} : environ $11,1$ (base faible, donc inférieur au pH d'une base forte de même concentration, qui vaudrait $13$) ;
\item \textbf{avant l'équivalence} : mélange tampon \ce{NH4+}/\ce{NH3}, pH qui décroît lentement ;
\item \textbf{à l'équivalence} : seuls restent \ce{NH4+} (acide faible) et \ce{Cl-} (spectateur) : la solution est \textbf{acide}, $pH_E < 7$ (environ $5,3$) ;
\item \textbf{à la demi-équivalence} ($V_a = V_{aE}/2$) : $[\ce{NH4+}] = [\ce{NH3}]$, donc
\[
pH_{1/2} = pK_a(\ce{NH4+}/\ce{NH3}) = 9,2 ;
\]
\item \textbf{après l'équivalence} : l'acide fort en excès impose un pH voisin de $1,7$.
\end{itemize}

\begin{popfigure}
\begin{center}
\begin{tikzpicture}
\begin{axis}[
  width=0.92\linewidth, height=7.2cm,
  xlabel={$V_a$ (mL)}, ylabel={pH},
  xmin=0, xmax=40, ymin=0, ymax=14,
  xtick={0,5,10,15,20,25,30,35,40},
  ytick={0,2,4,6,8,10,12,14},
  grid=both, grid style={popline!40},
  axis line style={popdark},
  clip=false
]
% Avant équivalence : zone tampon (Henderson-Hasselbalch pour base)
\addplot[popcurve, domain=0.5:19.5, samples=200] {9.2 + log10((20-x)/x)};
% Saut d'équivalence (linéaire approximatif)
\addplot[popcurve, domain=19.5:20.5, samples=50] {5.3 - 4.5*(x-20)};
% Après équivalence : excès d'acide fort
\addplot[popcurve, domain=20.5:40, samples=200] {-log10(0.1*(x-20)/(x+20))};
% Point initial (calculé : pOH = 1/2(pKb - log Cb) = 2.87 -> pH = 11.13)
\addplot[only marks, mark=*, mark size=2pt, popblue] coordinates {(0,11.13)};
% Équivalence
\addplot[dashed, poporange, thick] coordinates {(20,0) (20,5.3)};
\addplot[dashed, poporange, thick] coordinates {(0,5.3) (20,5.3)};
\node[poporange, font=\small\bfseries] at (axis cs:23.5,4.3) {E};
\node[poporange, font=\footnotesize] at (axis cs:26.5,6.1) {$pH_E \approx 5,3 < 7$};
% Demi-équivalence
\addplot[dashed, popgreen, thick] coordinates {(10,0) (10,9.2)};
\addplot[dashed, popgreen, thick] coordinates {(0,9.2) (10,9.2)};
\node[popgreen, font=\small\bfseries] at (axis cs:12.8,8.2) {$E_{1/2}$};
\node[popgreen, font=\footnotesize] at (axis cs:16,10) {$pH = pK_a = 9,2$};
\addplot[only marks, mark=*, mark size=2pt, popgreen] coordinates {(10,9.2)};
\addplot[only marks, mark=*, mark size=2pt, poporange] coordinates {(20,5.3)};
\end{axis}
\end{tikzpicture}
\end{center}
\legende{Dosage de \SI{20,0}{mL} d'ammoniac \SI{0,10}{mol.L^{-1}} par l'acide chlorhydrique \SI{0,10}{mol.L^{-1}} : la courbe est décroissante, $pH_E < 7$, et la demi-équivalence donne $pK_a(\ce{NH4+}/\ce{NH3}) = 9,2$.}
\end{popfigure}

\courssub{Comparaison : acide fort ou acide faible dosé par la même base forte}

Superposons les deux courbes de dosage : \ce{HCl} \SI{0,10}{mol.L^{-1}} et \ce{CH3COOH} \SI{0,10}{mol.L^{-1}}, dosés chacun par la soude \SI{0,10}{mol.L^{-1}}.

\begin{popfigure}
\begin{center}
\begin{tikzpicture}
\begin{axis}[
  width=0.92\linewidth, height=7.6cm,
  xlabel={$V_b$ (mL)}, ylabel={pH},
  xmin=0, xmax=40, ymin=0, ymax=14,
  xtick={0,5,10,15,20,25,30,35,40},
  ytick={0,2,4,6,8,10,12,14},
  grid=both, grid style={popline!40},
  axis line style={popdark},
  legend style={font=\footnotesize, at={(0.03,0.97)}, anchor=north west, draw=popline, fill=popcream},
  clip=false
]
% Acide fort : HCl 0.1M 20mL par NaOH 0.1M
% Avant équivalence
\addplot[popblue, very thick, domain=0.5:19.5, samples=150] {1 - log10((20-x)/(20+x))};
% Saut
\addplot[popblue, very thick, domain=19.5:20.5, samples=40] {7 + 5*(x-20)};
% Après équivalence
\addplot[popblue, very thick, domain=20.5:40, samples=150] {14 + log10(0.1*(x-20)/(x+20))};
% Point initial
\addplot[only marks, mark=*, mark size=2pt, popblue] coordinates {(0,1)};

% Acide faible : CH3COOH 0.1M 20mL par NaOH 0.1M
% Avant équivalence
\addplot[poporange, very thick, domain=0.5:19.5, samples=200] {4.8 + log10(x/(20-x))};
% Saut
\addplot[poporange, very thick, domain=19.5:20.5, samples=50] {8.72 + 4.5*(x-20)};
% Après équivalence
\addplot[poporange, very thick, domain=20.5:40, samples=200] {14 + log10(0.1*(x-20)/(x+20))};
% Point initial
\addplot[only marks, mark=*, mark size=2pt, poporange] coordinates {(0,2.87)};

\legend{{\ce{HCl} (acide fort)}, {\ce{CH3COOH} (acide faible)}}
% Ligne équivalence
\addplot[dashed, popmuted] coordinates {(20,0) (20,14)};
\node[popdark, font=\footnotesize] at (axis cs:20.6,13.2) {$V_{bE}$};
% Points équivalence
\addplot[only marks, mark=*, mark size=2pt, popblue] coordinates {(20,7)};
\addplot[only marks, mark=*, mark size=2pt, poporange] coordinates {(20,8.72)};
\node[popblue, font=\footnotesize] at (axis cs:24.5,6.4) {$pH_E = 7$};
\node[poporange, font=\footnotesize] at (axis cs:25.5,9.6) {$pH_E \approx 8,7$};
\end{axis}
\end{tikzpicture}
\end{center}
\legende{Superposition des courbes de dosage d'un acide fort (\ce{HCl}) et d'un acide faible (\ce{CH3COOH}) de même concentration par la même soude : même volume équivalent, mais pH initial, longueur du saut et $pH_E$ différents.}
\end{popfigure}

\begin{center}
\begin{tabularx}{\linewidth}{|l|X|X|}
\hline
\rowcolor{popblueL} \textbf{Critère} & \textbf{Acide fort (\ce{HCl})} & \textbf{Acide faible (\ce{CH3COOH})} \\
\hline
pH initial & $1,0$ & $2,9$ \\
\hline
Réaction de dosage & \ce{H3O+ + OH- -> 2H2O} & \ce{CH3COOH + OH- -> CH3COO- + H2O} \\
\hline
Zone tampon avant E & non & oui (couple \ce{CH3COOH}/\ce{CH3COO-}) \\
\hline
$pH_E$ & $= 7$ & $> 7$ ($\approx 8,7$) \\
\hline
Saut de pH & grand ($\approx 4$ à $10$) & plus court ($\approx 7$ à $10$) \\
\hline
$V_{bE}$ & \SI{20,0}{mL} & \SI{20,0}{mL} (identique) \\
\hline
Demi-équivalence & sans intérêt particulier & $pH = pK_a = 4,8$ \\
\hline
\end{tabularx}
\end{center}

\begin{exoresolu}[Exploitation complète d'une courbe de dosage]
On dose $V_a = \SI{25,0}{mL}$ d'une solution d'un acide faible \ce{AH} de concentration $C_a$ inconnue par de la soude à $C_b = \SI{0,050}{mol.L^{-1}}$. La courbe $pH = f(V_b)$ donne : équivalence à $V_{bE} = \SI{20,0}{mL}$ avec $pH_E = 8,4$ ; à $V_b = \SI{10,0}{mL}$, on lit $pH = 3,8$.
\begin{enumerate}
\item Écrire l'équation de la réaction de dosage.
\item Déterminer $C_a$.
\item Identifier le $pK_a$ du couple et proposer un nom pour l'acide.
\item Justifier le signe de $pH_E - 7$.
\end{enumerate}
\tcblower
\textbf{1. Équation de dosage.} L'acide faible réagit avec les ions hydroxyde selon une réaction quasi-totale :
\[ \ce{AH + OH- -> A- + H2O} \]
\textbf{2. Concentration $C_a$.} À l'équivalence, $n_{\text{versé}}(\ce{OH-}) = n_{\text{initial}}(\ce{AH})$, soit $C_a V_a = C_b V_{bE}$ :
\[
C_a = \frac{C_b \times V_{bE}}{V_a} = \frac{0,050 \times 20,0}{25,0} = \SI{0,040}{mol.L^{-1}}.
\]
\textbf{3. $pK_a$ du couple.} À la demi-équivalence, $V_b = \dfrac{V_{bE}}{2} = \SI{10,0}{mL}$, et l'on lit $pH = 3,8$. Or à la demi-équivalence $pH = pK_a$ :
\[
pK_a(\ce{AH}/\ce{A-}) = 3,8.
\]
D'après les tables, il pourrait s'agir de l'acide méthanoïque \ce{HCOOH} ($pK_a = 3,8$).
\textbf{4. Signe de $pH_E - 7$.} À l'équivalence, la solution contient la base conjuguée \ce{A-}, qui réagit avec l'eau : $\ce{A- + H2O <=> AH + OH-}$. La production d'ions \ce{OH-} rend la solution basique, d'où $pH_E = 8,4 > 7$.
\end{exoresolu}

\begin{exercice}[À toi de jouer]
On dose \SI{20,0}{mL} d'une solution d'ammoniac de concentration inconnue par de l'acide chlorhydrique à \SI{0,10}{mol.L^{-1}}. L'équivalence est obtenue pour $V_{aE} = \SI{15,0}{mL}$ et le pH vaut alors $5,4$.
\begin{enumerate}
\item Écrire l'équation de la réaction de dosage.
\item Calculer la concentration de la solution d'ammoniac.
\item À quel volume versé lit-on $pH = pK_a(\ce{NH4+}/\ce{NH3})$ ? Que vaut ce pH ?
\item Expliquer pourquoi $pH_E < 7$.
\end{enumerate}
\end{exercice}

\begin{corrige}[Exercice : dosage de l'ammoniac]
\textbf{1.} $\ce{NH3 + H3O+ -> NH4+ + H2O}$ (réaction quasi-totale, $K = 10^{9,2}$).
\textbf{2.} À l'équivalence : $C_b V_b = C_a V_{aE}$, d'où $C_b = \dfrac{0,10 \times 15,0}{20,0} = \SI{0,075}{mol.L^{-1}}$.
\textbf{3.} À la demi-équivalence, $V_a = \dfrac{V_{aE}}{2} = \SI{7,5}{mL}$, et $pH = pK_a(\ce{NH4+}/\ce{NH3}) = 9,2$.
\textbf{4.} À l'équivalence, la solution contient l'ion ammonium \ce{NH4+}, acide faible qui cède des protons à l'eau : $\ce{NH4+ + H2O <=> NH3 + H3O+}$. La solution est donc acide : $pH_E = 5,4 < 7$.
\end{corrige}

\begin{aretenir}[L'essentiel de la section]
\begin{itemize}
\item Acide faible dosé par une base forte : $\ce{AH + OH- -> A- + H2O}$ ; pH initial plus élevé, saut plus court, $pH_E > 7$ car la base conjuguée \ce{A-} basifie la solution à l'équivalence.
\item À la \cle{demi-équivalence} ($V_b = V_{bE}/2$) : $[\ce{AH}] = [\ce{A-}]$ et $pH = pK_a$ ; c'est aussi la zone de pouvoir tampon maximal.
\item Base faible dosée par un acide fort : $\ce{B + H3O+ -> BH+ + H2O}$ ; courbe décroissante, $pH_E < 7$, et $pH = pK_a(\ce{BH+}/\ce{B})$ à la demi-équivalence.
\item Le volume équivalent est le même pour un acide fort et un acide faible de même concentration : seule la \emph{forme} de la courbe change.
\end{itemize}
\end{aretenir}

\coursec{Dosage colorimétrique et choix de l'indicateur coloré}

Dans la section précédente, nous avons suivi le dosage d'un acide faible par une base forte \emph{au pH-mètre} : l'appareil mesure le pH après chaque ajout de soude, et l'équivalence se repère sur la courbe $pH = f(V_B)$ par la méthode des tangentes. C'est précis, mais cela exige un pH-mètre étalonné, une sonde fragile, du temps. Or, au laboratoire du lycée comme dans une savonnerie artisanale ou une brasserie, on a souvent besoin d'une réponse \emph{rapide} : « ai-je atteint l'équivalence, oui ou non ? ». C'est exactement le rôle du \cle{dosage colorimétrique} : remplacer l'œil électronique du pH-mètre par un œil chimique, l'\cle{indicateur coloré}, qui change de couleur au bon moment. Tout l'art du chimiste consiste alors à choisir \emph{le bon} indicateur — c'est-à-dire celui qui vire précisément là où se situe le saut de pH. C'est l'objet de cette section.

\courssub{Principe du dosage colorimétrique}

\textbf{Idée intuitive.}\par
Imagine que tu doses le jus de bissap (acide) par une solution de soude. Sans instrument, rien ne te dit quand l'équivalence est atteinte : le mélange reste limpide et sans odeur particulière. Mais si tu ajoutes au préalable quelques gouttes d'un colorant « intelligent » — une substance qui est \emph{elle-même} un couple acide/base et dont les deux formes ont des couleurs différentes — alors le virage de ce colorant, brutalement provoqué par le saut de pH à l'équivalence, te signale l'arrivée du point d'équivalence. On arrête alors la burette et on lit le volume versé $V_E$.

\begin{definition}[Dosage colorimétrique]
Un \cle{dosage colorimétrique} est un dosage acido-basique dans lequel l'équivalence est repérée par le \textbf{changement de couleur} d'un \cle{indicateur coloré} acido-basique ajouté en faible quantité (2 à 3 gouttes) dans la solution à doser. Le volume lu à la burette au moment du virage persistant est le volume équivalent $V_E$, qui permet d'exploiter la relation à l'équivalence :
\[ n_{\text{acide dosé}} = n_{\text{base versée à l'équivalence}} \quad\Longleftrightarrow\quad C_A V_A = C_B V_E \]
(pour un monoacide et une monobase).
\end{definition}

\textbf{Qu'est-ce qu'un indicateur coloré, au juste ?}\par
Un indicateur coloré acido-basique est un couple acide/base faible, noté $\ce{HIn}/\ce{In-}$, dont la particularité est que la forme acide $\ce{HIn}$ et la forme basique $\ce{In-}$ n'ont \textbf{pas la même couleur}. L'équilibre qui gouverne sa couleur est :
\[ \ce{HIn(aq) + H2O(l) <=> In-(aq) + H3O+(aq)} \qquad K_a(\ce{HIn}) = \frac{[\ce{In-}][\ce{H3O+}]}{[\ce{HIn}]} \]
Cet équilibre est \textbf{rapide} et se déplace selon le pH du milieu :
\begin{itemize}
  \item si le milieu est \textbf{acide} ($pH < pK_a - 1$), l'équilibre est déplacé vers la gauche : la forme $\ce{HIn}$ prédomine, la solution prend la \emph{teinte acide} ;
  \item si le milieu est \textbf{basique} ($pH > pK_a + 1$), l'équilibre est déplacé vers la droite : la forme $\ce{In-}$ prédomine, la solution prend la \emph{teinte basique} ;
  \item entre les deux, les deux formes coexistent en proportions comparables : on observe une \emph{teinte intermédiaire} (mélange des deux couleurs), appelée \textbf{teinte sensible}.
\end{itemize}

\begin{definition}[Zone de virage]
La \cle{zone de virage} d'un indicateur coloré est l'intervalle de pH, centré sur le $pK_a$ du couple $\ce{HIn}/\ce{In-}$ et d'étendue approximative $pK_a(\ce{HIn}) \pm 1$, dans lequel l'indicateur passe progressivement de sa teinte acide à sa teinte basique. C'est dans cette zone que l'œil perçoit le changement de couleur.
\end{definition}

\begin{aretenir}[Zones de virage des trois indicateurs usuels]
\begin{center}
\rowcolors{1}{popcream}{white}
\begin{tabularx}{\linewidth}{|l|X|X|X|}
\hline
\rowcolor{popblueL} \textbf{Indicateur} & \textbf{Teinte acide} & \textbf{Zone de virage} & \textbf{Teinte basique} \\
\hline
Hélianthine & rouge & 3,1 -- 4,4 & jaune orangé \\
Bleu de bromothymol (BBT) & jaune & 6,0 -- 7,6 & bleu \\
Phénolphtaléine & incolore & 8,2 -- 10,0 & rose (fuchsia) \\
\hline
\end{tabularx}
\end{center}
\end{aretenir}

\begin{popfigure}
\begin{center}
\begin{tikzpicture}[x=0.62cm]
% Echelle de pH
\draw[thick, popdark, -{Stealth}] (0,0) -- (14.8,0) node[below left=2pt and -30pt, font=\small] {};
\foreach \x in {0,1,...,14} {
  \draw[popdark] (\x,0.12) -- (\x,-0.12) node[below, font=\footnotesize] {\x};
}
\node[below, font=\small\itshape, popdark] at (7,-0.85) {pH};
% Hélianthine
\fill[red!70!poporange, opacity=0.85] (3.1,0.55) rectangle (4.4,1.15);
\draw[popdark] (3.1,0.55) rectangle (4.4,1.15);
\node[left, font=\small, align=right] at (3.0,0.85) {\textbf{Hélianthine}\\ 3,1 -- 4,4};
\node[font=\footnotesize] at (3.75,0.85) {virage};
% BBT
\fill[popblue!70, opacity=0.85] (6.0,1.55) rectangle (7.6,2.15);
\draw[popdark] (6.0,1.55) rectangle (7.6,2.15);
\node[left, font=\small, align=right] at (5.9,1.85) {\textbf{BBT}\\ 6,0 -- 7,6};
\node[font=\footnotesize, white] at (6.8,1.85) {virage};
% Phénolphtaléine
\fill[poppink!85!popink, opacity=0.85] (8.2,2.55) rectangle (10.0,3.15);
\draw[popdark] (8.2,2.55) rectangle (10.0,3.15);
\node[left, font=\small, align=right] at (8.1,2.85) {\textbf{Phénolphtaléine}\\ 8,2 -- 10,0};
\node[font=\footnotesize] at (9.1,2.85) {virage};
% pH = 7
\draw[dashed, poppurple, thick] (7,0) -- (7,3.4);
\node[above, font=\footnotesize, poppurple] at (7,3.4) {$pH = 7$};
\end{tikzpicture}
\end{center}
\legende{Zones de virage des trois indicateurs usuels, reportées sur une échelle de pH de 0 à 14. Chaque indicateur ne « voit » qu'une petite fenêtre de pH : il faut donc l'accorder au saut de pH du dosage.}
\end{popfigure}

\courssub{La règle d'or du choix de l'indicateur}

\textbf{Idée intuitive.}\par
L'indicateur est un « détecteur de pH » à fenêtre étroite. Pour qu'il signale l'équivalence au bon moment, il faut que sa fenêtre (sa zone de virage) soit placée \emph{là où le pH bascule brutalement}, c'est-à-dire dans le \cle{saut de pH} qui encadre le point équivalent. Si l'indicateur vire trop tôt ou trop tard, le volume $V_E$ lu sera faux, et toute la concentration calculée avec lui.

\begin{methode}[Choisir judicieusement un indicateur coloré]
\begin{enumerate}
  \item \textbf{Identifier la nature du dosage} (acide fort/base forte, acide faible/base forte, base faible/acide fort) et, si possible, \textbf{estimer le $pH_E$ à l'équivalence} ou tracer la courbe $pH = f(V)$.
  \item \textbf{Repérer le saut de pH} autour du point équivalent (la portion quasi verticale de la courbe).
  \item \textbf{Choisir l'indicateur dont la zone de virage contient le $pH_E$} ; à défaut, choisir celui dont la zone de virage est \textbf{entièrement comprise dans le saut de pH}.
  \item \textbf{Vérifier la teinte attendue} : le virage doit être net et facile à observer (par exemple, passage de l'incolore au rose persistant pour la phénolphtaléine).
  \item \textbf{N'ajouter que 2 à 3 gouttes} d'indicateur dans le bécher, agiter, puis doser jusqu'au virage persistant (au moins 30 secondes).
\end{enumerate}
\end{methode}

\begin{popfigure}
\begin{center}
\begin{tikzpicture}
\begin{axis}[
  width=0.92\linewidth, height=7.2cm,
  xlabel={$V_B$ (mL)}, ylabel={pH},
  xmin=0, xmax=25, ymin=0, ymax=14,
  xtick={0,5,10,15,20,25}, ytick={0,2,4,6,8,10,12,14},
  grid=both, grid style={popline!40},
  axis line style={popdark},
  label style={popdark},
]
% Bandes zones de virage
\fill[red!60!poporange, opacity=0.30] (axis cs:0,3.1) rectangle (axis cs:25,4.4);
\fill[popblue, opacity=0.25] (axis cs:0,6.0) rectangle (axis cs:25,7.6);
\fill[poppink!80!popink, opacity=0.30] (axis cs:0,8.2) rectangle (axis cs:25,10.0);
\node[font=\scriptsize, popdark] at (axis cs:22.5,3.75) {hélianthine};
\node[font=\scriptsize, popdark] at (axis cs:22.5,6.8) {BBT};
\node[font=\scriptsize, popdark] at (axis cs:22.5,9.1) {phénolphtaléine};
% Courbe dosage acide faible par base forte (pKa=4.8, Ve=10)
\addplot[popcurve, domain=0:9.6, samples=150] {4.8 + ln(x/(10-x))/ln(10)};
\addplot[popcurve, domain=9.6:10.4, samples=60] {4.8 + ln(x/(10-x))/ln(10)};
\addplot[popcurve, domain=10.4:25, samples=120] {14 + ln((x-10)/(x+10))/ln(10)};
% Point équivalent
\addplot[only marks, mark=*, mark size=2pt, poporange] coordinates {(10,8.7)};
\draw[dashed, poporange] (axis cs:10,0) -- (axis cs:10,8.7) -- (axis cs:0,8.7);
\node[font=\scriptsize, poporange, anchor=west] at (axis cs:10.3,7.9) {$E\;(pH_E \approx 8,7)$};
\node[font=\scriptsize, popdark, anchor=north] at (axis cs:10,0.1) {$V_E$};
\end{axis}
\end{tikzpicture}
\end{center}
\legende{Dosage d'un acide faible (vinaigre) par la soude : le saut de pH (environ 7 à 10) chevauche la zone de virage de la phénolphtaléine, qui convient parfaitement. Le BBT vire trop tôt, l'hélianthine bien trop tôt : ils sont à proscrire ici.}
\end{popfigure}

\courssub{Choix de l'indicateur selon la nature du dosage}

\textbf{Cas 1 : dosage d'un acide fort par une base forte (ou l'inverse).}\par
Le saut de pH est très ample (typiquement de $pH \approx 3$ à $pH \approx 11$) et le point équivalent se situe à $pH_E = 7$. Les trois indicateurs ont leur zone de virage incluse dans le saut : tous conviennent en théorie. Le \textbf{BBT} est l'indicateur \emph{idéal} car sa zone de virage (6,0 -- 7,6) contient $pH_E = 7$ : il vire du jaune au bleu (teinte sensible verte) exactement à l'équivalence. La \textbf{phénolphtaléine} reste acceptable car sa zone est encore dans le saut. L'hélianthine, en fin de saut, donne une précision un peu moindre mais reste tolérable.

\textbf{Cas 2 : dosage d'un acide faible par une base forte.}\par
À l'équivalence, la solution contient la base conjuguée de l'acide faible : le milieu est \textbf{basique}, $pH_E > 7$ (souvent entre 8 et 9). Le saut de pH est réduit et décalé vers le haut. Seule la \textbf{phénolphtaléine} (zone 8,2 -- 10,0) convient. Le BBT virerait \emph{avant} l'équivalence, et l'hélianthine bien plus tôt encore : le volume $V_E$ serait sous-estimé.

\textbf{Cas 3 : dosage d'une base faible par un acide fort.}\par
Situation symétrique : à l'équivalence, la solution contient l'acide conjugué de la base faible, donc $pH_E < 7$ (souvent entre 5 et 6). C'est l'\textbf{hélianthine} (zone 3,1 -- 4,4), ou parfois le BBT en limite basse, qui convient. La phénolphtaléine virerait \emph{avant} l'équivalence (virage prématuré) : elle est à proscrire.

\begin{aretenir}[Tableau récapitulatif du choix]
\begin{center}
\rowcolors{1}{popcream}{white}
\begin{tabularx}{\linewidth}{|X|X|X|}
\hline
\rowcolor{popblueL} \textbf{Dosage} & \textbf{$pH_E$} & \textbf{Indicateur conseillé} \\
\hline
Acide fort / base forte & $pH_E = 7$ & BBT (idéal), phénolphtaléine acceptable \\
Acide faible / base forte & $pH_E > 7$ & Phénolphtaléine \\
Base faible / acide fort & $pH_E < 7$ & Hélianthine \\
\hline
\end{tabularx}
\end{center}
\end{aretenir}

\begin{attention}[Erreur fréquente : la phénolphtaléine « passe-partout »]
Beaucoup d'élèves utilisent systématiquement la phénolphtaléine, parce que son virage incolore $\to$ rose est spectaculaire. \textbf{C'est une erreur grave} pour le dosage d'une base faible par un acide fort : le $pH_E$ étant acide (vers 5--6), la phénolphtaléine vire (devient incolore) vers $pH \approx 8,3$... c'est-à-dire \emph{avant} d'avoir versé assez d'acide pour atteindre l'équivalence. Le virage est \textbf{prématuré} (trop précoce), le volume équivalent $V_E$ lu est sous-estimé, et la concentration calculée est fausse. \textbf{Réflexe : toujours raisonner sur la position du $pH_E$ avant de choisir l'indicateur.}
\end{attention}

\courssub{Exemple chiffré : dosage du vinaigre par la soude}

\begin{exoresolu}[Vérification du degré d'acidité d'un vinaigre]
Énoncé. Un vinaigre commercial du marché de Mokolo affiche « $6^{\circ}$ », c'est-à-dire \SI{6}{g} d'acide éthanoïque \ce{CH3COOH} pour \SI{100}{g} de vinaigre. Pour vérifier cette étiquette, on dilue le vinaigre 10 fois, puis on dose $V_A = \SI{10,0}{mL}$ de la solution diluée par de la soude de concentration $C_B = \SI{0,10}{mol.L^{-1}}$. Le suivi pH-métrique donne un point équivalent à $V_E = \SI{10,2}{mL}$ avec $pH_E \approx 8,7$. (1) Quel indicateur coloré choisir pour un dosage colorimétrique ? Justifier. (2) Écrire l'équation-bilan de la réaction de dosage. (3) Calculer la concentration du vinaigre dilué, puis celle du vinaigre commercial, et vérifier l'étiquette. On donne $M(\ce{CH3COOH}) = \SI{60}{g.mol^{-1}}$ et on assimilera la densité du vinaigre à 1.
\tcblower
Solution détaillée.
\textbf{(1) Choix de l'indicateur.} Il s'agit du dosage d'un acide faible (\ce{CH3COOH}, $pK_a = 4,8$) par une base forte (\ce{NaOH}). Le point équivalent est basique : $pH_E \approx 8,7$. La zone de virage de la \textbf{phénolphtaléine} (8,2 -- 10,0) contient cette valeur : c'est l'indicateur adapté. On observera le passage de l'\emph{incolore} au \emph{rose persistant}. Le BBT (6,0 -- 7,6) virerait avant l'équivalence et l'hélianthine bien trop tôt : tous deux sous-estimeraient $V_E$.

\textbf{(2) Équation-bilan de la réaction de dosage.} La réaction entre l'acide éthanoïque et les ions hydroxyde est quasi totale ($K \gg 1$) et rapide :
\[ \ce{CH3COOH(aq) + OH-(aq) -> CH3COO-(aq) + H2O(l)} \]

\textbf{(3) Exploitation de l'équivalence.} À l'équivalence : $C_A V_A = C_B V_E$, d'où
\[ C_A = \frac{C_B V_E}{V_A} = \frac{0{,}10 \times 10{,}2}{10{,}0} = \SI{0,102}{mol.L^{-1}}. \]
Le vinaigre commercial est 10 fois plus concentré : $C = 10 \times 0{,}102 = \SI{1,02}{mol.L^{-1}}$, soit un titre massique :
\[ t = C \times M = 1{,}02 \times 60 \approx \SI{61}{g.L^{-1}} \approx \SI{6,1}{g} \text{ pour } \SI{100}{g}. \]
Le vinaigre affiche bien environ $6^{\circ}$ : \textbf{l'étiquette est conforme}. Le virage de la phénolphtaléine au rose persistant à $V_E \approx \SI{10,2}{mL}$ permet d'obtenir ce résultat sans pH-mètre.
\end{exoresolu}

\courssub{Limites du dosage colorimétrique}

Le dosage colorimétrique est rapide, économique et suffisant pour la plupart des contrôles de qualité, mais il faut connaître ses limites :

\begin{itemize}
  \item \textbf{Précision moindre.} La zone de virage s'étale sur 1 à 2 unités pH et l'œil humain apprécie le changement de teinte avec une incertitude d'environ une goutte à quelques gouttes de réactif. L'incertitude relative sur $V_E$ est typiquement de l'ordre de 1 à 2 \%, contre mieux que 0,5 \% en pH-métrie avec la méthode des tangentes ou de la dérivée.
  \item \textbf{Solutions colorées ou troubles.} Si la solution à doser est déjà fortement colorée (jus de bissap rouge foncé, vin de palme fermenté, extrait de pharmacopée), la teinte de l'indicateur devient \textbf{imperceptible} : le dosage colorimétrique est alors impossible, et seul le dosage pH-métrique (ou conductimétrique) reste utilisable.
  \item \textbf{Saut de pH trop faible.} Pour un acide très faible ou très dilué, le saut de pH est si réduit qu'aucune zone de virage ne s'y loge correctement : le virage est progressif et l'équivalence indétectable à l'œil.
  \item \textbf{Dosages acide faible / base faible.} La réaction n'étant pas totale et le saut quasi inexistant, le dosage colorimétrique (comme le pH-métrique d'ailleurs) est impraticable : on remplace l'un des réactifs par un réactif fort.
\end{itemize}

\begin{savaistu}[Pourquoi seulement 2 à 3 gouttes d'indicateur ?]
Un indicateur coloré est lui-même un \textbf{acide faible} ($\ce{HIn}$) : il consomme donc une petite quantité du réactif de dosage ! Si l'on en versait un volume important, il fausserait la mesure en décalant l'équivalence. On se limite à 2 ou 3 gouttes : c'est suffisant pour que la couleur soit visible (ces molécules sont des colorants très intenses) et négligeable devant les quantités dosées. C'est aussi pour cette raison que la phénolphtaléine, utilisée depuis le XIX\textsuperscript{e} siècle, reste irremplaçable : quelques microgrammes suffisent à colorer tout un erlenmeyer d'un rose éclatant. Dans les laboratoires de contrôle de la SONARA ou des brasseries camerounaises, ce même réflexe — « deux gouttes, pas plus » — fait partie des gestes de base transmis aux techniciens.
\end{savaistu}

\begin{exercice}[À toi de jouer]
On dose \SI{20,0}{mL} d'une solution d'ammoniac \ce{NH3} (base faible, $pK_a(\ce{NH4+}/\ce{NH3}) = 9,2$) par de l'acide chlorhydrique à \SI{0,050}{mol.L^{-1}}. Le point équivalent est atteint pour $V_E = \SI{18,4}{mL}$ et le pH à l'équivalence vaut environ 5,5.
\begin{enumerate}
  \item Écrire l'équation-bilan de la réaction de dosage et justifier son caractère quasi total.
  \item Quel indicateur coloré faut-il choisir parmi l'hélianthine, le BBT et la phénolphtaléine ? Justifier et préciser le changement de teinte observé.
  \item Calculer la concentration de la solution d'ammoniac.
\end{enumerate}
\end{exercice}

\begin{corrige}[Exercice : dosage de l'ammoniac]
\textbf{1.} $\ce{NH3(aq) + H3O+(aq) -> NH4+(aq) + H2O(l)}$. La constante de cette réaction est $K = \dfrac{1}{K_a(\ce{NH4+}/\ce{NH3})} = 10^{9,2} \gg 10^4$ : la réaction est \textbf{quasi totale} (et rapide), elle convient donc pour un dosage.

\textbf{2.} Le dosage est celui d'une \textbf{base faible par un acide fort} : $pH_E \approx 5,5 < 7$. Seule l'\textbf{hélianthine} (zone de virage 3,1 -- 4,4) a sa zone proche du saut de pH côté acide ; c'est elle qu'il faut choisir (le BBT, zone 6,0 -- 7,6, est acceptable en limite, la phénolphtaléine est à proscrire car elle virerait bien avant l'équivalence). On observe le passage du \emph{jaune orangé} au \emph{rouge-orangé} persistant.

\textbf{3.} À l'équivalence : $C_B V_B = C_A V_E$, d'où
\[ C_B = \frac{C_A V_E}{V_B} = \frac{0{,}050 \times 18{,}4}{20{,}0} = \SI{0,046}{mol.L^{-1}}. \]
\end{corrige}

\begin{aretenir}[L'essentiel de la section]
\begin{itemize}
  \item Le \cle{dosage colorimétrique} repère l'équivalence par le \textbf{virage d'un indicateur coloré}, couple $\ce{HIn}/\ce{In-}$ dont les deux formes ont des couleurs différentes.
  \item La \cle{zone de virage} ($pK_a(\ce{HIn}) \pm 1$) doit \textbf{contenir le $pH_E$} ou être \textbf{incluse dans le saut de pH}.
  \item Acide fort/base forte : \textbf{BBT} idéal ($pH_E = 7$). Acide faible/base forte : \textbf{phénolphtaléine} ($pH_E > 7$). Base faible/acide fort : \textbf{hélianthine} ($pH_E < 7$).
  \item On n'ajoute que \textbf{2 à 3 gouttes} d'indicateur, car c'est un acide faible qui consomme du réactif.
  \item Limites : précision moindre que la pH-métrie, impossibilité pour les solutions colorées ou les dosages à saut de pH trop faible.
\end{itemize}
\end{aretenir}

\coursec{Solutions tampons : définition, propriétés et préparation}

Dans la section précédente, nous avons vu que le choix d'un indicateur coloré repose sur la comparaison entre sa zone de virage et le pH à l'équivalence. Mais as-tu remarqué que, dans un dosage d'acide faible par une base forte, la courbe $pH = f(V_b)$ présente une zone presque horizontale autour de la demi-équivalence ? Dans cette région, on peut verser plusieurs millilitres de soude sans que le pH ne bouge presque. Ce comportement étonnant n'est pas un hasard : c'est la signature d'une \cle{solution tampon}. C'est précisément ce phénomène que nous allons maintenant étudier, comprendre et apprendre à exploiter.

\courssub{Qu'est-ce qu'une solution tampon ?}

\textbf{Une expérience qui fait réfléchir}\par

Imagine deux béchers. Dans le premier, tu verses \SI{100}{mL} d'eau pure (pH $= 7{,}0$). Dans le second, tu verses \SI{100}{mL} d'une solution contenant de l'acide éthanoïque et de l'éthanoate de sodium (pH $= 4{,}8$). Tu ajoutes maintenant dans chaque bécher \SI{1}{mL} d'acide chlorhydrique à \SI{1}{mol.L^{-1}}. Résultat : le pH de l'eau pure chute de $7{,}0$ à environ $2{,}0$, soit une variation de \textbf{5 unités de pH} ! Dans le second bécher, le pH passe de $4{,}8$ à environ $4{,}7$ : une variation de seulement \textbf{0,1 unité}. La seconde solution a « amorti » l'attaque acide, comme un pare-chocs absorbe un choc. C'est une solution tampon.

\begin{definition}[Solution tampon]
Une \cle{solution tampon} est une solution dont le \textbf{pH varie faiblement} lors de :
\begin{itemize}
\item l'ajout \textbf{modéré} d'un acide fort ;
\item l'ajout \textbf{modéré} d'une base forte ;
\item une \textbf{dilution modérée}.
\end{itemize}
Une solution tampon est constituée d'un mélange d'un \textbf{acide faible} et de sa \textbf{base conjuguée}, en concentrations voisines (couple $\ce{AH}/\ce{A-}$ avec $[\ce{AH}] \approx [\ce{A-}]$).
\end{definition}

\textbf{Pourquoi ce mélange résiste-t-il aux variations de pH ?}\par

Le secret réside dans la présence simultanée des deux partenaires du couple, chacun spécialisé dans la neutralisation d'un type d'« agresseur » :

\begin{itemize}
\item Si l'on ajoute un \textbf{acide fort} (ions \ce{H3O+}), ce sont les ions \ce{A-} qui interviennent, selon la réaction quasi totale :
\[ \ce{A- + H3O+ -> AH + H2O} \]
Les ions oxonium ajoutés sont consommés et remplacés par de l'acide faible \ce{AH}, beaucoup moins agressif pour le pH.
\item Si l'on ajoute une \textbf{base forte} (ions \ce{HO-}), c'est l'acide faible \ce{AH} qui réagit, selon la réaction quasi totale :
\[ \ce{AH + HO- -> A- + H2O} \]
Les ions hydroxyde sont consommés et remplacés par la base faible \ce{A-}.
\item En cas de \textbf{dilution modérée}, les concentrations $[\ce{AH}]$ et $[\ce{A-}]$ diminuent dans la même proportion : leur rapport reste inchangé, donc le pH aussi (voir la relation de Henderson-Hasselbalch ci-dessous).
\end{itemize}

\begin{popfigure}
\begin{center}
\begin{tikzpicture}[scale=0.9, every node/.style={font=\small}]
% Bécher eau pure
\draw[thick, popblue] (-1.2,0) -- (-1.2,2.2) (1.2,0) -- (1.2,2.2) (-1.2,0) -- (1.2,0);
\fill[popblue!25] (-1.15,0.05) rectangle (1.15,1.2);
\node[popdark] at (0,0.6) {eau pure};
\node[popdark] at (0,-0.5) {pH $= 7{,}0$};
% flèche ajout acide
\draw[-{Stealth}, very thick, poppink] (0,3.4) -- (0,2.5);
\node[poppink, align=center] at (0,4.0) {$+ \SI{1}{mL}$ de \ce{HCl}};
% résultat
\draw[-{Stealth}, very thick, popdark] (1.6,1.1) -- (3.0,1.1);
\node[poppink, font=\bfseries, align=center] at (4.6,1.1) {pH $= 2{,}0$\\ $\Delta pH \approx 5$ !};
% Bécher tampon
\begin{scope}[xshift=10.5cm]
\draw[thick, popgreen] (-1.2,0) -- (-1.2,2.2) (1.2,0) -- (1.2,2.2) (-1.2,0) -- (1.2,0);
\fill[popgreen!25] (-1.15,0.05) rectangle (1.15,1.2);
\node[popdark] at (0,0.6) {tampon \ce{CH3COOH}/\ce{CH3COO-}};
\node[popdark] at (0,-0.5) {pH $= 4{,}8$};
\draw[-{Stealth}, very thick, poppink] (0,3.4) -- (0,2.5);
\node[poppink, align=center] at (0,4.0) {$+ \SI{1}{mL}$ de \ce{HCl}};
\draw[-{Stealth}, very thick, popdark] (1.6,1.1) -- (3.0,1.1);
\node[popgreen!60!black, font=\bfseries, align=center] at (4.6,1.1) {pH $= 4{,}7$\\ $\Delta pH \approx 0{,}1$};
\end{scope}
\end{tikzpicture}
\end{center}
\legende{Ajout du même volume d'acide fort dans l'eau pure et dans une solution tampon : la variation de pH est environ 50 fois plus faible dans le tampon.}
\end{popfigure}

\courssub{Le pH d'une solution tampon : relation de Henderson-Hasselbalch}

Considérons le couple acide faible / base conjuguée $\ce{AH}/\ce{A-}$, dont l'équilibre en solution s'écrit :
\[ \ce{AH + H2O <=> A- + H3O+} \]
Cet équilibre est régi par la constante d'acidité :
\[ K_a = \frac{[\ce{A-}][\ce{H3O+}]}{[\ce{AH}]} \]
En isolant $[\ce{H3O+}]$ puis en prenant l'opposé du logarithme décimal, on obtient la \cle{relation de Henderson-Hasselbalch} :

\begin{propriete}[Relation de Henderson-Hasselbalch]
Pour toute solution contenant un couple acide faible / base conjuguée $\ce{AH}/\ce{A-}$ :
\[ pH = pK_a + \log\frac{[\ce{A-}]}{[\ce{AH}]} \]
Conséquence fondamentale : lorsque $[\ce{AH}] = [\ce{A-}]$, le logarithme vaut zéro et
\[ pH = pK_a \]
C'est exactement la situation de la \textbf{demi-équivalence} rencontrée dans le dosage d'un acide faible par une base forte : c'est là que la courbe est la plus plate, donc là que l'effet tampon est \textbf{maximal}.
\end{propriete}

\begin{aretenir}[Les trois idées forces]
\begin{itemize}
\item Un tampon = un couple $\ce{AH}/\ce{A-}$ présent en quantités voisines.
\item Son pH est voisin du $pK_a$ du couple : $pH = pK_a + \log\dfrac{[\text{base}]}{[\text{acide}]}$.
\item L'effet tampon est maximal quand $[\ce{AH}] = [\ce{A-}]$, c'est-à-dire quand $pH = pK_a$.
\end{itemize}
\end{aretenir}

\begin{exemplebox}[Fabrication artisanale d'un tampon éthanoïque]
On mélange $V_a = \SI{50}{mL}$ d'une solution d'acide éthanoïque \ce{CH3COOH} à $C = \SI{0,10}{mol.L^{-1}}$ et $V_b = \SI{50}{mL}$ d'une solution d'éthanoate de sodium \ce{CH3COONa} à $C = \SI{0,10}{mol.L^{-1}}$ (on rappelle $pK_a(\ce{CH3COOH}/\ce{CH3COO-}) = 4{,}8$).

Quantités mélangées : $n(\ce{CH3COOH}) = \num{0,10} \times \num{0,050} = \SI{5,0e-3}{mol}$ et $n(\ce{CH3COO-}) = \SI{5,0e-3}{mol}$.

Les volumes étant égaux et les concentrations identiques, on a exactement $[\ce{CH3COOH}] = [\ce{CH3COO-}]$ dans le mélange, donc :
\[ pH = pK_a + \log(1) = 4{,}8 \]
On a préparé \SI{100}{mL} d'une solution tampon de pH $= 4{,}8$, capable de résister aux ajouts modérés d'acide ou de base.
\end{exemplebox}

\courssub{Préparation d'une solution tampon}

\begin{methode}[Préparer un tampon de pH $= pK_a$ (méthode du mélange direct)]
\begin{enumerate}
\item Choisir un couple $\ce{AH}/\ce{A-}$ dont le $pK_a$ est égal (ou très proche) du pH désiré.
\item Préparer deux solutions de \textbf{même concentration} $C$ : l'une de l'acide faible \ce{AH}, l'autre du sel de la base conjuguée (par exemple \ce{NaA}, qui libère \ce{A-} par dissolution totale : $\ce{NaA -> Na+ + A-}$).
\item Mélanger des \textbf{volumes égaux} $V$ des deux solutions. Alors $[\ce{AH}] = [\ce{A-}]$ et $pH = pK_a$.
\item Vérifier le pH obtenu au pH-mètre et ajuster si nécessaire par ajout d'une goutte d'acide ou de base diluée.
\end{enumerate}
\end{methode}

\begin{popfigure}
\begin{center}
\begin{tikzpicture}[scale=0.85, every node/.style={font=\small}]
% Bécher 1 : acide
\draw[thick, poporange] (-1,0) -- (-1,1.8) (1,0) -- (1,1.8) (-1,0) -- (1,0);
\fill[poporange!30] (-0.95,0.05) rectangle (0.95,0.9);
\node[popdark, align=center] at (0,0.45) {\ce{CH3COOH}\\ $C = \SI{0,10}{mol.L^{-1}}$};
\node[popdark] at (0,-0.5) {$V = \SI{50}{mL}$};
% Bécher 2 : base conjuguée
\begin{scope}[xshift=4.5cm]
\draw[thick, popblue] (-1,0) -- (-1,1.8) (1,0) -- (1,1.8) (-1,0) -- (1,0);
\fill[popblue!30] (-0.95,0.05) rectangle (0.95,0.9);
\node[popdark, align=center] at (0,0.45) {\ce{CH3COONa}\\ $C = \SI{0,10}{mol.L^{-1}}$};
\node[popdark] at (0,-0.5) {$V = \SI{50}{mL}$};
\end{scope}
% Flèches vers le grand bécher
\draw[-{Stealth}, very thick, popdark] (0,2.1) .. controls (3.5,3.2) .. (5.6,2.6);
\draw[-{Stealth}, very thick, popdark] (4.5,2.1) .. controls (5.2,2.6) .. (6.2,2.7);
% Grand bécher
\begin{scope}[xshift=7.6cm]
\draw[thick, poppurple] (-1.3,0) -- (-1.3,2.4) (1.3,0) -- (1.3,2.4) (-1.3,0) -- (1.3,0);
\fill[poppurple!25] (-1.25,0.05) rectangle (1.25,1.5);
\node[popdark, align=center] at (0,0.75) {\textbf{Tampon}\\ $[\ce{CH3COOH}] = [\ce{CH3COO-}]$};
\node[poppurple, font=\bfseries] at (0,-0.5) {pH $= pK_a = 4{,}8$};
\end{scope}
\end{tikzpicture}
\end{center}
\legende{Préparation d'une solution tampon de pH $= pK_a$ : mélange de volumes égaux de deux solutions de même concentration, l'une de l'acide faible, l'autre de sa base conjuguée.}
\end{popfigure}

\textbf{Autre méthode : la demi-équivalence d'un dosage}\par

Il existe une seconde stratégie, très élégante, qui fait le lien avec les sections précédentes. On part d'une solution d'acide faible seul et on la dose par une base forte, mais \textbf{on s'arrête à la demi-équivalence}. À ce stade, la moitié de l'acide a été transformée en base conjuguée selon :
\[ \ce{AH + HO- -> A- + H2O} \quad\text{(réaction quasi totale, } K = 10^{pK_a(\text{eau}) - pK_a(\ce{AH})} \gg 1\text{)} \]
Il reste donc dans le bécher $n(\ce{AH}) = n(\ce{A-}) = \dfrac{n_0}{2}$ : on a fabriqué \emph{in situ} un tampon de pH $= pK_a$. C'est pourquoi le point de demi-équivalence, repéré sur la courbe $pH = f(V_b)$ par la condition $V_b = \dfrac{V_E}{2}$, permet à la fois de mesurer le $pK_a$ et de préparer un tampon.

\begin{exoresolu}[Préparation d'un tampon par demi-dosage]
On dispose de $V_a = \SI{100}{mL}$ d'une solution d'acide éthanoïque de concentration $C_a = \SI{0,20}{mol.L^{-1}}$ et d'une solution de soude de concentration $C_b = \SI{0,50}{mol.L^{-1}}$. Quel volume de soude faut-il verser pour obtenir une solution tampon de pH $= 4{,}8$ ?
\tcblower
\textbf{Étape 1 : volume équivalent.} À l'équivalence, $C_a V_a = C_b V_E$, d'où :
\[ V_E = \frac{C_a V_a}{C_b} = \frac{\num{0,20} \times \num{100}}{\num{0,50}} = \SI{40}{mL} \]
\textbf{Étape 2 : demi-équivalence.} Pour obtenir $[\ce{CH3COOH}] = [\ce{CH3COO-}]$, il faut verser la moitié du volume équivalent :
\[ V_b = \frac{V_E}{2} = \SI{20}{mL} \]
\textbf{Vérification.} Après ajout de \SI{20}{mL} de soude : $n(\ce{HO-}) = \num{0,50} \times \num{0,020} = \SI{1,0e-2}{mol}$ a transformé \SI{1,0e-2}{mol} d'acide en éthanoate. Il reste $n(\ce{CH3COOH}) = \num{0,020} - \num{0,010} = \SI{1,0e-2}{mol}$ et il s'est formé $n(\ce{CH3COO-}) = \SI{1,0e-2}{mol}$. Les quantités sont égales, donc $pH = pK_a = 4{,}8$. Le tampon est prêt.
\end{exoresolu}

\begin{attention}[Le pH d'un tampon n'est pas figé !]
Erreur fréquente au baccalauréat : écrire qu'une solution tampon a un pH « constant » ou « invariable ». C'est faux. Le pH d'un tampon \textbf{varie peu, mais il varie}. De plus, l'effet tampon n'est garanti que pour des ajouts \textbf{modérés} : si l'on verse assez d'acide fort pour consommer la quasi-totalité des ions \ce{A-} (ou assez de base pour consommer tout l'acide \ce{AH}), le tampon est « saturé » et le pH s'effondre brutalement. On dit que la \cle{capacité tampon} est dépassée. Retiens : un tampon amortit les chocs, il ne les supprime pas.
\end{attention}

\courssub{Importance des solutions tampons}

Les solutions tampons ne sont pas une curiosité de laboratoire : elles sont partout, y compris en toi.

\begin{itemize}
\item \textbf{Le sang} : son pH doit rester compris entre $7{,}35$ et $7{,}45$. Il est tamponné principalement par le couple $\ce{H2CO3}/\ce{HCO3-}$ (acide carbonique / ion hydrogénocarbonate), secondé par les couples des phosphates et des protéines.
\item \textbf{Les médicaments injectables} : une solution injectée dans le sang doit être tamponnée au voisinage de pH $7{,}4$ pour ne pas provoquer de lésions ni de douleurs violentes.
\item \textbf{L'étalonnage des pH-mètres} : tout pH-mètre est étalonné avec des solutions tampons étalons commerciales (souvent pH $= 4{,}0$, $7{,}0$ et $10{,}0$), dont le pH est parfaitement connu et stable dans le temps.
\item \textbf{Les fermentations} : la fabrication de la bière, du vin de palme, du bili-bili ou du yaourt exige un pH maîtrisé ; les milieux de fermentation possèdent des systèmes tampons naturels qui protègent les levures et bactéries des variations brutales d'acidité.
\item \textbf{Les cosmétiques et les shampooings} « au pH neutre » ou « pH équilibré » sont tamponnés pour respecter la peau et les cheveux.
\end{itemize}

\begin{savaistu}[Ton sang, un tampon qui te garde en vie]
Le sang humain est tamponné principalement par le couple $\ce{H2CO3}/\ce{HCO3-}$, dont le $pK_a$ vaut $6{,}1$. Le rapport $[\ce{HCO3-}]/[\ce{H2CO3}]$ d'environ 20 déplace le pH sanguin vers $7{,}4$. Lors d'un effort musculaire intense, l'acide lactique produit par les muscles est neutralisé par les ions hydrogénocarbonate :
\[ \ce{CH3-CH(OH)-COOH + HCO3- -> CH3-CH(OH)-COO- + H2CO3} \]
puis l'acide carbonique formé est éliminé par les poumons sous forme de \ce{CO2} : c'est pour cela que tu respires plus vite quand tu cours ! Une variation du pH sanguin de seulement \textbf{0,2 unité} (acidose ou alcalose) met la vie en danger : sans ses tampons, ton organisme ne survivrait pas à un simple sprint.
\end{savaistu}

\begin{exercice}[À toi de jouer]
On souhaite préparer \SI{200}{mL} d'une solution tampon de pH $= 9{,}2$ à partir du couple $\ce{NH4+}/\ce{NH3}$ ($pK_a = 9{,}2$). On dispose d'une solution de chlorure d'ammonium \ce{NH4Cl} à \SI{0,10}{mol.L^{-1}} et d'une solution d'ammoniac \ce{NH3} à \SI{0,10}{mol.L^{-1}}.
\begin{enumerate}
\item Quels volumes de chaque solution faut-il mélanger ?
\item Calculer le pH du tampon obtenu si l'on mélange finalement \SI{120}{mL} de la solution d'ammoniac et \SI{80}{mL} de la solution de chlorure d'ammonium.
\end{enumerate}
\end{exercice}

\begin{corrige}[Exercice]
\begin{enumerate}
\item Pour obtenir $pH = pK_a = 9{,}2$, il faut $[\ce{NH4+}] = [\ce{NH3}]$. Les concentrations étant égales, on mélange des volumes égaux : $V = \SI{100}{mL}$ de chaque solution, pour un volume total de \SI{200}{mL}.
\item Quantités mélangées : $n(\ce{NH3}) = \num{0,10} \times \num{0,120} = \SI{1,2e-2}{mol}$ et $n(\ce{NH4+}) = \num{0,10} \times \num{0,080} = \SI{8,0e-3}{mol}$. Dans le même volume total, le rapport des concentrations est égal au rapport des quantités :
\[ pH = pK_a + \log\frac{n(\ce{NH3})}{n(\ce{NH4+})} = 9{,}2 + \log\frac{\num{1,2e-2}}{\num{8,0e-3}} = 9{,}2 + \log(1{,}5) \approx 9{,}2 + 0{,}18 \approx 9{,}4 \]
Le tampon obtenu a un pH légèrement supérieur au $pK_a$, car la base conjuguée est en excès.
\end{enumerate}
\end{corrige}

\begin{aretenir}[L'essentiel de la section]
\begin{itemize}
\item Une \cle{solution tampon} voit son pH varier \textbf{faiblement} lors d'ajouts modérés d'acide, de base ou d'eau.
\item Elle contient un couple acide faible / base conjuguée en concentrations voisines ; son pH est donné par $pH = pK_a + \log\dfrac{[\text{base}]}{[\text{acide}]}$.
\item Pour préparer un tampon de $pH = pK_a$ : mélanger des \textbf{volumes égaux} de solutions d'acide faible et de base conjuguée de \textbf{même concentration}, ou doser un acide faible jusqu'à la \textbf{demi-équivalence}.
\item L'effet tampon est maximal à $pH = pK_a$ et disparaît si la capacité tampon est dépassée.
\item Les tampons sont indispensables : sang, médicaments injectables, étalonnage des pH-mètres, fermentations.
\end{itemize}
\end{aretenir}

\newpage

\coursec{Travaux pratiques}

\courssub{Objectifs du TP}

À l'issue de cette séance de travaux pratiques, tu seras capable de :

\begin{itemize}
\item réaliser un \cle{dosage pH-métrique} de l'acide éthanoïque contenu dans un vinaigre commercial par une solution d'hydroxyde de sodium (soude) de concentration connue ;
\item tracer la courbe $pH = f(V_b)$ et déterminer le \cle{volume équivalent} $V_{bE}$ par la \cle{méthode des tangentes} ;
\item vérifier expérimentalement le \cle{degré d'acidité} annoncé sur l'étiquette du vinaigre ;
\item retrouver la valeur du $pK_a$ du couple \ce{CH3COOH}/\ce{CH3COO-} par lecture du pH à la \cle{demi-équivalence} ;
\item préparer une \cle{solution tampon} de $pH = pK_a$ et vérifier expérimentalement sa propriété essentielle : la stabilité du pH lors de l'ajout modéré d'acide ou de base.
\end{itemize}

\courssub{Matériel et produits}

\begin{center}
\begin{tabularx}{\linewidth}{|l|X|}
\hline
\rowcolor{popblueL} \textbf{Matériel} & \textbf{Produits} \\
\hline
Burette graduée de \SI{25}{mL} & Vinaigre commercial (degré indiqué sur l'étiquette, ex. $6^{\circ}$) \\
\hline
Pipette jaugée de \SI{10,0}{mL} + poire aspirante & Solution d'hydroxyde de sodium à \SI{0,10}{mol.L^{-1}} \\
\hline
Fiole jaugée de \SI{100,0}{mL} & Solution d'éthanoate de sodium à \SI{0,10}{mol.L^{-1}} \\
\hline
Béchers de \SI{100}{mL} et \SI{250}{mL} & Eau distillée \\
\hline
pH-mètre étalonné (ou papier pH) & Phénolphtaléine (si dosage colorimétrique) \\
\hline
Agitateur magnétique + barreau (ou agitateur en verre) & Acide chlorhydrique dilué (quelques gouttes) \\
\hline
Support de burette, pince, entonnoir & \\
\hline
\end{tabularx}
\end{center}

\begin{savaistu}[Alternatives avec du matériel local]
Pas de pH-mètre au laboratoire ? Utilise la \cle{phénolphtaléine} comme indicateur coloré : sa zone de virage (8,2 -- 10,0) contient le pH à l'équivalence du dosage d'un acide faible par une base forte. Pas de solution d'éthanoate de sodium ? On peut la préparer en dissolvant des cristaux d'éthanoate de sodium \ce{CH3COONa} (environ \SI{0,82}{g} pour \SI{100}{mL} de solution à \SI{0,10}{mol.L^{-1}}). Le vinaigre blanc du marché convient parfaitement ; on peut aussi utiliser du jus de citron dilué pour un dosage d'entraînement.
\end{savaistu}

\courssub{Sécurité}

\begin{attention}[Consignes de sécurité]
\begin{itemize}
\item La \cle{soude} (hydroxyde de sodium) est \textbf{corrosive} (pictogramme : deux gouttes qui attaquent une main et une barre métallique). Elle provoque des brûlures graves de la peau et des yeux. Porte obligatoirement la \textbf{blouse}, les \textbf{gants} et les \textbf{lunettes de protection}.
\item En cas de contact avec la peau : rincer abondamment à l'eau claire pendant plusieurs minutes et prévenir l'enseignant.
\item \textbf{Ne jamais pipeter à la bouche} : utilise toujours la poire aspirante (propipette).
\item Le vinaigre concentré dégage des vapeurs irritantes : ne pas le sentir directement, travailler dans un espace aéré.
\item Avant le dosage, vérifier que la burette est bien fixée et que le robinet ne fuit pas.
\item Les solutions usées se versent dans le bidon de récupération prévu, jamais dans l'évier sans avis de l'enseignant.
\end{itemize}
\end{attention}

\courssub{Schéma du dispositif}

\begin{popfigure}
\begin{center}
\begin{tikzpicture}[scale=1.0, every node/.style={font=\small}]
% Support
\draw[popdark, line width=1.2pt] (-3.4,0) -- (-3.4,6.8);
\draw[popdark, line width=1.2pt] (-4.2,0) -- (-2.6,0);
% Pince
\draw[popdark, line width=1pt] (-3.4,5.6) -- (-0.6,5.6);
\draw[popdark, line width=1pt] (-0.6,5.4) rectangle (-0.2,5.8);
% Burette
\draw[popblue, line width=1pt] (-0.55,6.6) -- (-0.55,1.9);
\draw[popblue, line width=1pt] (-0.25,6.6) -- (-0.25,1.9);
\draw[popblue!60, fill=popblue!20] (-0.53,5.9) rectangle (-0.27,2.0);
\draw[popblue, line width=1pt] (-0.55,1.9) -- (-0.48,1.4) -- (-0.32,1.4) -- (-0.25,1.9);
\draw[popdark] (-0.75,1.55) -- (-0.05,1.55);
\draw[popdark, line width=0.8pt] (-0.4,1.4) -- (-0.4,1.0);
% graduations
\foreach \y in {2.4,2.9,3.4,3.9,4.4,4.9,5.4} \draw[popmuted] (-0.55,\y) -- (-0.45,\y);
\node[popblue, anchor=west] at (0.4,4.2) {Burette : soude \SI{0,10}{mol.L^{-1}}};
\draw[popmuted, ->] (0.4,4.0) -- (-0.2,3.6);
% goutte
\draw[popblue!70, fill=popblue!40] (-0.4,0.72) circle (0.05);
% Bécher
\draw[popdark, line width=1.1pt] (-1.5,0) -- (-1.5,1.6) (-0.9,0) arc (180:360:0.6 and 0.12) (0.7,0) -- (0.7,1.6);
\draw[poporange!70, fill=poporange!25] (-1.42,0.06) -- (-1.42,0.85) -- (0.62,0.85) -- (0.62,0.06) -- cycle;
\draw[poporange] (-1.42,0.85) -- (0.62,0.85);
\node[poporange!80!black, anchor=west] at (1.0,0.5) {Bécher : \SI{10,0}{mL} de vinaigre dilué + eau distillée};
% barreau magnétique
\draw[popdark, fill=popdark] (-0.55,0.10) rectangle (-0.05,0.20);
% sonde pH
\draw[popgreen!70!black, line width=1.2pt] (0.35,2.3) -- (0.35,0.55);
\draw[popgreen!70!black, fill=popgreen!30] (0.25,0.35) rectangle (0.45,0.6);
\draw[popgreen!70!black] (0.35,2.3) -- (1.6,2.6);
\node[popgreen!60!black, anchor=west] at (1.5,2.7) {Sonde du pH-mètre};
% pH-mètre
\draw[popdark, fill=popgreenL, rounded corners=2pt] (1.6,2.3) rectangle (3.4,3.3);
\node[popdark] at (2.5,3.0) {pH-mètre};
\node[popdark, font=\footnotesize] at (2.5,2.6) {pH = 2,9};
% agitateur
\draw[popmuted, dashed] (-1.7,-0.35) rectangle (0.9,-0.05);
\node[popmuted, font=\footnotesize] at (-0.4,-0.55) {agitateur magnétique};
\end{tikzpicture}
\end{center}
\legende{Dispositif de dosage pH-métrique du vinaigre dilué par la soude.}
\end{popfigure}

\courssub{Protocole}

\textbf{Partie A : dilution du vinaigre}

\begin{enumerate}
\item À l'aide de la pipette jaugée munie de la poire aspirante, prélever \SI{10,0}{mL} de vinaigre commercial.
\item Les introduire dans la fiole jaugée de \SI{100,0}{mL} et compléter avec de l'eau distillée jusqu'au trait de jauge. Boucher et homogénéiser. On obtient la solution $S$, diluée \textbf{10 fois}.
\end{enumerate}

\textbf{Partie B : dosage pH-métrique}

\begin{enumerate}
\item Rincer puis remplir la burette avec la solution de soude à $C_b = \SI{0,10}{mol.L^{-1}}$. Ajuster le zéro.
\item Pipeter $V_a = \SI{10,0}{mL}$ de la solution $S$ et les verser dans le bécher de \SI{250}{mL}. Ajouter environ \SI{20}{mL} d'eau distillée pour que la sonde du pH-mètre trempe correctement (cette eau ne modifie pas l'équivalence).
\item Introduire la sonde et le barreau magnétique. Mettre en marche l'agitation douce.
\item Relever le pH initial, puis verser la soude \textbf{mL par mL} jusqu'à \SI{8}{mL}, puis \textbf{goutte à goutte (0,2 mL)} autour du saut de pH, puis à nouveau mL par mL jusqu'à \SI{20}{mL}. Noter le pH après chaque ajout.
\item Tracer la courbe $pH = f(V_b)$ et déterminer $V_{bE}$ par la méthode des tangentes.
\end{enumerate}

\textbf{Partie C : préparation et test de la solution tampon}

\begin{enumerate}
\item Dans un bécher propre, mélanger \SI{25,0}{mL} de la solution $S$ (vinaigre dilué) et \SI{25,0}{mL} de la solution d'éthanoate de sodium de même concentration.
\item Mesurer le pH du mélange : c'est la solution tampon.
\item Ajouter \textbf{5 gouttes} d'acide chlorhydrique dilué, agiter, mesurer à nouveau le pH. Comparer avec la variation de pH qu'aurait provoquée le même ajout dans \SI{50}{mL} d'eau distillée.
\end{enumerate}

\courssub{Observations et résultats attendus}

Voici un exemple de tableau de mesures obtenu avec un vinaigre à $6^{\circ}$ dilué 10 fois ($C_a \approx \SI{0,10}{mol.L^{-1}}$) :

\begin{center}
\begin{tabular}{|c|*{9}{c}|}
\hline
\rowcolor{popblueL} $V_b$ (mL) & 0 & 2 & 4 & 6 & 8 & 9 & 9,6 & 10,0 & 10,4 \\
\hline
pH & 2,9 & 4,2 & 4,6 & 4,9 & 5,4 & 5,8 & 6,5 & 8,7 & 10,8 \\
\hline
\rowcolor{popblueL} $V_b$ (mL) & 11 & 12 & 14 & 16 & 18 & 20 & \multicolumn{3}{c|}{} \\
\hline
pH & 11,4 & 11,7 & 12,0 & 12,2 & 12,3 & 12,4 & \multicolumn{3}{c|}{} \\
\hline
\end{tabular}
\end{center}

Observations attendues :
\begin{itemize}
\item la courbe $pH = f(V_b)$ est \textbf{croissante} et présente un \textbf{saut de pH} modéré (environ de 6,5 à 10,8) autour de $V_{bE} \approx \SI{10,0}{mL}$ ;
\item le pH initial ($\approx 2,9$) est supérieur à celui d'un acide fort de même concentration : le vinaigre contient un \textbf{acide faible} ;
\item le pH à l'équivalence ($\approx 8,7$) est \textbf{supérieur à 7} (solution basique, l'ion éthanoate est une base faible) ;
\item à la demi-équivalence ($V_b = \SI{5,0}{mL}$), on lit $pH \approx 4,8 = pK_a$ du couple \ce{CH3COOH}/\ce{CH3COO-} ;
\item avec la phénolphtaléine, la solution vire au \textbf{rose persistant} dès la première goutte en excès après l'équivalence ;
\item le pH du tampon préparé est voisin de $4,8$ et varie à peine ($\Delta pH < 0,2$) après l'ajout des 5 gouttes d'acide, alors que le même ajout dans l'eau distillée ferait chuter le pH de plusieurs unités.
\end{itemize}

\courssub{Exploitation}

\begin{exoresolu}[Exploitation guidée du TP]
\textbf{Questions.} (1) Écrire l'équation-bilan de la réaction de dosage. (2) À partir de $V_{bE} = \SI{10,0}{mL}$, calculer la concentration $C_a$ de la solution diluée $S$, puis celle du vinaigre commercial. (3) En déduire le degré d'acidité du vinaigre et le comparer à l'étiquette ($6^{\circ}$). Données : $M(\ce{CH3COOH}) = \SI{60}{g.mol^{-1}}$ ; le degré d'acidité est la masse d'acide éthanoïque en grammes pour \SI{100}{g} de vinaigre (on prendra $\rho \approx \SI{1,0}{g.mL^{-1}}$). (4) Justifier le choix de la phénolphtaléine. (5) Expliquer pourquoi le mélange de la partie C est une solution tampon de $pH = pK_a$.
\tcblower
\textbf{Solution détaillée.}

(1) La réaction entre l'acide éthanoïque (acide faible) et la soude (base forte) est quasi totale :
\[ \ce{CH3COOH + OH- -> CH3COO- + H2O} \]

(2) À l'équivalence, les réactifs ont été mélangés dans les proportions stœchiométriques :
\[ C_a \times V_a = C_b \times V_{bE} \quad\Longrightarrow\quad C_a = \frac{C_b \times V_{bE}}{V_a} = \frac{0{,}10 \times 10{,}0}{10{,}0} = \SI{0,10}{mol.L^{-1}}. \]
Le vinaigre commercial est 10 fois plus concentré (dilution au dixième) :
\[ C_0 = 10 \times C_a = \SI{1,0}{mol.L^{-1}}. \]

(3) Concentration massique du vinaigre :
\[ C_m = C_0 \times M = 1{,}0 \times 60 = \SI{60}{g.L^{-1}}. \]
Un litre de vinaigre a une masse d'environ $\SI{1000}{mL} \times \SI{1,0}{g.mL^{-1}} = \SI{1000}{g}$ et contient \SI{60}{g} d'acide éthanoïque, soit \SI{6}{g} pour \SI{100}{g} de vinaigre : le degré d'acidité est donc $\approx 6^{\circ}$, \textbf{en accord avec l'étiquette}.

(4) Le pH à l'équivalence ($\approx 8,7$) appartient à la zone de virage de la phénolphtaléine (8,2 -- 10,0) : c'est l'indicateur judicieux. L'hélianthine (3,1 -- 4,4) virerait bien avant l'équivalence et donnerait un résultat faux.

(5) Le mélange contient des quantités égales d'acide éthanoïque et d'ions éthanoate : $[\ce{CH3COOH}] = [\ce{CH3COO-}]$. D'après la relation de Henderson-Hasselbalch :
\[ pH = pK_a + \log\frac{[\ce{CH3COO-}]}{[\ce{CH3COOH}]} = pK_a + \log 1 = pK_a \approx 4{,}8. \]
C'est bien une solution tampon : l'acide du couple consomme les ions \ce{OH-} ajoutés, la base du couple consomme les ions \ce{H3O+} ajoutés, d'où la faible variation de pH observée.
\end{exoresolu}

\courssub{Conclusion}

Ce TP t'a permis de vérifier, par un dosage pH-métrique rigoureux, que le vinaigre commercial étudié contient bien environ \SI{6}{g} d'acide éthanoïque pour \SI{100}{g}, conformément à son étiquette : la chimie est un outil de \textbf{contrôle de qualité} accessible au laboratoire du lycée. Tu as également retrouvé expérimentalement le $pK_a$ du couple \ce{CH3COOH}/\ce{CH3COO-} ($\approx 4,8$) par lecture du pH à la demi-équivalence, puis préparé une solution tampon dont tu as vérifié la propriété fondamentale : son pH varie très peu lors de l'ajout modéré d'acide ou de base. C'est ce même principe qui stabilise le pH de ton sang et celui de nombreuses préparations pharmaceutiques vendues dans les pharmacies du Cameroun.

\newpage

\coursec{Méthodes et exercices résolus}

\courssub{Fiches méthodes et exercices résolus}

\begin{methode}[Écrire l'équation-bilan d'une réaction de dosage]
\begin{enumerate}
\item \cle{Identifier} la nature des réactifs : acide fort ou faible ? base forte ou faible ? 
      Un acide fort est totalement dissocié en solution aqueuse : on écrit l'ion \ce{H3O+}. 
      Une base forte (soude \ce{NaOH}, potasse \ce{KOH}) fournit l'ion \ce{HO-}. 
      Un acide ou une base faible s'écrit sous sa \textbf{forme moléculaire} (\ce{CH3COOH}, \ce{NH3}, \ce{HCOOH}...).
\item \cle{Choisir} les couples acide/base mis en jeu et écrire les demi-équations de transfert de proton \ce{H+}.
\item \cle{Combiner} les demi-équations pour obtenir l'équation-bilan. Cas types :
\begin{itemize}
\item acide fort + base forte : \ce{H3O+ + HO- -> 2 H2O} ;
\item acide faible + base forte : \ce{AH + HO- -> A- + H2O} ;
\item acide fort + base faible : \ce{H3O+ + B -> BH+ + H2O}.
\end{itemize}
\item \cle{Vérifier} la conservation des éléments et des charges, puis préciser que la réaction est \textbf{totale} (constante d'équilibre $K \gg 1$), condition indispensable pour un dosage.
\end{enumerate}
\textbf{Réflexe :} ne jamais écrire \ce{CH3COOH + NaOH} : \ce{Na+} est un ion spectateur, il ne doit pas apparaître dans l'équation-bilan. De même, \ce{Cl-} ne s'écrit pas pour un dosage à l'acide chlorhydrique.
\end{methode}

\begin{exoresolu}[Dosage du vinaigre : équation-bilan]
On dose un volume $V_A = \SI{10,0}{mL}$ de vinaigre (acide éthanoïque \ce{CH3COOH}) par une solution d'hydroxyde de sodium de concentration $C_B = \SI{0,10}{mol.L^{-1}}$.
\begin{enumerate}
\item Identifier les couples acide/base mis en jeu.
\item Écrire l'équation-bilan de la réaction de dosage et justifier qu'elle peut servir de support à un dosage.
\end{enumerate}
On donne : $pK_a(\ce{CH3COOH}/\ce{CH3COO-}) = 4,8$ ; $pK_a(\ce{H2O}/\ce{HO-}) = 14$ (à \SI{25}{\celsius}).
\tcblower
\textbf{1. Couples mis en jeu.} L'acide éthanoïque est un acide \emph{faible} : il participe au couple \ce{CH3COOH}/\ce{CH3COO-}. L'ion hydroxyde \ce{HO-} est une base \emph{forte} : il participe au couple \ce{H2O}/\ce{HO-}.

\textbf{2. Équation-bilan.} L'acide \ce{CH3COOH} cède un proton à la base \ce{HO-} :
\[ \ce{CH3COOH + HO- -> CH3COO- + H2O} \]
Vérification : atomes conservés (C : 2, H : 5, O : 3 de chaque côté) ; charge $-1$ de chaque côté.

\textbf{Justification.} La constante de la réaction vaut :
\[ K = 10^{pK_a(\ce{H2O}/\ce{HO-}) - pK_a(\ce{CH3COOH}/\ce{CH3COO-})} = 10^{14 - 4,8} = 10^{9,2} \approx \num{1,6e9} \gg 1 \]
La réaction est donc \textbf{totale} ; de plus elle est \textbf{rapide} et \textbf{unique} : elle possède les trois qualités exigées d'une réaction de dosage.

\textbf{Conclusion :} le dosage du vinaigre par la soude repose sur la réaction totale \ce{CH3COOH + HO- -> CH3COO- + H2O}.
\end{exoresolu}

\begin{methode}[Tracer et exploiter la courbe pH = f(V)]
\begin{enumerate}
\item \cle{Repérer} les points particuliers : pH initial ($V_B = 0$), point de demi-équivalence ($V_{1/2} = V_E/2$), point équivalent $E$, pH final (excès de titrant).
\item \cle{Tracer} la courbe en reliant les points par une courbe lisse ; le \textbf{saut de pH} apparaît autour de l'équivalence.
\item \cle{Déterminer} $V_E$ par la \textbf{méthode des tangentes} :
\begin{itemize}
\item tracer deux tangentes parallèles à la courbe, de part et d'autre du saut de pH ;
\item tracer la parallèle à ces tangentes, équidistante des deux ;
\item son intersection avec la courbe donne le point équivalent $E$ ; on lit $V_E$ sur l'axe des abscisses et $pH_E$ sur l'axe des ordonnées.
\end{itemize}
\item \cle{Conclure} : à l'équivalence, les réactifs ont été mélangés dans les proportions stœchiométriques.
\end{enumerate}
\textbf{Réflexe :} pour un acide fort dosé par une base forte, $pH_E = 7$ (à \SI{25}{\celsius}) ; pour un acide faible dosé par une base forte, $pH_E > 7$ (solution basique à l'équivalence, hydrolyse de la base conjuguée) ; pour une base faible dosée par un acide fort, $pH_E < 7$ (solution acide, hydrolyse de l'acide conjugué).
\end{methode}

\begin{exoresolu}[Lecture graphique d'un dosage acide fort -- base forte]
On dose $V_A = \SI{20,0}{mL}$ d'une solution d'acide chlorhydrique de concentration $C_A$ inconnue par de la soude à $C_B = \SI{0,050}{mol.L^{-1}}$. Le suivi pH-métrique donne la courbe $pH = f(V_B)$ ci-dessous.

\begin{popfigure}
\begin{center}
\begin{tikzpicture}
\begin{axis}[
width=0.9\linewidth, height=7cm,
xlabel={$V_B$ (mL)}, ylabel={pH},
xmin=0, xmax=30, ymin=0, ymax=14,
xtick={0,5,10,15,20,25,30}, ytick={0,2,4,6,8,10,12,14},
grid=both, grid style={popline!40},
legend style={at={(0.03,0.97)}, anchor=north west, font=\small},
clip=false
]
% Courbe sigmoïde centrée sur 15 mL, pH 7
\addplot[popcurve, domain=0:30, samples=300] {7 + 6*tanh(1.2*(x-15))};
% Droites tangentes (approximation visuelle)
\addplot[poporange, dashed, thick, domain=0:15] {7 + 6*tanh(1.2*(7.5-15)) + (6*1.2*(1-tanh(1.2*(7.5-15))^2))*(x-7.5)};
\addplot[poporange, dashed, thick, domain=15:30] {7 + 6*tanh(1.2*(22.5-15)) + (6*1.2*(1-tanh(1.2*(22.5-15))^2))*(x-22.5)};
% Droite médiane
\addplot[poporange, dashed, thick, domain=0:30] {7 + (6*1.2*(1-tanh(1.2*(7.5-15))^2))*(x-15)};
% Point E
\addplot[only marks, mark=*, mark size=3pt, poporange] coordinates {(15,7)};
\node[poporange, anchor=south west, font=\small] at (axis cs:15.5,7.2) {$E$};
\end{axis}
\end{tikzpicture}
\end{center}
\legende{Courbe de dosage pH-métrique de l'acide chlorhydrique par la soude.}
\end{popfigure}

\begin{enumerate}
\item Déterminer graphiquement le volume équivalent $V_E$ et le pH à l'équivalence.
\item Ce résultat est-il cohérent avec la nature des réactifs ?
\end{enumerate}
\tcblower
\textbf{1. Détermination graphique.} Par la méthode des tangentes (ou en repérant le point d'inflexion au milieu du saut de pH), le point équivalent se situe aux coordonnées :
\[ V_E = \SI{15,0}{mL} \qquad pH_E = 7,0 \]

\textbf{2. Cohérence.} L'acide chlorhydrique est un acide fort et la soude une base forte. À l'équivalence, la solution ne contient que des ions spectateurs (\ce{Na+}, \ce{Cl-}) et de l'eau : la solution est \textbf{neutre}. On doit donc trouver $pH_E = 7$ à \SI{25}{\celsius}, ce qui est exactement le cas.

\textbf{Conclusion :} $V_E = \SI{15,0}{mL}$ et $pH_E = 7,0$, valeur caractéristique du dosage d'un acide fort par une base forte.
\end{exoresolu}

\begin{methode}[Exploiter la relation à l'équivalence]
\begin{enumerate}
\item \cle{Écrire} l'équation-bilan du dosage et vérifier la stœchiométrie (le plus souvent 1 pour 1).
\item \cle{Traduire} l'équivalence : les réactifs ont été introduits dans les proportions stœchiométriques. Pour \ce{AH + HO- -> A- + H2O} :
\[ n_A(\text{initial}) = n_B(\text{versé à l'équivalence}) \quad \Longleftrightarrow \quad C_A \times V_A = C_B \times V_E \]
\item \cle{Isoler} l'inconnue (souvent $C_A$) : $C_A = \dfrac{C_B \times V_E}{V_A}$.
\item \cle{Calculer} en convertissant tous les volumes dans la même unité (mL partout ou L partout) et en respectant les chiffres significatifs.
\item \cle{Conclure} par une phrase complète avec l'unité.
\end{enumerate}
\textbf{Réflexe :} si un prélèvement a été dilué avant dosage, remonter à la concentration de la solution mère en multipliant par le facteur de dilution ($C_{\text{mère}} = f_{\text{dil}} \times C_{\text{diluée}}$).
\end{methode}

\begin{exoresolu}[Vérification de l'étiquette d'un détartrant]
L'étiquette d'un détartrant commercial annonce un pourcentage massique en acide chlorhydrique de $\SI{9}{\percent}$ (densité $d = 1,04$). Pour vérifier, on dilue 10 fois le détartrant (solution $S$), puis on dose $V_A = \SI{20,0}{mL}$ de $S$ par de la soude à $C_B = \SI{0,100}{mol.L^{-1}}$. L'équivalence est obtenue pour $V_E = \SI{16,4}{mL}$.
\begin{enumerate}
\item Écrire l'équation-bilan du dosage.
\item Calculer la concentration $C_A$ de la solution diluée $S$, puis celle $C_0$ du détartrant.
\item Calculer le pourcentage massique réel et conclure. On donne $M(\ce{HCl}) = \SI{36,5}{g.mol^{-1}}$.
\end{enumerate}
\tcblower
\textbf{1. Équation-bilan.} Acide fort + base forte :
\[ \ce{H3O+ + HO- -> 2 H2O} \]

\textbf{2. Concentrations.} À l'équivalence : $C_A \times V_A = C_B \times V_E$, d'où :
\[ C_A = \frac{C_B \times V_E}{V_A} = \frac{\SI{0,100}{mol.L^{-1}} \times \SI{16,4}{mL}}{\SI{20,0}{mL}} = \SI{8,20e-2}{mol.L^{-1}} \]
Le détartrant a été dilué 10 fois (facteur de dilution $f = 10$) :
\[ C_0 = f \times C_A = 10 \times \SI{8,20e-2}{mol.L^{-1}} = \SI{0,820}{mol.L^{-1}} \]

\textbf{3. Pourcentage massique.} Un litre de détartrant a une masse $m = d \times \SI{1000}{g} = \SI{1040}{g}$ et contient :
\[ m(\ce{HCl}) = C_0 \times M = 0,820 \times 36,5 = \SI{29,9}{g} \]
\[ \%\,m = \frac{29,9}{1040} \times 100 \approx \SI{2,9}{\percent} \]

\textbf{Conclusion :} le pourcentage réel ($\SI{2,9}{\percent}$) est très inférieur à celui annoncé ($\SI{9}{\percent}$) : l'étiquette n'est \textbf{pas conforme}. (Remarque : si le fabricant exprime le pourcentage en solution commerciale concentrée diluée d'usage, il faut comparer à la notice ; ici l'écart est trop grand pour être une simple incertitude.)
\end{exoresolu}

\begin{methode}[Exploiter la demi-équivalence pour déterminer un pKa]
\begin{enumerate}
\item \cle{Repérer} le volume équivalent $V_E$ sur la courbe (méthode des tangentes).
\item \cle{Calculer} le volume de demi-équivalence : $V_{1/2} = \dfrac{V_E}{2}$.
\item \cle{Lire} le pH de la courbe pour $V_B = V_{1/2}$ : à la demi-équivalence, $[\ce{AH}] = [\ce{A-}]$, donc
\[ pH_{1/2} = pK_a \]
\item \cle{Identifier} l'acide en comparant le $pK_a$ trouvé aux valeurs tabulées.
\end{enumerate}
\textbf{Réflexe :} cette méthode ne fonctionne que pour le dosage d'un acide \emph{faible} (ou d'une base faible) ; elle est sans objet pour un acide fort, totalement dissocié.
\end{methode}

\begin{exoresolu}[Identification d'un acide faible]
Un élève dose $V_A = \SI{25,0}{mL}$ d'une solution d'un acide faible \ce{AH} par de la soude à $C_B = \SI{0,10}{mol.L^{-1}}$. La courbe $pH = f(V_B)$ donne : $V_E = \SI{12,5}{mL}$ et, pour $V_B = \SI{6,25}{mL}$, $pH = 4,8$.
\begin{enumerate}
\item Déterminer le $pK_a$ du couple \ce{AH}/\ce{A-} et identifier l'acide parmi : acide méthanoïque ($pK_a = 3,8$), acide éthanoïque ($pK_a = 4,8$), acide benzoïque ($pK_a = 4,2$).
\item Calculer la concentration $C_A$ de la solution dosée.
\end{enumerate}
\tcblower
\textbf{1. Détermination du $pK_a$.} Le volume de demi-équivalence est :
\[ V_{1/2} = \frac{V_E}{2} = \frac{12,5}{2} = \SI{6,25}{mL} \]
Or, à la demi-équivalence, $pH = pK_a$. La courbe donne $pH = 4,8$ pour $V_B = \SI{6,25}{mL}$, donc :
\[ pK_a = 4,8 \]
L'acide dosé est l'\textbf{acide éthanoïque} \ce{CH3COOH}.

\textbf{2. Concentration.} À l'équivalence : $C_A V_A = C_B V_E$ :
\[ C_A = \frac{0,10 \times 12,5}{25,0} = \SI{5,0e-2}{mol.L^{-1}} \]

\textbf{Conclusion :} l'acide est l'acide éthanoïque ($pK_a = 4,8$) et sa concentration vaut $C_A = \SI{5,0e-2}{mol.L^{-1}}$.
\end{exoresolu}

\begin{methode}[Choisir judicieusement un indicateur coloré]
\begin{enumerate}
\item \cle{Déterminer} le pH à l'équivalence $pH_E$ (lecture graphique ou raisonnement sur la nature des réactifs).
\item \cle{Comparer} $pH_E$ aux zones de virage des indicateurs disponibles :
\begin{center}
\begin{tabularx}{\linewidth}{|l|c|X|}
\hline
\rowcolor{popblueL} \textbf{Indicateur} & \textbf{Zone de virage} & \textbf{Couleurs (milieu acide / milieu basique)} \\
\hline
Hélianthine & 3,1 -- 4,4 & rouge / jaune \\
\hline
Bleu de bromothymol (BBT) & 6,0 -- 7,6 & jaune / bleu \\
\hline
Phénolphtaléine & 8,2 -- 10,0 & incolore / rose \\
\hline
\end{tabularx}
\end{center}
\item \cle{Retenir} l'indicateur dont la zone de virage \textbf{contient} $pH_E$ (ou, au minimum, est entièrement comprise dans le saut de pH).
\item \cle{Justifier} le choix en une phrase : « la zone de virage de l'indicateur encadre le pH à l'équivalence ».
\end{enumerate}
\textbf{Réflexe :} acide fort -- base forte ($pH_E = 7$) : BBT ; acide faible -- base forte ($pH_E > 7$) : phénolphtaléine ; base faible -- acide fort ($pH_E < 7$) : hélianthine (ou mieux, rouge de méthyle zone 4,2--6,2).
\end{methode}

\begin{exoresolu}[Dosage colorimétrique de l'ammoniac]
On dose une solution d'ammoniac \ce{NH3} (base faible, $pK_a(\ce{NH4+}/\ce{NH3}) = 9,2$) par de l'acide chlorhydrique. Le pH à l'équivalence vaut $pH_E = 5,5$.
\begin{enumerate}
\item Écrire l'équation-bilan du dosage.
\item Choisir, parmi l'hélianthine, le BBT et la phénolphtaléine, l'indicateur le mieux adapté. Justifier.
\end{enumerate}
\tcblower
\textbf{1. Équation-bilan.} Acide fort + base faible :
\[ \ce{H3O+ + NH3 -> NH4+ + H2O} \]

\textbf{2. Choix de l'indicateur.} À l'équivalence, la solution contient l'ion ammonium \ce{NH4+}, acide faible : la solution est acide ($pH_E = 5,5$). Comparons les zones de virage :
\begin{itemize}
\item hélianthine : 3,1 -- 4,4 : le virage survient \emph{avant} l'équivalence (zone trop acide) ;
\item BBT : 6,0 -- 7,6 : le virage survient \emph{après} l'équivalence (zone trop basique) ;
\item phénolphtaléine : 8,2 -- 10,0 : totalement inadaptée.
\end{itemize}
Aucun des trois indicateurs n'encadre parfaitement $pH_E = 5,5$. Le saut de pH pour ce dosage s'étend typiquement de pH $\approx 7$ à pH $\approx 3$. Les zones de virage de l'hélianthine et du BBT sont toutes deux \emph{à l'intérieur} du saut, mais aucune ne contient le point d'équivalence. En pratique, on utilise le \textbf{rouge de méthyle} (zone 4,2 -- 6,2) qui contient $pH_E = 5,5$. Parmi les trois proposés, l'hélianthine est souvent retenue car son virage (jaune $\to$ rouge) est très net dans la partie acide du saut.

\textbf{Conclusion :} l'indicateur adapté doit avoir une zone de virage contenant $pH_E = 5,5$ : le rouge de méthyle est le meilleur choix ; parmi les trois proposés, l'hélianthine est le moins mauvais car sa zone de virage se situe dans le saut de pH.
\end{exoresolu}

\begin{methode}[Préparer une solution tampon de pH = pKa]
\begin{enumerate}
\item \cle{Choisir} le couple acide faible / base conjuguée dont le $pK_a$ est le plus proche du pH désiré (idéalement $pH = pK_a$).
\item \cle{Respecter} la condition fondamentale : $[\ce{AH}] = [\ce{A-}]$, car alors (équation de Henderson-Hasselbalch)
\[ pH = pK_a + \log\frac{[\ce{A-}]}{[\ce{AH}]} = pK_a + \log 1 = pK_a \]
\item \cle{Préparer} par l'une des trois voies :
\begin{itemize}
\item mélanger des volumes égaux de solution d'acide \ce{AH} et de solution de base conjuguée \ce{A-} de \textbf{même concentration} ;
\item doser un acide faible par une base forte et s'\textbf{arrêter à la demi-équivalence} ($V_B = V_E/2$) ;
\item dissoudre un mélange équimolaire d'acide faible et de son sel (ex. \ce{CH3COOH} + \ce{CH3COONa}).
\end{itemize}
\item \cle{Vérifier} le pH au pH-mètre et ajuster si nécessaire par ajout goutte à goutte d'acide ou de base forte.
\end{enumerate}
\textbf{Réflexe :} une solution tampon modère les variations de pH lors de l'ajout \emph{modéré} d'acide ou de base, et lors d'une dilution modérée. Sa capacité tampon est maximale pour $pH = pK_a$.
\end{methode}

\begin{exoresolu}[Préparation d'un tampon éthanoate de pH 4,8]
Un technicien de laboratoire à la SONARA doit préparer $V = \SI{200}{mL}$ d'une solution tampon de $pH = 4,8$ à partir d'une solution d'acide éthanoïque à $C = \SI{0,20}{mol.L^{-1}}$ et d'une solution d'éthanoate de sodium \ce{CH3COONa} à la même concentration $C$. On donne $pK_a(\ce{CH3COOH}/\ce{CH3COO-}) = 4,8$.
\begin{enumerate}
\item Expliquer pourquoi ce couple convient.
\item Déterminer les volumes $V_A$ d'acide et $V_B$ de sel à mélanger.
\item Que devient le pH si l'on ajoute une goutte d'acide chlorhydrique concentré ? Justifier qualitativement.
\end{enumerate}
\tcblower
\textbf{1. Choix du couple.} Le pH désiré (4,8) est exactement égal au $pK_a$ du couple \ce{CH3COOH}/\ce{CH3COO-} : ce couple est donc idéal, car il suffit d'obtenir $[\ce{CH3COOH}] = [\ce{CH3COO-}]$ pour avoir $pH = pK_a$.

\textbf{2. Volumes à mélanger.} Les deux solutions mères ont la même concentration $C$ ; pour obtenir des quantités de matière égales $n(\ce{CH3COOH}) = n(\ce{CH3COO-})$, il faut prendre des volumes égaux :
\[ V_A = V_B = \frac{V}{2} = \SI{100}{mL} \]
Vérification : dans le mélange final, 
\[ [\ce{CH3COOH}] = [\ce{CH3COO-}] = \frac{C \times \SI{100}{mL}}{\SI{200}{mL}} = \SI{0,10}{mol.L^{-1}} \]
\[ pH = pK_a + \log\frac{[\ce{CH3COO-}]}{[\ce{CH3COOH}]} = 4,8 + \log 1 = 4,8 \]

\textbf{3. Effet d'un ajout d'acide.} Les ions \ce{H3O+} apportés réagissent totalement avec la base du tampon :
\[ \ce{CH3COO- + H3O+ -> CH3COOH + H2O} \]
La concentration de la base \ce{CH3COO-} diminue légèrement et celle de l'acide \ce{CH3COOH} augmente d'autant : le rapport $\dfrac{[\ce{CH3COO-}]}{[\ce{CH3COOH}]}$ varie très peu, donc le pH \textbf{diminue très faiblement} : c'est l'effet tampon.

\textbf{Conclusion :} en mélangeant \SI{100}{mL} de chaque solution, on obtient \SI{200}{mL} de tampon de $pH = pK_a = 4,8$, capable de modérer les variations de pH.
\end{exoresolu}

\begin{attention}[Les 5 erreurs qui coûtent des points au Bac]
\begin{enumerate}
\item \textbf{Écrire \ce{Na+} ou \ce{Cl-} dans l'équation-bilan.} Ce sont des ions spectateurs : l'équation d'un dosage s'écrit uniquement avec les espèces qui échangent le proton (\ce{H3O+ + HO- -> 2 H2O}, \ce{AH + HO- -> A- + H2O}).
\item \textbf{Confondre équivalence et demi-équivalence.} À l'équivalence : $C_A V_A = C_B V_E$ (réactifs en proportions stœchiométriques). À la demi-équivalence : $V_B = V_E/2$ et $pH = pK_a$. Ne jamais écrire $pH_E = pK_a$ !
\item \textbf{Affirmer que le pH à l'équivalence vaut toujours 7.} $pH_E = 7$ uniquement pour le dosage acide fort -- base forte (à \SI{25}{\celsius}). Acide faible dosé par une base forte : $pH_E > 7$ ; base faible dosée par un acide fort : $pH_E < 7$.
\item \textbf{Choisir l'indicateur coloré « au hasard ».} Le choix se justifie obligatoirement : la zone de virage doit \textbf{contenir le pH à l'équivalence} (ou être comprise dans le saut de pH). Une justification absente ou fausse fait perdre le point, même si l'indicateur est le bon.
\item \textbf{Oublier la dilution dans les calculs de concentration.} Quand l'échantillon a été dilué avant dosage, la concentration calculée est celle de la solution \emph{diluée} : il faut multiplier par le facteur de dilution pour retrouver la concentration de la solution commerciale. Penser aussi à convertir les volumes dans la même unité et à respecter les chiffres significatifs.
\end{enumerate}
\end{attention}

\newpage

\coursec{Exercices}

\coursec{Leçon 08 : Dosages acido-basiques et solutions tampons}

\courssub{Je vérifie mes connaissances}

\begin{exercice}[QCM de synthèse]
Pour chaque question, une seule réponse est exacte. Entoure la bonne lettre.

\textbf{1.} À l'équivalence d'un dosage acido-basique :
\begin{itemize}
\item[a)] le pH est toujours égal à 7 ;
\item[b)] les réactifs ont été mélangés dans les proportions stoechiométriques ;
\item[c)] l'indicateur coloré est toujours incolore.
\end{itemize}

\textbf{2.} Lors du dosage d'un acide faible par une base forte, le pH à l'équivalence est :
\begin{itemize}
\item[a)] inférieur à 7 ; \item[b)] égal à 7 ; \item[c)] supérieur à 7.
\end{itemize}

\textbf{3.} À la demi-équivalence du dosage d'un acide faible \ce{AH} par la soude :
\begin{itemize}
\item[a)] $pH = pK_a$ du couple \ce{AH}/\ce{A-} ; \item[b)] $pH = 7$ ; \item[c)] $pH = 14 - pK_a$.
\end{itemize}

\textbf{4.} Pour doser une solution d'ammoniac (base faible) par l'acide chlorhydrique ($pH_E \approx 5$), l'indicateur le mieux adapté est :
\begin{itemize}
\item[a)] la phénolphtaléine (zone de virage 8,2 -- 10) ; \item[b)] l'hélianthine (zone de virage 3,1 -- 4,4) ; \item[c)] le bleu de bromothymol (zone de virage 6,0 -- 7,6).
\end{itemize}

\textbf{5.} Une solution tampon :
\begin{itemize}
\item[a)] a toujours un pH égal à 7 ;
\item[b)] voit son pH varier très peu lors de l'ajout modéré d'acide, de base ou d'eau ;
\item[c)] ne contient que des espèces fortes.
\end{itemize}
\end{exercice}

\begin{exercice}[Vrai ou faux, à justifier]
Réponds par \textbf{vrai} ou \textbf{faux}, puis justifie ta réponse en une ou deux phrases.
\begin{enumerate}
\item À l'équivalence du dosage d'un acide fort par une base forte, le pH vaut 7 à \SI{25}{\celsius}.
\item L'équivalence acido-basique signifie que la solution obtenue est neutre.
\item La méthode des tangentes permet de repérer graphiquement le point équivalent sur une courbe $pH = f(V)$.
\item Pour préparer une solution tampon de $pH = pK_a$, il faut mélanger un acide faible et sa base conjuguée avec $[\text{base}] = [\text{acide}]$.
\item Diluer dix fois une solution tampon divise son pH par dix.
\end{enumerate}
\end{exercice}

\begin{exercice}[Texte à trous]
Recopie et complète le texte suivant avec les mots ou expressions de la liste : \emph{équivalence} ; \emph{tangentes} ; \emph{stoechiométriques} ; \emph{demi-équivalence} ; \emph{tampon} ; \emph{zone de virage} ; \emph{$pK_a$}.

\medskip
Lors d'un dosage, le point \ldots\ldots\ldots\ldots est le point de la courbe $pH = f(V)$ où le pH varie le plus brutalement ; on le détermine graphiquement par la méthode des \ldots\ldots\ldots\ldots. À ce point, les réactifs ont été mélangés dans les proportions \ldots\ldots\ldots\ldots. Au point de \ldots\ldots\ldots\ldots, la moitié de l'acide faible a réagi et le pH est égal au \ldots\ldots\ldots du couple. Pour un dosage colorimétrique, on choisit un indicateur dont la \ldots\ldots\ldots\ldots encadre le pH à l'équivalence. Un mélange équimolaire d'un acide faible et de sa base conjuguée constitue une solution \ldots\ldots\ldots\ldots.
\end{exercice}

\begin{exercice}[Définitions et relations]
\begin{enumerate}
\item Définis l'\cle{équivalence acido-basique}.
\item Écris la relation entre les quantités de matière à l'équivalence pour le dosage d'un acide (concentration $C_a$, volume $V_a$) par une base (concentration $C_b$, volume équivalent $V_{bE}$).
\item Donne la définition d'une \cle{solution tampon} et cite deux de ses propriétés.
\item Écris l'équation-bilan de la réaction entre l'acide éthanoïque \ce{CH3COOH} et l'hydroxyde de sodium \ce{NaOH}.
\end{enumerate}
\end{exercice}

\courssub{Je m'entraîne}

\begin{aretenir}[Relation à l'équivalence]
Pour une réaction acido-basique $\nu_a \ce{AH} + \nu_b \ce{B} \rightarrow \text{produits}$, à l'équivalence :
\[
\frac{n_{a,0}}{\nu_a} = \frac{n_{b,E}}{\nu_b} \quad \text{soit} \quad \frac{C_a V_a}{\nu_a} = \frac{C_b V_{bE}}{\nu_b}
\]
Cette relation permet de calculer une concentration inconnue.
\end{aretenir}

\begin{methode}[Méthode des tangentes (détermination graphique de $V_E$)]
\begin{enumerate}
\item Tracer la courbe $pH = f(V)$ sur papier millimétré avec des échelles adaptées.
\item Tracer la tangente à la branche montante (avant le saut) et la tangente à la branche plate (après le saut).
\item Tracer la bissectrice de l'angle formé par ces deux tangentes (ou tracer les deux tangentes et trouver leur intersection).
\item L'abscisse du point d'intersection (ou de la bissectrice avec la courbe) donne le volume équivalent $V_E$. L'ordonnée donne $pH_E$.
\end{enumerate}
\end{methode}

\begin{experience}[Dosage pH-métrique : dispositif et protocole]
\begin{center}
\begin{tikzpicture}[scale=0.7, every node/.style={transform shape}, >=Stealth]
% Support
\draw[thick, popdark] (0,0) -- (0,6) -- (4,6) -- (4,5.5);
% Burette
\draw[thick, popblue] (4,5.5) -- (4,4.5) -- (5,4.5) -- (5,1) -- (4,1) -- (4,0.5);
\draw[thick, popblue] (3.5,4.5) -- (5,4.5); \draw[thick, popblue] (3.5,1) -- (5,1);
\foreach \y in {1,2,3,4} \draw[popblue] (3.7,\y) -- (4.3,\y);
\node[popblue, right] at (5.2, 2.7) {Burette (titrant)};
% Robinet
\draw[thick, popblue] (4,0.5) -- (4,0);
\draw[thick, popblue, rotate=30] (4,0) -- ++(0.4,0);
% Bécher
\draw[thick, popdark] (1.5,0) -- (1,-0.5) -- (3,-0.5) -- (2.5,0);
\draw[thick, popdark] (1,-0.5) -- (1,-3) -- (3,-3) -- (3,-0.5);
% Barre aimantée
\draw[poporange, thick] (2,-2) ellipse (0.4 and 0.1);
% Sonde pH
\draw[thick, poppurple] (2,0.5) -- (2,-2.5);
\draw[thick, poppurple] (1.7,-2.5) -- (2.3,-2.5);
\node[poppurple, right] at (2.5, -1) {Sonde pH};
% pH-mètre
\draw[thick, popdark] (5.5,-1.5) rectangle (7.5,-3);
\node at (6.5, -2.25) {pH-mètre};
\draw[thick, poppurple] (2,0.5) .. controls (3,1.5) and (5,1.5) .. (6.5,-1.5);
% Agitateur
\draw[thick, popgray] (0.5,-3.5) rectangle (3.5,-3);
\node at (2, -3.25) {Agitateur magnétique};
% Goutte
\draw[popblue, thick] (4,0) ellipse (0.15 and 0.3);
\end{tikzpicture}
\end{center}
\legende{Dispositif de dosage pH-métrique : la burette verse le titrant dans le bécher contenant le titré, la sonde pH reliée au pH-mètre suit l'évolution du pH sous agitation magnétique.}
\end{experience}

\begin{exercice}[Exploitation de courbes anonymes]
On a réalisé trois dosages pH-métriques à \SI{25}{\celsius} et tracé les courbes $pH = f(V)$ ci-dessous (sens de l'ajout : réactif versé dans la burette). Les caractéristiques lues sur les courbes sont :

\begin{center}
\begin{tabularx}{\linewidth}{|l|X|X|X|}
\hline
\rowcolor{popblueL} Courbe & pH initial & $pH_E$ & Allure du saut \\
\hline
(1) & 2,9 & 8,7 & saut marqué, une seule branche \\
\hline
(2) & 1,0 & 7,0 & saut très marqué \\
\hline
(3) & 11,1 & 5,3 & courbe décroissante \\
\hline
\end{tabularx}
\end{center}

\begin{enumerate}
\item Associe à chaque courbe le dosage correspondant parmi : dosage d'un acide fort par une base forte ; dosage d'un acide faible par une base forte ; dosage d'une base faible par un acide fort. Justifie chaque choix à l'aide du pH initial et de $pH_E$.
\item Pour chaque dosage, propose un indicateur coloré adapté parmi : hélianthine (3,1 -- 4,4), BBT (6,0 -- 7,6), phénolphtaléine (8,2 -- 10).
\end{enumerate}
\end{exercice}

\begin{exercice}[Calcul de concentration à l'équivalence]
On dose $V_a = \SI{20,0}{mL}$ d'une solution d'acide chlorhydrique \ce{(H3O+ + Cl-)} par une solution d'hydroxyde de sodium de concentration $C_b = \SI{0,10}{mol.L^{-1}}$. Le suivi pH-métrique donne un volume équivalent $V_{bE} = \SI{15,0}{mL}$.
\begin{enumerate}
\item Écris l'équation-bilan de la réaction de dosage.
\item Écris la relation à l'équivalence et calcule la concentration molaire $C_a$ de la solution d'acide chlorhydrique.
\item Calcule la concentration massique $C_m$ de cette solution. Donnée : $M(\ce{HCl}) = \SI{36,5}{g.mol^{-1}}$.
\item Quel est le pH de la solution à l'équivalence ? Justifie.
\end{enumerate}
\end{exercice}

\begin{exercice}[Détermination d'un $pK_a$]
On dose $V_a = \SI{25,0}{mL}$ d'une solution d'acide méthanoïque \ce{HCOOH} de concentration $C_a = \SI{0,080}{mol.L^{-1}}$ par la soude à $C_b = \SI{0,10}{mol.L^{-1}}$.
\begin{enumerate}
\item Calcule le volume équivalent $V_{bE}$ attendu.
\item Sur la courbe tracée, on lit un pH égal à 3,8 à la demi-équivalence. En déduire le $pK_a$ du couple \ce{HCOOH}/\ce{HCOO-}.
\item Calcule la constante d'acidité $K_a$ de ce couple.
\item Calcule le volume de soude versé à la demi-équivalence.
\end{enumerate}
\end{exercice}

\begin{savaistu}[L'acide méthanoïque et les fourmis]
L'acide méthanoïque (\ce{HCOOH}) tire son nom du latin \emph{formica} (fourmi). Il est présent dans le venin de certaines fourmis et dans les orties (piqûres). C'est l'acide carboxylique le plus simple. Au Cameroun, on le retrouve à l'état de traces dans certains miels ou jus de fruits fermentés. Son $pK_a$ (3,75--3,8) en fait un acide faible typique, bien plus fort que l'acide éthanoïque ($pK_a=4,8$).
\end{savaistu}

\begin{exercice}[Préparation d'une solution tampon]
Au laboratoire du lycée, on dispose d'une solution d'acide éthanoïque \ce{CH3COOH} à \SI{0,20}{mol.L^{-1}} et d'une solution d'éthanoate de sodium \ce{CH3COONa} à \SI{0,20}{mol.L^{-1}}. On veut préparer $V = \SI{100}{mL}$ d'une solution tampon de $pH = 4,8$. Donnée : $pK_a(\ce{CH3COOH}/\ce{CH3COO-}) = 4,8$.
\begin{enumerate}
\item Quelle relation lie $pH$, $pK_a$, $[\ce{CH3COOH}]$ et $[\ce{CH3COO-}]$ dans ce mélange ?
\item Montre qu'il faut mélanger des volumes égaux des deux solutions.
\item Donne les volumes $V_1$ et $V_2$ à prélever et décris le mode opératoire.
\item Qu'advient-il du pH si l'on ajoute quelques gouttes d'acide chlorhydrique dilué à ce tampon ? Justifie par une équation-bilan.
\end{enumerate}
\end{exercice}

\begin{methode}[Préparation d'une solution tampon de pH = pKa]
\begin{enumerate}
\item Choisir un couple acide/base conjuguée dont $pK_a \approx pH$ visé.
\item Préparer deux solutions mères de même concentration $C$ : l'acide faible \ce{AH} et sa base conjuguée \ce{A-} (souvent sel de sodium/potassium).
\item Calculer le volume total $V$ souhaité. Prendre $V/2$ de chaque solution mère.
\item Mélanger dans un flacon jaugé de volume $V$, compléter à l'eau distillée au trait de jauge, homogénéiser.
\item Vérifier le pH au pH-mètre étalonné.
\end{enumerate}
\end{methode}

\courssub{Je me prépare au Bac}

\begin{exercice}[Vérification de l'étiquette d'une solution d'acide chlorhydrique]
Un laboratoire de Douala reçoit un flacon d'acide chlorhydrique dont l'étiquette indique : « solution à \SI{0,10}{mol.L^{-1}} ». Pour vérifier cette indication, un technicien prélève $V_a = \SI{20,0}{mL}$ de la solution commerciale, les introduit dans un bécher et dose par une solution d'hydroxyde de sodium de concentration $C_b = \SI{0,100}{mol.L^{-1}}$, en suivant le pH après chaque ajout. Il obtient les résultats suivants :

\begin{center}
\begin{tabularx}{\linewidth}{|l|X|X|X|X|X|X|X|X|X|}
\hline
\rowcolor{popblueL} $V_b$ (mL) & 0 & 2 & 4 & 6 & 8 & 9 & 9,5 & 10 & 10,5 \\
\hline
pH & 1,1 & 1,2 & 1,3 & 1,5 & 1,9 & 2,4 & 2,9 & 7,0 & 11,1 \\
\hline
\end{tabularx}
\end{center}
\begin{center}
\begin{tabularx}{\linewidth}{|l|X|X|X|X|X|X|}
\hline
\rowcolor{popblueL} $V_b$ (mL) & 11 & 12 & 14 & 16 & 18 & 20 \\
\hline
pH & 11,4 & 11,7 & 12,0 & 12,2 & 12,3 & 12,4 \\
\hline
\end{tabularx}
\end{center}

\begin{enumerate}
\item Fais le schéma légendé du dispositif de dosage pH-métrique.
\item Écris l'équation-bilan de la réaction de dosage et donne ses caractéristiques (totale, rapide, exothermique).
\item Trace la courbe $pH = f(V_b)$ sur papier millimétré. Échelle : \SI{1}{cm} pour \SI{2}{mL} en abscisse et \SI{1}{cm} pour 1 unité de pH en ordonnée.
\item Détermine graphiquement les coordonnées du point équivalent $E$ par la méthode des tangentes.
\item Justifie la valeur du pH à l'équivalence.
\item Calcule la concentration molaire $C_a$ de la solution commerciale, puis sa concentration massique $C_m$. Donnée : $M(\ce{HCl}) = \SI{36,5}{g.mol^{-1}}$.
\item L'indication de l'étiquette est-elle exacte ? Conclus.
\item Parmi l'hélianthine (3,1 -- 4,4), le BBT (6,0 -- 7,6) et la phénolphtaléine (8,2 -- 10), quel(s) indicateur(s) aurait-on pu utiliser pour un dosage colorimétrique ? Justifie.
\end{enumerate}
\end{exercice}

\begin{exercice}[Dosage de l'ammoniac d'un produit ménager]
Une solution commerciale d'ammoniac \ce{NH3}, utilisée comme produit de nettoyage, affiche une concentration d'environ \SI{1}{mol.L^{-1}}. Pour la contrôler, on dilue dix fois la solution commerciale (solution $S$), puis on dose $V_b = \SI{20,0}{mL}$ de $S$ par une solution d'acide chlorhydrique de concentration $C_a = \SI{0,10}{mol.L^{-1}}$. La courbe $pH = f(V_a)$ fait apparaître un point équivalent $E(V_{aE} = \SI{19,6}{mL}\,;\, pH_E = 5,4)$. Données : $pK_a(\ce{NH4+}/\ce{NH3}) = 9,2$ ; $M(\ce{NH3}) = \SI{17,0}{g.mol^{-1}}$.

\begin{enumerate}
\item Pourquoi a-t-on dilué la solution commerciale avant le dosage ?
\item Écris l'équation-bilan de la réaction de dosage.
\item Explique, à l'aide des espèces présentes à l'équivalence, pourquoi $pH_E < 7$.
\item Calcule la concentration molaire $C_b$ de la solution diluée $S$, puis celle $C_0$ de la solution commerciale. L'affichage du fabricant est-il correct ?
\item Calcule la concentration massique de la solution commerciale.
\item Quel indicateur coloré faut-il choisir pour un dosage colorimétrique : hélianthine (3,1 -- 4,4), BBT (6,0 -- 7,6) ou phénolphtaléine (8,2 -- 10) ? Justifie.
\item À la demi-équivalence, quelle relation existe-t-il entre $[\ce{NH4+}]$ et $[\ce{NH3}]$ ? Quel est alors le pH de la solution ? Quel nom donne-t-on à une telle solution ?
\end{enumerate}
\end{exercice}

\begin{savaistu}[Ammoniac et lessive traditionnelle]
Au Cameroun, la lessive artisanale (savon noir, \emph{ndolé} lessive) utilise souvent de la potasse (carbonate de potassium) ou de l'ammoniac pour dégraisser. L'ammoniac commercial (\emph{eau d'ammoniac}) est une solution de \ce{NH3} dans l'eau, typiquement 10 à 30 \% massique (environ 5 à 15 M). Le dosage par l'acide chlorhydrique est le contrôle qualité standard à la SONARA ou dans les savonneries industrielles (CIMENCAM pour le nettoyage des fours). L'odeur caractéristique (piquante) est due à \ce{NH3} gazeux en équilibre.
\end{savaistu}

\begin{exercice}[Acide éthanoïque du vinaigre et préparation d'un tampon]
Le vinaigre est une solution aqueuse d'acide éthanoïque \ce{CH3COOH}. Un groupe d'élèves de Terminale TI dose $V_a = \SI{10,0}{mL}$ de vinaigre dilué dix fois par la soude à $C_b = \SI{0,10}{mol.L^{-1}}$. La courbe pH-métrique obtenue donne : point équivalent $E(V_{bE} = \SI{12,0}{mL}\,;\, pH_E = 8,4)$ et pH à la demi-équivalence égal à 4,8. Données : $M(\ce{CH3COOH}) = \SI{60,0}{g.mol^{-1}}$ ; produit ionique de l'eau $K_e = 10^{-14}$ à \SI{25}{\celsius}.

\begin{enumerate}
\item Écris l'équation-bilan de la réaction de dosage.
\item Détermine le $pK_a$ du couple \ce{CH3COOH}/\ce{CH3COO-} en justifiant ta démarche, puis calcule $K_a$.
\item Justifie que $pH_E > 7$ en identifiant l'espèce prédominante à l'équivalence.
\item Calcule la concentration molaire $C_a$ du vinaigre dilué, puis celle du vinaigre commercial.
\item Le degré d'acidité d'un vinaigre est la masse d'acide éthanoïque (en grammes) contenue dans \SI{100}{g} de vinaigre (on assimilera la masse volumique du vinaigre à \SI{1,0}{g.mL^{-1}}). Calcule le degré d'acidité de ce vinaigre.
\item Les élèves veulent préparer \SI{200}{mL} d'une solution tampon de $pH = 4,8$ à partir de solutions de \ce{CH3COOH} et de \ce{CH3COONa}, toutes deux à \SI{0,10}{mol.L^{-1}}. Calcule les volumes à mélanger.
\item Cite deux propriétés de la solution tampon obtenue et une application concrète des solutions tampons dans la vie courante ou en biologie.
\end{enumerate}
\end{exercice}

\begin{savaistu}[Le vinaigre au Cameroun : du vin de palme à la conservation]
Le vinaigre traditionnel camerounais s'obtient par double fermentation du vin de palme (vin $\xrightarrow{\text{levures}}$ alcool $\xrightarrow{\text{acétobacters}}$ acide éthanoïque). Le degré d'acidité (5--8\%) assure la conservation des aliments (pickles, condiments) en inhibant les bactéries pathogènes (pH < 4,5). Les savonneries artisanales (région Ouest, Nord) utilisent parfois l'acide éthanoïque pour neutraliser l'excès de soude dans le savon (réaction acide-base). Le contrôle du degré d'acidité par dosage est une compétence clé en agroalimentaire (normes CODEX/ANOR).
\end{savaistu}

\begin{aretenir}[Résumé : Choix de l'indicateur coloré]
Pour un dosage colorimétrique, l'indicateur doit changer de couleur \textbf{dans la zone de saut de pH} à l'équivalence.
\begin{itemize}
\item \textbf{Acide fort / Base forte} : Saut large (pH 3--11). \textbf{BBT} (6,0--7,6) ou \textbf{Phénolphtaléine} (8,2--10).
\item \textbf{Acide faible / Base forte} : $pH_E > 7$ (ex: 8--10). Saut dans le basique. \textbf{Phénolphtaléine} (8,2--10).
\item \textbf{Base faible / Acide fort} : $pH_E < 7$ (ex: 3--6). Saut dans l'acide. \textbf{Hélianthine} (3,1--4,4) ou \textbf{Rouge de méthyle} (4,2--6,3).
\end{itemize}
Ne jamais utiliser un indicateur dont la zone de virage est en dehors du saut (erreur systématique).
\end{aretenir}

\begin{methode}[Calcul du degré d'acidité (vinaigre)]
\begin{enumerate}
\item Déterminer $C_a$ (mol/L) du vinaigre commercial par dosage.
\item Calculer $C_m = C_a \times M(\ce{CH3COOH})$ (g/L).
\item Comme $\rho \approx 1$ g/mL, \SI{1}{L} = \SI{1000}{g}.
\item Masse d'acide dans \SI{100}{g} = $\frac{C_m}{10}$.
\item Degré d'acidité (\% massique) = $\frac{C_m}{10}$.
\end{enumerate}
Exemple : $C_m = 72$ g/L $\rightarrow$ Degré = 7,2\%.
\end{methode}

\newpage

\coursec{Corrigés des exercices}

\courssub{Corrigés détaillés}

\begin{corrige}[Exercice 1 — QCM de synthèse]
\begin{enumerate}
\item \textbf{b)} L'équivalence correspond au mélange des réactifs en proportions stoechiométriques. Le pH n'est égal à 7 que pour un dosage acide fort / base forte ; un indicateur coloré change de teinte, il ne devient pas incolore.
\item \textbf{c)} À l'équivalence, seule la base conjuguée \ce{A-} est présente ; elle réagit avec l'eau : \ce{A- + H2O <=> AH + HO-}. Des ions \ce{HO-} apparaissent, donc $pH_E > 7$.
\item \textbf{a)} À la demi-équivalence, $n_{\ce{AH}} = n_{\ce{A-}}$, donc $[\ce{AH}] = [\ce{A-}]$ (même volume de solution). La relation de Henderson-Hasselbalch donne : $pH = pK_a + \log\dfrac{[\ce{A-}]}{[\ce{AH}]} = pK_a + \log 1 = pK_a$.
\item \textbf{b)} Pour le dosage base faible / acide fort, $pH_E \approx 5 < 7$ et le saut de pH s'étend environ de 3 à 7. Seule l'hélianthine (zone de virage 3,1 -- 4,4) vire dans le saut. Le BBT et la phénolphtaléine vireraient avant l'équivalence.
\item \textbf{b)} C'est la définition même d'une solution tampon : son pH varie très peu par ajout modéré d'acide, de base ou par dilution.
\end{enumerate}
\end{corrige}

\begin{corrige}[Exercice 2 — Vrai ou faux, à justifier]
\begin{enumerate}
\item \textbf{Vrai.} À l'équivalence d'un dosage acide fort / base forte, il ne reste que de l'eau et des ions spectateurs (\ce{Na+}, \ce{Cl-}) qui ne réagissent pas avec l'eau. D'où $[\ce{H3O+}] = [\ce{HO-}] = \SI{1,0e-7}{mol.L^{-1}}$ et $pH = 7$ à \SI{25}{\celsius}.
\item \textbf{Faux.} L'équivalence signifie que les réactifs ont été mélangés en proportions stoechiométriques, pas que la solution est neutre. Pour un acide faible dosé par une base forte, $pH_E > 7$ (hydrolyse de la base conjuguée formée).
\item \textbf{Vrai.} Le point équivalent est le point d'inflexion de la courbe $pH = f(V)$. La méthode des tangentes (deux tangentes parallèles de part et d'autre du saut, puis droite équidistante) permet de repérer $V_E$ graphiquement.
\item \textbf{Vrai.} D'après Henderson-Hasselbalch, $pH = pK_a + \log\dfrac{[\text{base}]}{[\text{acide}]}$. Si $[\text{base}] = [\text{acide}]$, alors $\log 1 = 0$ et $pH = pK_a$.
\item \textbf{Faux.} La dilution divise par dix les deux concentrations, donc le rapport $[\text{base}]/[\text{acide}]$ est inchangé : le pH reste pratiquement constant. C'est précisément une propriété des solutions tampons.
\end{enumerate}
\end{corrige}

\begin{corrige}[Exercice 3 — Texte à trous]
Lors d'un dosage, le point d'\textbf{équivalence} est le point de la courbe $pH = f(V)$ où le pH varie le plus brutalement ; on le détermine graphiquement par la méthode des \textbf{tangentes}. À ce point, les réactifs ont été mélangés dans les proportions \textbf{stoechiométriques}. Au point de \textbf{demi-équivalence}, la moitié de l'acide faible a réagi et le pH est égal au \textbf{$pK_a$} du couple. Pour un dosage colorimétrique, on choisit un indicateur dont la \textbf{zone de virage} encadre le pH à l'équivalence. Un mélange équimolaire d'un acide faible et de sa base conjuguée constitue une solution \textbf{tampon}.
\end{corrige}

\begin{corrige}[Exercice 4 — Définitions et relations]
\begin{enumerate}
\item L'\cle{équivalence acido-basique} est l'état du mélange pour lequel l'acide et la base ont été introduits dans les proportions stoechiométriques de l'équation-bilan de la réaction de dosage : chaque réactif a alors été entièrement consommé.
\item Pour une réaction 1/1 : $n_a = n_b$, c'est-à-dire
\[ C_a \times V_a = C_b \times V_{bE} \]
Plus généralement, pour $\nu_a \ce{AH} + \nu_b \ce{B} \ce{->}$ produits : $\dfrac{C_a V_a}{\nu_a} = \dfrac{C_b V_{bE}}{\nu_b}$.
\item Une \cle{solution tampon} est une solution dont le pH varie très peu lors de l'ajout modéré d'un acide ou d'une base, ou lors d'une dilution. Propriétés : (1) le pH est quasi invariant par ajout modéré d'acide ou de base ; (2) le pH est quasi invariant par dilution ; (3) le pouvoir tampon est maximal lorsque $pH = pK_a$.
\item \[ \ce{CH3COOH + HO- -> CH3COO- + H2O} \]
(réaction totale, rapide, exothermique ; on n'écrit pas \ce{Na+}, ion spectateur).
\end{enumerate}
\end{corrige}

\begin{corrige}[Exercice 5 — Exploitation de courbes anonymes]
\begin{enumerate}
\item \textbf{Courbe (1) : dosage d'un acide faible par une base forte.} Le pH initial de 2,9 est trop élevé pour un acide fort de concentration usuelle (un acide fort à \SI{0,1}{mol.L^{-1}} donnerait pH = 1) : l'acide est donc faible, partiellement dissocié. De plus $pH_E = 8,7 > 7$, caractéristique de la présence de la base conjuguée \ce{A-} à l'équivalence.

\textbf{Courbe (2) : dosage d'un acide fort par une base forte.} Le pH initial de 1,0 correspond à un acide fort totalement dissocié ($pH = -\log C_a$ avec $C_a = \SI{0,1}{mol.L^{-1}}$) et $pH_E = 7,0$ : à l'équivalence, seuls des ions spectateurs et l'eau subsistent.

\textbf{Courbe (3) : dosage d'une base faible par un acide fort.} Le pH initial de 11,1 indique une base (si elle était forte à \SI{0,1}{mol.L^{-1}}, on aurait pH = 13 ; c'est donc une base faible). La courbe décroissante confirme que l'on verse un acide, et $pH_E = 5,3 < 7$ traduit l'hydrolyse de l'acide conjugué formé (\ce{BH+}).
\item Choix des indicateurs (la zone de virage doit être incluse dans le saut de pH) :
\begin{itemize}
\item Courbe (1), $pH_E = 8,7$ : \textbf{phénolphtaléine} (8,2 -- 10), la seule qui vire dans le saut basique.
\item Courbe (2), $pH_E = 7,0$ : le saut est très large (environ 3 à 11) ; le \textbf{BBT} (6,0 -- 7,6) convient parfaitement, la \textbf{phénolphtaléine} aussi.
\item Courbe (3), $pH_E = 5,3$ : \textbf{hélianthine} (3,1 -- 4,4), seule zone de virage incluse dans le saut acide.
\end{itemize}
\end{enumerate}
\textbf{Point de vigilance :} ne jamais choisir un indicateur dont la zone de virage est hors du saut de pH : le changement de teinte surviendrait avant ou après l'équivalence, ce qui fausserait le volume repéré.
\end{corrige}

\begin{corrige}[Exercice 6 — Calcul de concentration à l'équivalence]
\begin{enumerate}
\item Équation-bilan (ions spectateurs \ce{Na+} et \ce{Cl-} omis) :
\[ \ce{H3O+ + HO- -> 2 H2O} \]
\item À l'équivalence, les réactifs sont en proportions stoechiométriques (réaction 1/1) :
\[ n_{\ce{H3O+},0} = n_{\ce{HO-},E} \iff C_a V_a = C_b V_{bE} \]
\[ C_a = \frac{C_b \times V_{bE}}{V_a} = \frac{0,10 \times 15,0}{20,0} = \SI{0,075}{mol.L^{-1}} \]
\item Concentration massique :
\[ C_m = C_a \times M(\ce{HCl}) = 0,075 \times 36,5 = \SI{2,74}{g.L^{-1}} \]
\item À l'équivalence, la solution ne contient que de l'eau et les ions \ce{Na+} et \ce{Cl-}, qui ne réagissent pas avec l'eau (conjuguant d'une base forte et d'un acide fort). La solution est donc neutre : $pH_E = 7$ à \SI{25}{\celsius}.
\end{enumerate}
\textbf{Point de vigilance :} les volumes peuvent rester en mL dans la relation $C_aV_a = C_bV_{bE}$ à condition d'utiliser la même unité des deux côtés. Ne pas convertir seulement l'un des deux volumes en litres !
\end{corrige}

\begin{corrige}[Exercice 7 — Détermination d'un $pK_a$]
\begin{enumerate}
\item À l'équivalence : $C_a V_a = C_b V_{bE}$, donc
\[ V_{bE} = \frac{C_a V_a}{C_b} = \frac{0,080 \times 25,0}{0,10} = \SI{20,0}{mL} \]
\item À la demi-équivalence, la moitié de l'acide \ce{HCOOH} a été transformée en \ce{HCOO-} : $[\ce{HCOOH}] = [\ce{HCOO-}]$. La relation de Henderson-Hasselbalch donne alors :
\[ pH = pK_a + \log\frac{[\ce{HCOO-}]}{[\ce{HCOOH}]} = pK_a + \log 1 = pK_a \]
On lit $pH_{1/2E} = 3,8$, donc $\boxed{pK_a = 3,8}$ pour le couple \ce{HCOOH}/\ce{HCOO-}.
\item Constante d'acidité :
\[ K_a = 10^{-pK_a} = 10^{-3,8} \approx \SI{1,6e-4}{} \]
\item Le volume versé à la demi-équivalence est la moitié du volume équivalent :
\[ V_{1/2E} = \frac{V_{bE}}{2} = \frac{20,0}{2} = \SI{10,0}{mL} \]
\end{enumerate}
\textbf{Point de vigilance :} la propriété $pH = pK_a$ à la demi-équivalence n'est valable que pour le dosage d'un acide \emph{faible} (ou d'une base faible). Pour un acide fort, il n'existe pas de palier de demi-équivalence exploitable.
\end{corrige}

\begin{corrige}[Exercice 8 — Préparation d'une solution tampon]
\begin{enumerate}
\item Relation de Henderson-Hasselbalch pour le couple \ce{CH3COOH}/\ce{CH3COO-} :
\[ pH = pK_a + \log\frac{[\ce{CH3COO-}]}{[\ce{CH3COOH}]} \]
\item On veut $pH = 4,8 = pK_a$. Il faut donc $\log\dfrac{[\ce{CH3COO-}]}{[\ce{CH3COOH}]} = 0$, c'est-à-dire $[\ce{CH3COO-}] = [\ce{CH3COOH}]$ dans le mélange final. Après mélange, les concentrations sont $[\ce{CH3COOH}] = \dfrac{0,20\,V_1}{V_1+V_2}$ et $[\ce{CH3COO-}] = \dfrac{0,20\,V_2}{V_1+V_2}$. Comme les deux solutions ont la même concentration, l'égalité des concentrations finales impose $V_1 = V_2$.
\item $V_1 = V_2 = \dfrac{100}{2} = \SI{50,0}{mL}$.

\textbf{Mode opératoire :} prélever \SI{50,0}{mL} de solution d'acide éthanoïque à la pipette jaugée et les verser dans une fiole jaugée de \SI{100}{mL} ; prélever de même \SI{50,0}{mL} de solution d'éthanoate de sodium et les ajouter ; homogénéiser ; vérifier le pH au pH-mètre étalonné.
\item Le pH varie très peu. Les ions \ce{H3O+} apportés sont consommés par la base du tampon selon la réaction totale :
\[ \ce{CH3COO- + H3O+ -> CH3COOH + H2O} \]
La quantité de \ce{CH3COO-} diminue légèrement et celle de \ce{CH3COOH} augmente légèrement : le rapport $\dfrac{[\ce{CH3COO-}]}{[\ce{CH3COOH}]}$ ne change que faiblement, donc le pH reste quasi constant.
\end{enumerate}
\textbf{Point de vigilance :} un tampon n'est efficace que pour des ajouts \emph{modérés} ; si l'on ajoute une quantité d'acide supérieure à la quantité de \ce{CH3COO-} présente, le tampon est « rompu » et le pH chute brutalement.
\end{corrige}

\begin{corrige}[Exercice 9 — Vérification de l'étiquette d'une solution d'acide chlorhydrique]
\begin{enumerate}
\item Schéma : burette graduée contenant la soude (\ce{Na+} + \ce{HO-}) à \SI{0,100}{mol.L^{-1}}, fixée sur un support au-dessus d'un bécher contenant les \SI{20,0}{mL} d'acide chlorhydrique et un barreau aimanté ; la sonde du pH-mètre plonge dans le bécher ; l'ensemble repose sur un agitateur magnétique. (Se reporter au schéma du cours.)
\item Équation-bilan :
\[ \ce{H3O+ + HO- -> 2 H2O} \]
Caractéristiques : réaction \textbf{totale} (constante d'équilibre $K = 10^{14} \gg 1$), \textbf{rapide} (quasi instantanée) et \textbf{exothermique} : elle peut servir de réaction de dosage.
\item Courbe : on reporte les points $(V_b\,;\,pH)$ du tableau. La courbe croît lentement de pH 1,1 à 2,9, présente un saut quasi vertical autour de $V_b = \SI{10}{mL}$ (le pH passe de 2,9 à 11,1 entre 9,5 et \SI{10,5}{mL}), puis un palier vers pH 12,4.
\item Méthode des tangentes : on trace deux tangentes parallèles de part et d'autre du saut, puis la droite équidistante qui coupe la courbe au point équivalent. On lit :
\[ E\,(V_{bE} = \SI{10,0}{mL}\,;\, pH_E = 7,0) \]
\item À l'équivalence, la solution ne contient que de l'eau et les ions spectateurs \ce{Na+} et \ce{Cl-}, sans action sur l'eau : la solution est neutre, d'où $pH_E = 7$ à \SI{25}{\celsius}. C'est cohérent avec la lecture graphique.
\item À l'équivalence : $C_a V_a = C_b V_{bE}$, donc
\[ C_a = \frac{C_b V_{bE}}{V_a} = \frac{0,100 \times 10,0}{20,0} = \SI{0,050}{mol.L^{-1}} \]
\[ C_m = C_a \times M(\ce{HCl}) = 0,050 \times 36,5 = \SI{1,83}{g.L^{-1}} \]
\item L'étiquette annonce \SI{0,10}{mol.L^{-1}} alors que la concentration mesurée est \SI{0,050}{mol.L^{-1}}, soit deux fois moins. \textbf{L'indication de l'étiquette est fausse} (la solution a probablement été diluée ou mal préparée).
\item Le saut de pH s'étend environ de 3 à 11. Les indicateurs dont la zone de virage est incluse dans ce saut conviennent : le \textbf{BBT} (6,0 -- 7,6) et la \textbf{phénolphtaléine} (8,2 -- 10). L'hélianthine (3,1 -- 4,4) virerait dès le début du saut : elle est déconseillée.
\end{enumerate}
\textbf{Barème indicatif (sur 10) :} schéma légendé 1 ; équation + caractéristiques 1,5 ; courbe soignée 2 ; tangentes et $E$ 1,5 ; justification $pH_E$ 1 ; calculs $C_a$ et $C_m$ 2 ; conclusion et indicateur 1.

\textbf{Points de vigilance :} bien lire $V_{bE}$ sur la courbe et non dans le tableau ; ne pas confondre concentration molaire et concentration massique ; vérifier l'homogénéité des unités de volume.
\end{corrige}

\begin{corrige}[Exercice 10 — Dosage de l'ammoniac d'un produit ménager]
\begin{enumerate}
\item La solution commerciale ($\approx \SI{1}{mol.L^{-1}}$) est environ dix fois plus concentrée que l'acide titrant. Sans dilution, le volume équivalent serait d'environ \SI{200}{mL}, très supérieur à la capacité d'une burette (\SI{50}{mL}). La dilution au dixième ramène $V_{aE}$ à environ \SI{20}{mL}, mesurable avec précision.
\item Équation-bilan :
\[ \ce{NH3 + H3O+ -> NH4+ + H2O} \]
(réaction totale et rapide, utilisable pour un dosage).
\item À l'équivalence, tout l'ammoniac a été converti en ion ammonium \ce{NH4+}. Or \ce{NH4+} est l'acide conjugué de la base faible \ce{NH3} ($pK_a = 9,2$) : c'est un acide faible qui réagit partiellement avec l'eau :
\[ \ce{NH4+ + H2O <=> NH3 + H3O+} \]
Il se forme des ions \ce{H3O+}, donc $pH_E < 7$ (ici 5,4).
\item À l'équivalence : $C_b V_b = C_a V_{aE}$, donc
\[ C_b = \frac{C_a V_{aE}}{V_b} = \frac{0,10 \times 19,6}{20,0} = \SI{0,098}{mol.L^{-1}} \]
La solution commerciale est dix fois plus concentrée :
\[ C_0 = 10 \times C_b = \SI{0,98}{mol.L^{-1}} \approx \SI{1,0}{mol.L^{-1}} \]
\textbf{L'affichage du fabricant est correct} (écart de 2 \%, compatible avec les incertitudes expérimentales).
\item Concentration massique :
\[ C_m = C_0 \times M(\ce{NH3}) = 0,98 \times 17,0 = \SI{16,7}{g.L^{-1}} \]
\item $pH_E = 5,4$ : le saut de pH se situe dans la zone acide (environ 3 à 7). Seule l'\textbf{hélianthine} (zone de virage 3,1 -- 4,4) convient. Le BBT (6,0 -- 7,6) et la phénolphtaléine (8,2 -- 10) vireraient nettement avant l'équivalence.
\item À la demi-équivalence, la moitié de \ce{NH3} a été transformée en \ce{NH4+} : $[\ce{NH4+}] = [\ce{NH3}]$. Alors :
\[ pH = pK_a + \log\frac{[\ce{NH3}]}{[\ce{NH4+}]} = 9,2 + \log 1 = 9,2 \]
La solution contient un mélange équimolaire d'une base faible et de son acide conjugué : c'est une \textbf{solution tampon} (de pH 9,2).
\end{enumerate}
\textbf{Barème indicatif (sur 10) :} justification de la dilution 1 ; équation-bilan 1 ; explication de $pH_E < 7$ 1,5 ; calculs $C_b$ et $C_0$ + conclusion 2,5 ; $C_m$ 1 ; indicateur 1,5 ; demi-équivalence et tampon 1,5.

\textbf{Points de vigilance :} ne pas oublier de remonter à la concentration commerciale en multipliant par le facteur de dilution ; à la demi-équivalence, c'est le $pK_a$ du couple \ce{NH4+}/\ce{NH3} (9,2) qu'il faut utiliser, et non 7.
\end{corrige}

\begin{corrige}[Exercice 11 — Acide éthanoïque du vinaigre et préparation d'un tampon]
\begin{enumerate}
\item Équation-bilan :
\[ \ce{CH3COOH + HO- -> CH3COO- + H2O} \]
\item À la demi-équivalence, $[\ce{CH3COOH}] = [\ce{CH3COO-}]$, donc $pH = pK_a$. La courbe donne $pH_{1/2E} = 4,8$, d'où $\boxed{pK_a = 4,8}$ pour le couple \ce{CH3COOH}/\ce{CH3COO-}.
\[ K_a = 10^{-pK_a} = 10^{-4,8} \approx \SI{1,6e-5}{} \]
\item À l'équivalence, l'espèce prédominante est l'ion éthanoate \ce{CH3COO-}, base faible conjuguée de \ce{CH3COOH}. Elle réagit avec l'eau :
\[ \ce{CH3COO- + H2O <=> CH3COOH + HO-} \]
La présence d'ions \ce{HO-} rend la solution basique : $pH_E = 8,4 > 7$.
\item À l'équivalence : $C_a V_a = C_b V_{bE}$, donc pour le vinaigre dilué :
\[ C_a = \frac{0,10 \times 12,0}{10,0} = \SI{0,12}{mol.L^{-1}} \]
Le vinaigre commercial est dix fois plus concentré :
\[ C_0 = 10 \times 0,12 = \SI{1,2}{mol.L^{-1}} \]
\item Concentration massique du vinaigre commercial :
\[ C_m = C_0 \times M(\ce{CH3COOH}) = 1,2 \times 60,0 = \SI{72}{g.L^{-1}} \]
Avec $\rho = \SI{1,0}{g.mL^{-1}}$, \SI{100}{g} de vinaigre occupent \SI{100}{mL} = \SI{0,100}{L}. Masse d'acide dans \SI{100}{g} :
\[ m = 72 \times 0,100 = \SI{7,2}{g} \]
Le \textbf{degré d'acidité} du vinaigre est donc \textbf{7,2} (soit 7,2 \%), valeur typique d'un vinaigre de table.
\item Pour obtenir $pH = 4,8 = pK_a$, il faut $[\ce{CH3COO-}] = [\ce{CH3COOH}]$. Les deux solutions ayant la même concentration (\SI{0,10}{mol.L^{-1}}), il faut mélanger des volumes égaux :
\[ V_{\ce{CH3COOH}} = V_{\ce{CH3COONa}} = \frac{200}{2} = \SI{100}{mL} \]
\item Propriétés : le pH de cette solution varie très peu par ajout modéré d'acide ou de base, et il est pratiquement insensible à la dilution. Applications : le sang est tamponné à pH 7,4 par le couple \ce{H2CO3}/\ce{HCO3-} ; les collyres et sirops pharmaceutiques sont tamponnés pour ne pas irriter ; les jus (bissap) et boissons fermentées doivent leur stabilité gustative à des tampons naturels.
\end{enumerate}
\textbf{Barème indicatif (sur 10) :} équation-bilan 1 ; $pK_a$ justifié + $K_a$ 2 ; explication $pH_E > 7$ 1,5 ; $C_a$ et $C_0$ 2 ; degré d'acidité 1,5 ; volumes du tampon 1 ; propriétés et application 1.

\textbf{Points de vigilance :} le degré d'acidité se calcule sur le vinaigre \emph{commercial} (concentration $C_0$, pas $C_a$) ; ne pas oublier le facteur de dilution 10 ; à la demi-équivalence, le pH lu donne directement le $pK_a$ sans calcul supplémentaire.
\end{corrige}

\newpage

\coursec{Fiche bilan}

\coursec{Réactions de dosage et équivalence acido-basique}

\begin{savaistu}[Contexte camerounais : le contrôle qualité]
Au laboratoire de la \textbf{SONARA} (Limbe), on vérifie la teneur en acide sulfurique des électrolytes de batteries. À la brasserie, on dose l'acidité du moût pour garantir la fermentation. Au marché, le \emph{bili-bili} ou le vin de palme doivent avoir une acidité maîtrisée pour la conservation et le goût. Partout, le \textbf{dosage acido-basique} est l'outil de référence : rapide, peu coûteux et rigoureux.
\end{savaistu}

\begin{definition}[Dosage (ou titrage) acido-basique]
C'est une opération qui permet de \textbf{déterminer la concentration inconnue} ($C_X$) d'une solution (le \cle{dosé}) en la faisant réagir avec une solution de \textbf{concentration connue} (le \cle{doseur}, ou étalon), grâce à une réaction acido-basique.
\end{definition}

\begin{propriete}[Conditions de faisabilité d'un dosage]
Pour qu'un dosage soit exploitable, la réaction mise en jeu doit être :
\begin{enumerate}
  \item \textbf{Rapide} : l'équilibre est atteint quasi instantanément après chaque ajout.
  \item \textbf{Totale} (ou quasi totale) : la réaction va jusqu'à consommation complète du réactif limitant ($K \gg 1$, en pratique $K > 10^4$).
  \item \textbf{Unique} : une seule réaction se produit entre le dosé et le doseur (pas d'interférence).
\end{enumerate}
\end{propriete}

\courssub{Les quatre cas de figure et leurs équations-bilan}

On note \ce{AH} l'acide (fort ou faible) et \ce{B} la base (forte ou faible). Les équations-bilan s'écrivent avec une \textbf{flèche simple} $\ce{->}$ si la réaction est totale, $\ce{<=>}$ si elle est limitée.

\begin{center}
\begin{tabularx}{\linewidth}{|l|X|c|}
\hline
\rowcolor{popblueL} \textbf{Cas} & \textbf{Équation-bilan} & \textbf{Totalité} \\
\hline
Acide fort / Base forte & $\ce{H3O+ + HO- -> 2 H2O}$ & Totale ($K = 10^{14}$) \\
\hline
Acide faible / Base forte & $\ce{AH + HO- -> A- + H2O}$ & Totale ($K = 10^{pK_a(\ce{AH}) + 14} \gg 1$) \\
\hline
Acide fort / Base faible & $\ce{B + H3O+ -> BH+ + H2O}$ & Totale ($K = 10^{14 - pK_a(\ce{BH+})} \gg 1$) \\
\hline
Acide faible / Base faible & $\ce{AH + B <=> A- + BH+}$ & \textbf{Limitée} ($K = 10^{pK_a(\ce{BH+}) - pK_a(\ce{AH})}$) \\
\hline
\end{tabularx}
\end{center}

\begin{attention}[Piège fréquent : charges et atomes]
Vérifie \textbf{systématiquement} l'équilibrage des charges et des atomes.
\begin{itemize}
  \item $\ce{H3O+ + HO- -> 2 H2O}$ : charges $+1-1=0$ à gauche, $0$ à droite. OK.
  \item $\ce{AH + HO- -> A- + H2O}$ : charges $0-1=-1$ à gauche, $-1+0=-1$ à droite. OK.
  \item $\ce{B + H3O+ -> BH+ + H2O}$ : charges $0+1=+1$ à gauche, $+1+0=+1$ à droite. OK.
\end{itemize}
Ne confonds pas \ce{AH} (acide faible, neutre le plus souvent) et \ce{H3O+} (ion oxonium, chargé).
\end{attention}

\courssub{Notion d'équivalence et relation fondamentale}

\begin{definition}[Équivalence acido-basique]
L'\cle{équivalence} est l'état du système lorsque les réactifs ont été introduits dans les \textbf{proportions stœchiométriques exactes} définies par l'équation-bilan. À l'équivalence, \textbf{les deux réactifs sont entièrement consommés} (quantités de matière nulles pour les réactifs limitants).
\end{definition}

\begin{propriete}[Relation à l'équivalence (monoacide / monobase)]
Soit un acide \ce{AH} (concentration $C_A$, volume $V_A$) dosé par une base \ce{B} (concentration $C_B$, volume versé $V_B$). À l'équivalence ($V_B = V_{BE}$) :
\[
n_A(\text{apporté}) = n_B(\text{apporté}) \quad \Rightarrow \quad \boxed{C_A V_A = C_B V_{BE}}
\]
\emph{Unités :} $C$ en \SI{}{\mol\per\liter} (ou \SI{}{\mol\per\L}), $V$ en \SI{}{\liter} (ou \SI{}{\mL} si les deux volumes sont dans la même unité).
\end{propriete}

\begin{methode}[Écrire l'équation-bilan et la relation à l'équivalence]
\begin{enumerate}
  \item Identifier les espèces réactives réelles : \ce{H3O+} (acide fort) ou \ce{AH} (acide faible) ; \ce{HO-} (base forte) ou \ce{B} (base faible).
  \item Écrire l'équation-bilan en respectant la stœchiométrie (coefficients 1 pour les monoacides/monobases).
  \item Déduire la relation $n_{\text{acide}} = n_{\text{base}}$ à l'équivalence, puis $C_A V_A = C_B V_{BE}$ (en adaptant les coefficients si polyacide/polybase).
  \item Vérifier l'homogénéité des unités avant tout calcul.
\end{enumerate}
\end{methode}

\begin{exoresolu}[Dosage d'une solution commerciale de soude]
\textbf{Énoncé :} On dose \SI{20,0}{\mL} d'une solution de soude (hydroxyde de sodium, base forte) par une solution d'acide chlorhydrique (acide fort) de concentration $C_A = \SI{0,100}{\mol\per\L}$. L'équivalence est atteinte pour $V_{BE} = \SI{18,5}{\mL}$ d'acide versé. Calculer la concentration $C_B$ de la solution de soude.

\textbf{Solution détaillée :}
\begin{enumerate}
  \item \textbf{Équation-bilan :} $\ce{H3O+ + HO- -> 2 H2O}$ (acide fort / base forte).
  \item \textbf{Relation à l'équivalence :} $n_{\ce{H3O+}} = n_{\ce{HO-}}$ $\Rightarrow$ $C_A V_{AE} = C_B V_{BE}$.
    \emph{Attention :} ici l'acide est le doseur (versé à la burette), la base est le dosé (dans le bécher). Le volume à l'équivalence est $V_{AE} = \SI{18,5}{\mL}$ pour l'acide. La relation s'écrit $C_A V_{AE} = C_B V_B$.
  \item \textbf{Calcul :}
  \[
  C_B = \frac{C_A V_{AE}}{V_B} = \frac{\SI{0,100}{\mol\per\L} \times \SI{18,5}{\mL}}{\SI{20,0}{\mL}} = \SI{0,0925}{\mol\per\L}
  \]
  \item \textbf{Résultat :} La solution de soude a une concentration de $\boldsymbol{\SI{0,0925}{\mol\per\L}}$.
\end{enumerate}
\end{exoresolu}

\begin{exercice}[Entraînement : dosage vinaigre]
On dose \SI{10,0}{\mL} de vinaigre (acide acétique \ce{CH3COOH}, acide faible, $pK_a = 4,8$) par une solution de soude \SI{0,100}{\mol\per\L}. L'équivalence est atteinte pour $V_{BE} = \SI{12,0}{\mL}$.
\begin{enumerate}
  \item Écrire l'équation-bilan de la réaction.
  \item Calculer la concentration molaire de l'acide acétique dans le vinaigre.
  \item En déduire la concentration massique en \SI{}{\gram\per\L} ($M_{\ce{CH3COOH}} = \SI{60,0}{\gram\per\mol}$).
\end{enumerate}
\end{exercice}

\begin{corrige}[Exercice : dosage vinaigre]
\begin{enumerate}
  \item $\ce{CH3COOH + HO- -> CH3COO- + H2O}$.
  \item $C_A V_A = C_B V_{BE} \Rightarrow C_A = \frac{\SI{0,100}{\mol\per\L} \times \SI{12,0}{\mL}}{\SI{10,0}{\mL}} = \SI{0,120}{\mol\per\L}$.
  \item $\rho_m = C_A \times M = \SI{0,120}{\mol\per\L} \times \SI{60,0}{\gram\per\mol} = \SI{7,20}{\gram\per\L}$.
\end{enumerate}
\end{corrige}

\coursec{Dosage pH-métrique d'un acide fort par une base forte}

\begin{experience}[Dispositif et protocole pH-métrique]
\textbf{Matériel :} Bécher \SI{100}{\mL}, burette \SI{50}{\mL} (base forte \ce{HO-}), pH-mètre étalonné (deux points : pH 4,01 et 7,00 ou 9,21), barre magnétique + agitateur, solution acide fort \ce{H3O+} (inconnue $C_A$, volume $V_A$).
\textbf{Protocole :}
\begin{enumerate}
  \item Verser $V_A$ de l'acide dans le bécher, immerger l'électrode, agiter doucement.
  \item Noter le pH initial ($V_B = 0$).
  \item Ajouter la base par fractions (\SI{1}{\mL} puis \SI{0,1}{\mL} près de l'équivalence), noter $V_B$ et pH après stabilisation.
  \item Continuer bien après l'équivalence (excès de base).
  \item Tracer la courbe $\mathrm{pH} = f(V_B)$.
\end{enumerate}
\end{experience}

\begin{popfigure}
\begin{center}
\begin{tikzpicture}[scale=0.9]
\begin{axis}[
  width=\linewidth, height=8cm,
  axis lines=middle,
  xlabel={$V_B$ (\si{\mL})}, ylabel={pH},
  xmin=0, xmax=25, ymin=0, ymax=14,
  xtick={0,5,10,12,15,20,25}, ytick={0,1,2,3,4,5,6,7,8,9,10,11,12,13,14},
  grid=major, grid style={popline},
  clip=false
]
% Courbe théorique pH = f(V) pour dosage fort/fort
% Paramètres : C_A = 0.1 M, V_A = 10 mL, C_B = 0.1 M -> V_BE = 10 mL
% Avant E : pH = -log( C_A*V_A - C_B*V_B / (V_A+V_B) ) ... approximation simplifiée pour tracé
% On trace une sigmoïde centrée sur 10 mL, pH 7
\addplot[popcurve, domain=0:9.9, samples=100] 
  {-log10( max(1e-14, (1*10 - 1*x)/(10+x) ) )}; % Approx [H3O+]
\addplot[popcurve, domain=10.1:25, samples=100] 
  {14 + log10( max(1e-14, (1*x - 1*10)/(10+x) ) )}; % Approx pOH -> pH
% Point équivalence
\addplot[only marks, mark=*, mark size=2.5, popink] coordinates {(10,7)};
\node[popink, anchor=south west] at (axis cs:10,7) {$E(10; 7)$};
% Tangentes
\addplot[poporange, dashed, domain=0:25] {0.5*x + 2}; % Tangente avant
\addplot[poporange, dashed, domain=0:25] {0.5*x - 5}; % Tangente après
% Droite équidistante (verticale ici car tangentes parallèles symétriques)
\addplot[popgreen, thick, dashed, domain=0:14] ({10}, {x});
\node[popgreen, anchor=south east] at (axis cs:10,13) {Droite équidistante};
\end{axis}
\end{tikzpicture}
\end{center}
\legende{Courbe pH-métrique type : dosage acide fort / base forte ($C_A=C_B=\SI{0,1}{\mol\per\L}$, $V_A=\SI{10}{\mL}$). Saut de pH vertical en $E$ : pH$_E = 7$ à \SI{25}{\celsius}.}
\end{popfigure}

\courssub{Étude analytique de la courbe : trois zones}

On note $n_A^0 = C_A V_A$ (quantité initiale d'acide), $n_B = C_B V_B$ (quantité de base versée). Volume total $V_T = V_A + V_B$.

\begin{enumerate}
  \item \textbf{Avant l'équivalence} ($V_B < V_{BE}$, excès d'acide) :
    $n_{\ce{H3O+}} = n_A^0 - n_B = C_A V_A - C_B V_B$.
    \[
    [\ce{H3O+}] = \frac{C_A V_A - C_B V_B}{V_A + V_B} \quad \Rightarrow \quad \mathrm{pH} = -\log\left(\frac{C_A V_A - C_B V_B}{V_A + V_B}\right)
    \]
  \item \textbf{À l'équivalence} ($V_B = V_{BE}$) :
    Réactifs totalement consommés. Solution neutre (eau + sel neutre \ce{NaCl}, \ce{KNO3}...).
    \[
    \mathrm{pH}_E = 7 \quad (\text{à } \SI{25}{\celsius})
    \]
  \item \textbf{Après l'équivalence} ($V_B > V_{BE}$, excès de base) :
    $n_{\ce{HO-}} = n_B - n_A^0 = C_B V_B - C_A V_A$.
    \[
    [\ce{HO-}] = \frac{C_B V_B - C_A V_A}{V_A + V_B} \quad \Rightarrow \quad \mathrm{pOH} = -\log[\ce{HO-}] \quad \Rightarrow \quad \mathrm{pH} = 14 - \mathrm{pOH}
    \]
\end{enumerate}

\begin{propriete}[Caractéristiques de la courbe fort/fort]
\begin{itemize}
  \item Forme sigmoïde avec un \textbf{saut de pH très grand} (environ 3 à 11) centré sur $V_{BE}$.
  \item pH à l'équivalence \textbf{strictement égal à 7} (à \SI{25}{\celsius}).
  \item La courbe est symétrique par rapport au point $E$ si $C_A = C_B$.
\end{itemize}
\end{propriete}

\courssub{Détermination graphique de l'équivalence : méthode des tangentes}

\begin{methode}[Méthode des tangentes (classique au Bac)]
\begin{enumerate}
  \item Tracer la courbe $\mathrm{pH} = f(V_B)$ point par point.
  \item Tracer \textbf{deux tangentes} : une sur la branche montante avant le saut, une sur la branche montante après le saut. Elles doivent être \textbf{parallèles}.
  \item Tracer la \textbf{droite équidistante} de ces deux tangentes (médiatrice de tout segment perpendiculaire les reliant). Cette droite est verticale si les tangentes sont parallèles et symétriques.
  \item L'intersection de cette droite équidistante avec la courbe donne le point \textbf{E} : on lit $V_{BE}$ (abscisse) et $\mathrm{pH}_E$ (ordonnée).
\end{enumerate}
\end{methode}

\begin{methode}[Méthode de la dérivée (plus précise)]
\begin{enumerate}
  \item Calculer numériquement la dérivée $\frac{\Delta \mathrm{pH}}{\Delta V_B}$ (ou $\frac{d\mathrm{pH}}{dV_B}$) pour chaque intervalle de mesure.
  \item Tracer la courbe $\frac{d\mathrm{pH}}{dV_B} = f(V_B)$.
  \item Le \textbf{maximum} de cette courbe correspond à l'équivalence. L'abscisse du maximum donne $V_{BE}$.
\end{enumerate}
\end{methode}

\begin{attention}[Erreurs fréquentes sur les tangentes]
\begin{itemize}
  \item Ne pas tracer les tangentes sur la partie \textbf{verticale} du saut (pente infinie).
  \item Les tangentes doivent être tracées \textbf{loin du saut}, sur les parties courbes bien définies (zones tamponnées par l'excès de réactif).
  \item Si les concentrations sont très différentes ($C_A \ll C_B$ ou inverse), le saut est moins vertical ; la méthode de la dérivée est alors plus fiable.
\end{itemize}
\end{attention}

\begin{exoresolu}[Exploitation d'une courbe pH-métrique fort/fort]
\textbf{Énoncé :} On dose \SI{20,0}{\mL} d'une solution d'acide chlorhydrique par une solution de soude \SI{0,100}{\mol\per\L}. La courbe $\mathrm{pH}=f(V_B)$ donne $V_{BE} = \SI{15,0}{\mL}$ et $\mathrm{pH}_E = 7,0$.
\begin{enumerate}
  \item Calculer $C_A$.
  \item Calculer le pH pour $V_B = \SI{5,0}{\mL}$ et $V_B = \SI{25,0}{\mL}$.
\end{enumerate}
\textbf{Solution :}
\begin{enumerate}
  \item $C_A = \frac{C_B V_{BE}}{V_A} = \frac{0,100 \times 15,0}{20,0} = \SI{0,0750}{\mol\per\L}$.
  \item $V_B = \SI{5,0}{\mL} < V_{BE}$ : excès acide. $n_{\ce{H3O+}} = (0,0750 \times 20,0 - 0,100 \times 5,0) \times 10^{-3} = \SI{1,00e-3}{\mol}$. $V_T = \SI{25,0}{\mL}$. $[\ce{H3O+}] = \frac{1,00 \times 10^{-3}}{25,0 \times 10^{-3}} = \SI{0,0400}{\mol\per\L}$. $\mathrm{pH} = -\log(0,0400) = \mathbf{1,40}$.
  \item $V_B = \SI{25,0}{\mL} > V_{BE}$ : excès base. $n_{\ce{HO-}} = (0,100 \times 25,0 - 0,0750 \times 20,0) \times 10^{-3} = \SI{1,00e-3}{\mol}$. $V_T = \SI{45,0}{\mL}$. $[\ce{HO-}] = \frac{1,00 \times 10^{-3}}{45,0 \times 10^{-3}} = \SI{0,0222}{\mol\per\L}$. $\mathrm{pOH} = 1,65$. $\mathrm{pH} = 14 - 1,65 = \mathbf{12,35}$.
\end{enumerate}
\end{exoresolu}

\coursec{Dosage d'un acide faible par une base forte (et réciproques)}

\courssub{Dosage d'un acide faible \ce{AH} par une base forte \ce{HO-}}

\begin{experience}[Courbe pH-métrique acide faible / base forte]
Même dispositif. Dosé : acide acétique \ce{CH3COOH} \SI{0,1}{\mol\per\L} (\SI{20}{\mL}). Doseur : \ce{NaOH} \SI{0,1}{\mol\per\L}. $pK_a = 4,8$.
Observations : pH initial plus élevé (acide faible), zone plate (tampon), demi-équivalence, saut de pH réduit (8-10), équivalence basique.
\end{experience}

\begin{popfigure}
\begin{center}
\begin{tikzpicture}[scale=0.9]
\begin{axis}[
  width=\linewidth, height=8cm,
  axis lines=middle,
  xlabel={$V_B$ (\si{\mL})}, ylabel={pH},
  xmin=0, xmax=25, ymin=0, ymax=14,
  xtick={0,5,10,12,15,20,25}, ytick={0,1,2,3,4,5,6,7,8,9,10,11,12,13,14},
  grid=major, grid style={popline},
  clip=false
]
% Courbe théorique acide faible (pKa=4.8) / base forte
% Paramètres : C_A = 0.1 M, V_A = 10 mL, C_B = 0.1 M -> V_BE = 10 mL
% Avant E : Henderson-Hasselbalch pH = pKa + log(n_A-/n_AH)
% n_AH = 1 - 0.1*x ; n_A- = 0.1*x (mmol) pour x en mL
\addplot[popcurve, domain=0.1:9.9, samples=200] 
  {4.8 + log10( (0.1*x) / (1 - 0.1*x) )};
% Point demi-équivalence
\addplot[only marks, mark=*, mark size=2.5, poporange] coordinates {(5, 4.8)};
\node[poporange, anchor=south west] at (axis cs:5, 4.8) {$1/2E$};
% Équivalence : hydrolyse A- -> pH > 7. Calcul approx : [A-] = 0.05 M. Kb = 10^-9.2. [OH-] = sqrt(Kb*C) = sqrt(10^-9.2 * 0.05) ~ 5.6e-6. pOH=5.25 -> pH=8.75
\addplot[only marks, mark=*, mark size=2.5, popink] coordinates {(10, 8.75)};
\node[popink, anchor=south west] at (axis cs:10, 8.75) {$E(10; 8,75)$};
% Après E : excès HO-
\addplot[popcurve, domain=10.1:25, samples=100] 
  {14 + log10( max(1e-14, (1*x - 1*10)/(10+x) ) )};
% Tangentes
\addplot[poporange, dashed, domain=0:25] {0.4*x + 2.8};
\addplot[poporange, dashed, domain=0:25] {0.4*x - 1.2};
\addplot[popgreen, thick, dashed, domain=0:14] ({10}, {x});
\end{axis}
\end{tikzpicture}
\end{center}
\legende{Courbe pH-métrique : acide faible ($pK_a=4,8$) / base forte. Zone tampon (plate) avant $E$, demi-équivalence $1/2E$ (pH $= pK_a$), saut de pH réduit ($\approx$ 7-10), pH$_E > 7$.}
\end{popfigure}

\begin{propriete}[Équation-bilan et tableau d'avancement (acide faible / base forte)]
\[
\ce{AH + HO- -> A- + H2O} \quad (K = 1/K_b(\ce{A-}) = 10^{pK_a+14} \gg 1)
\]
\begin{center}
\begin{tabular}{|c|c|c|c|c|}
\hline
\rowcolor{popblueL} & \ce{AH} & \ce{HO-} & \ce{A-} & \ce{H2O} \\
\hline
État initial & $n_A^0$ & $n_B$ & 0 & -- \\
\hline
Avancement & $-x$ & $-x$ & $+x$ & -- \\
\hline
État final & $n_A^0-x$ & $n_B-x$ & $x$ & -- \\
\hline
\end{tabular}
\end{center}
À l'équivalence ($x = n_A^0 = n_B$) : $n_{\ce{AH}} = 0$, $n_{\ce{HO-}} = 0$, $n_{\ce{A-}} = n_A^0$. Le milieu contient \textbf{uniquement la base conjuguée \ce{A-}}.
\end{propriete}

\begin{propriete}[pH à la demi-équivalence]
À la demi-équivalence ($V_B = V_{BE}/2$), on a versé la moitié de la base nécessaire. D'après le tableau : $n_{\ce{AH}} = n_{\ce{A-}} = n_A^0/2$.
D'après l'équation de Henderson-Hasselbalch (valable dans la zone tampon) :
\[
\mathrm{pH} = pK_a + \log\frac{[\ce{A-}]}{[\ce{AH}]} = pK_a + \log 1 = \boxed{pK_a}
\]
C'est une méthode expérimentale fondamentale pour \textbf{déterminer le $pK_a$ d'un acide faible}.
\end{propriete}

\begin{propriete}[pH à l'équivalence (acide faible / base forte)]
Le milieu ne contient que \ce{A-} (base conjuguée de \ce{AH}), qui hydrolyse :
\[
\ce{A- + H2O <=> AH + HO-} \quad K_b = \frac{K_e}{K_a} = 10^{pK_a - 14}
\]
Le milieu est \textbf{basique} : $\mathrm{pH}_E > 7$ (à \SI{25}{\celsius}).
Calcul approché (si $K_b$ petit) : $[\ce{HO-}] \approx \sqrt{K_b \times C_{\ce{A-}}}$ avec $C_{\ce{A-}} = \frac{n_A^0}{V_A + V_{BE}}$.
\end{propriete}

\courssub{Dosage d'une base faible \ce{B} par un acide fort \ce{H3O+}}

C'est le cas symétrique (ex: dosage ammoniaque \ce{NH3} par \ce{HCl}).
\begin{itemize}
  \item \textbf{Équation-bilan :} $\ce{B + H3O+ -> BH+ + H2O}$ (totale).
  \item \textbf{Demi-équivalence :} $V_A = V_{AE}/2$, $n_{\ce{B}} = n_{\ce{BH+}}$, $\mathrm{pH} = pK_a(\ce{BH+}) = 14 - pK_b(\ce{B})$.
  \item \textbf{Équivalence :} Milieu ne contient que \ce{BH+} (acide conjugué de \ce{B}), qui hydrolyse : $\ce{BH+ + H2O <=> B + H3O+}$. Milieu \textbf{acide} : $\mathrm{pH}_E < 7$.
  \item \textbf{Saut de pH :} Centré sur pH acide (ex: 4-6 pour \ce{NH3}/\ce{NH4+}, $pK_a = 9,2$).
\end{itemize}

\courssub{Dosage acide faible / base faible}

\begin{propriete}[Réaction limitée]
$\ce{AH + B <=> A- + BH+}$ avec $K = \frac{K_a(\ce{AH})}{K_a(\ce{BH+})} = 10^{pK_a(\ce{BH+}) - pK_a(\ce{AH})}$.
\begin{itemize}
  \item Si $\Delta pK_a > 4$ ($K > 10^4$) : réaction quasi totale $\rightarrow$ dosage possible (saut de pH observable).
  \item Si $\Delta pK_a$ faible : réaction peu déplacée $\rightarrow$ pas de saut de pH net $\rightarrow$ \textbf{dosage impossible} directement.
\end{itemize}
\end{propriete}

\begin{exoresolu}[Dosage acide acétique / soude : détermination de $pK_a$]
\textbf{Énoncé :} Dosage de \SI{20,0}{\mL} d'acide acétique \SI{0,100}{\mol\per\L} par \ce{NaOH} \SI{0,100}{\mol\per\L}. On relève : $V_{BE} = \SI{20,0}{\mL}$, $\mathrm{pH}_E = 8,72$. À $V_B = \SI{10,0}{\mL}$, $\mathrm{pH} = 4,75$.
\begin{enumerate}
  \item Vérifier que $V_B = \SI{10,0}{\mL}$ correspond à la demi-équivalence.
  \item En déduire $pK_a$ de l'acide acétique.
  \item Justifier que $\mathrm{pH}_E > 7$.
\end{enumerate}
\textbf{Solution :}
\begin{enumerate}
  \item $V_{BE}/2 = 20,0/2 = \SI{10,0}{\mL} = V_B$. C'est bien la demi-équivalence.
  \item À la demi-équivalence, $\mathrm{pH} = pK_a = \mathbf{4,75}$.
  \item À l'équivalence, seul \ce{CH3COO-} reste. $C_{\ce{CH3COO-}} = \frac{0,100 \times 20,0}{40,0} = \SI{0,0500}{\mol\per\L}$. $K_b = 10^{-14}/10^{-4,75} = 10^{-9,25}$. $[\ce{HO-}] = \sqrt{10^{-9,25} \times 0,0500} \approx \SI{1,88e-6}{\mol\per\L}$. $\mathrm{pOH} = 5,73$, $\mathrm{pH} = 8,27$ (proche de 8,72, différence due aux approximations d'activité). Le milieu est basique.
\end{enumerate}
\end{exoresolu}

\begin{exercice}[Dosage ammoniaque]
On dose \SI{25,0}{\mL} d'une solution d'ammoniaque \ce{NH3} (base faible, $pK_b = 4,75$) par de l'acide chlorhydrique \SI{0,100}{\mol\per\L}. L'équivalence est atteinte pour $V_{AE} = \SI{22,0}{\mL}$.
\begin{enumerate}
  \item Écrire l'équation-bilan. Calculer $C_B$ (concentration de \ce{NH3}).
  \item Quel est le pH à la demi-équivalence ? ($pK_a(\ce{NH4+}) = 9,25$).
  \item Prédire le pH à l'équivalence ($<7$, $=7$ ou $>7$ ?) et justifier.
\end{enumerate}
\end{exercice}

\begin{corrige}[Exercice : dosage ammoniaque]
\begin{enumerate}
  \item $\ce{NH3 + H3O+ -> NH4+ + H2O}$. $C_B = \frac{C_A V_{AE}}{V_B} = \frac{0,100 \times 22,0}{25,0} = \SI{0,0880}{\mol\per\L}$.
  \item Demi-équivalence : $V_A = \SI{11,0}{mL}$. $n_{\ce{NH3}} = n_{\ce{NH4+}}$. $\mathrm{pH} = pK_\mathrm{a}(\ce{NH4+}) = \mathbf{9,25}$.
  \item À l'équivalence, seul \ce{NH4+} (acide faible) reste. Hydrolyse : $\ce{NH4+ + H2O <=> NH3 + H3O+}$. Milieu acide : $\mathbf{\mathrm{pH}_E < 7}$.
\end{enumerate}
\end{corrige}

\coursec{Dosage colorimétrique et choix de l'indicateur coloré}

\begin{definition}[Indicateur coloré]
C'est un \textbf{couple acide/base conjugués} (\ce{HIn}/\ce{In-}) dont les deux formes ont des \textbf{couleurs distinctes} dans le domaine visible. L'indicateur est lui-même un acide faible (ou une base faible).
\[
\ce{HIn <=> H+ + In-} \quad \text{Couleur 1} \quad \rightleftharpoons \quad \text{Couleur 2}
\]
\end{definition}

\begin{propriete}[Zone de virage]
L'indicateur change de couleur sur un intervalle de pH centré sur son $pK_a$ :
\[
\text{Zone de virage} = [pK_a - 1 ; pK_a + 1]
\]
\begin{itemize}
  \item pH $< pK_a - 1$ : forme acide \ce{HIn} prédominante (couleur 1).
  \item pH $> pK_a + 1$ : forme basique \ce{In-} prédominante (couleur 2).
  \item Entre les deux : mélange des deux couleurs (couleur intermédiaire).
\end{itemize}
\end{propriete}

\begin{center}
\begin{tabularx}{\linewidth}{|l|c|c|c|X|}
\hline
\rowcolor{popblueL} \textbf{Indicateur} & \textbf{$pK_a$} & \textbf{Zone de virage} & \textbf{Couleur acide} & \textbf{Couleur basique} \\
\hline
Hélianthine & 1,5 & 0,5 -- 2,5 & Rouge & Jaune \\
\hline
Bromothymol blue (B.T.B.) & 7,0 & 6,0 -- 8,0 & Jaune & Bleu \\
\hline
Phénolphtaléine (P.Ph.) & 9,3 & 8,3 -- 10,0 & Incolore & Rose/Fuchsia \\
\hline
\end{tabularx}
\end{center}

\begin{methode}[Choix judicieux de l'indicateur coloré]
\begin{enumerate}
  \item Déterminer (ou estimer) le \textbf{pH à l'équivalence} $\mathrm{pH}_E$ du dosage prévu.
  \item Choisir un indicateur dont la \textbf{zone de virage contient $\mathrm{pH}_E$}.
  \item Vérifier que le changement de couleur est net et visible (contraste fort).
\end{enumerate}
\end{methode}

\begin{center}
\begin{tabularx}{\linewidth}{|l|X|c|}
\hline
\rowcolor{popblueL} \textbf{Type de dosage} & \textbf{pH$_E$} & \textbf{Indicateur recommandé} \\
\hline
Acide fort / Base forte & 7,0 & \textbf{B.T.B.} (zone 6-8) ou Phénolphtaléine (zone 8,3-10) \\
\hline
Acide faible / Base forte & $> 7$ (ex: 8-10) & \textbf{Phénolphtaléine} (zone 8,3-10) \\
\hline
Acide fort / Base faible & $< 7$ (ex: 4-6) & \textbf{Hélianthine} (zone 0,5-2,5) ou \textbf{B.T.B.} (zone 6-8) si pH$_E \approx 6$ \\
\hline
\end{tabularx}
\end{center}

\begin{attention}[Pièges du dosage colorimétrique]
\begin{itemize}
  \item \textbf{Indicateur mal choisi :} Si la zone de virage ne contient pas pH$_E$, le virage se produit \textbf{avant} ou \textbf{après} l'équivalence $\rightarrow$ \textbf{erreur systématique} sur $V_{BE}$ (surestimation ou sous-estimation de la concentration).
  \item \textbf{Excès d'indicateur :} L'indicateur consomme lui-même du doseur (c'est un acide/base). Mettre \textbf{2 à 3 gouttes maxi}.
  \item \textbf{Couleur du milieu :} Si le dosé est coloré (ex: jus de bissap, thé), le virage est masqué $\rightarrow$ privilégier le \textbf{dosage pH-métrique}.
\end{itemize}
\end{attention}

\begin{savaistu}[Le bissap et les anthocyanes]
Le \emph{bissap} (fleurs d'hibiscus) contient des \textbf{anthocyanes}, pigments naturels qui changent de couleur avec le pH (rouge acide, bleu/violet basique). C'est un indicateur coloré naturel ! Au village, on teste l'acidité du vin de palme en y trempant une feuille d'hibiscus séchée. En chimie analytique moderne, on utilise des indicateurs de synthèse (phénolphtaléine, BTB) plus reproductibles.
\end{savaistu}

\coursec{Solutions tampons : définition, propriétés et préparation}

\begin{definition}[Solution tampon (ou solution tamponnée)]
C'est une solution dont le \textbf{pH varie très peu} lors de l'ajout d'une \textbf{faible quantité} d'acide fort, de base forte, ou lors d'une \textbf{dilution}. Elle est constituée d'un \textbf{couple acide faible / base conjuguée} (\ce{AH}/\ce{A-}) (ou base faible / acide conjugué \ce{B}/\ce{BH+}) en \textbf{concentrations comparables} (généralement rapport entre 0,1 et 10).
\end{definition}

\begin{propriete}[Équation de Henderson-Hasselbalch]
Pour un tampon \ce{AH}/\ce{A-} :
\[
\boxed{\mathrm{pH} = pK_a(\ce{AH}) + \log\frac{[\ce{A-}]}{[\ce{AH}]}}
\]
\begin{itemize}
  \item Valable si les concentrations sont les concentrations \textbf{d'équilibre} (souvent approximées par les concentrations formelles/apportées car dissociation faible).
  \item Si $[\ce{A-}] = [\ce{AH}]$, alors $\mathrm{pH} = pK_a$.
  \item Le pH ne dépend \textbf{pas de la dilution} (le rapport des concentrations reste constant), mais la \textbf{capacité tampon} diminue.
\end{itemize}
\end{propriete}

\begin{propriete}[Mécanisme d'action (effet tampon)]
\begin{itemize}
  \item \textbf{Ajout d'acide fort (\ce{H3O+}) :} $\ce{H3O+ + A- -> AH + H2O}$. La base conjuguée \ce{A-} neutralise l'acide ajouté. $[\ce{A-}]$ baisse, $[\ce{AH}]$ monte $\rightarrow$ pH baisse faiblement.
  \item \textbf{Ajout de base forte (\ce{HO-}) :} $\ce{HO- + AH -> A- + H2O}$. L'acide faible \ce{AH} neutralise la base ajoutée. $[\ce{AH}]$ baisse, $[\ce{A-}]$ monte $\rightarrow$ pH monte faiblement.
\end{itemize}
\end{propriete}

\begin{definition}[Capacité tampon]
C'est la quantité de matière d'acide fort ou de base forte qu'il faut ajouter à \SI{1}{\L} de solution tampon pour faire varier son pH d'une unité. Elle est \textbf{maximale} pour $\mathrm{pH} = pK_a$ (c'est-à-dire $[\ce{AH}] = [\ce{A-}]$) et augmente avec la \textbf{concentration totale} $C_{\text{tot}} = [\ce{AH}] + [\ce{A-}]$.
\end{definition}

\courssub{Préparation d'une solution tampon}

\begin{methode}[Préparer un tampon de pH = pKa (cas le plus courant au Bac)]
\textbf{Objectif :} Obtenir $[\ce{AH}] = [\ce{A-}]$.
\begin{enumerate}
  \item \textbf{Méthode par mélange (la plus simple) :} Mélanger \textbf{volumes égaux} ($V$) de deux solutions mères de \textbf{même concentration} $C_0$ : l'acide faible \ce{AH} et sa base conjuguée \ce{A-} (souvent apportée par un sel, ex: \ce{CH3COONa}).
    \begin{itemize}
      \item Concentrations finales : $[\ce{AH}] = [\ce{A-}] = C_0/2$.
      \item pH final = $pK_a$.
    \end{itemize}
  \item \textbf{Méthode par demi-neutralisation (dosage) :} Dosage d'un volume $V_A$ d'acide faible \ce{AH} ($C_A$) par une base forte \ce{HO-} ($C_B = C_A$) \textbf{jusqu'à la demi-équivalence} ($V_B = V_{BE}/2$).
    \begin{itemize}
      \item À $1/2E$ : $n_{\ce{AH}} = n_{\ce{A-}}$.
      \item Le mélange réactionnel \textbf{est} la solution tampon de pH $= pK_a$.
    \end{itemize}
\end{enumerate}
\end{methode}

\begin{methode}[Préparer un tampon de pH imposé $\neq pK_a$ (Équation de Henderson-Hasselbalch)]
\begin{enumerate}
  \item Choisir un couple \ce{AH}/\ce{A-} dont le $pK_a$ est \textbf{proche du pH visé} (idéalement $|pH - pK_a| < 1$).
  \item Calculer le rapport $R = \frac{[\ce{A-}]}{[\ce{AH}]} = 10^{\mathrm{pH} - pK_a}$.
  \item Choisir une concentration totale $C_{\text{tot}} = [\ce{AH}] + [\ce{A-}]$ (ex: \SI{0,1}{\mol\per\L}).
  \item Calculer les concentrations cibles : $[\ce{AH}] = \frac{C_{\text{tot}}}{1+R}$, $[\ce{A-}] = \frac{R \cdot C_{\text{tot}}}{1+R}$.
  \item Mélanger les volumes appropriés de solutions mères d'acide et de sel (ou faire un dosage partiel).
\end{enumerate}
\end{methode}

\begin{exoresolu}[Préparation tampon acétate pH = 5,0]
\textbf{Énoncé :} Préparer \SI{100}{\mL} de solution tampon \ce{CH3COOH}/\ce{CH3COO-} de pH = 5,00 et de concentration totale \SI{0,100}{\mol\per\L}. $pK_a = 4,75$. On dispose de solutions mères : acide acétique \SI{0,200}{\mol\per\L} et acétate de sodium \SI{0,200}{\mol\per\L}.
\textbf{Solution :}
\begin{enumerate}
  \item $R = 10^{5,00 - 4,75} = 10^{0,25} \approx 1,78$.
  \item $[\ce{AH}] = \frac{0,100}{1+1,78} = \SI{0,0360}{\mol\per\L}$ ; $[\ce{A-}] = \frac{1,78 \times 0,100}{2,78} = \SI{0,0640}{\mol\per\L}$.
  \item Quantités pour \SI{100}{\mL} : $n_{\ce{AH}} = \SI{3,60e-3}{\mol}$, $n_{\ce{A-}} = \SI{6,40e-3}{\mol}$.
  \item Volumes solutions mères (\SI{0,200}{\mol\per\L}) : $V_{\ce{AH}} = \frac{3,60 \times 10^{-3}}{0,200} = \SI{18,0}{\mL}$ ; $V_{\ce{A-}} = \frac{6,40 \times 10^{-3}}{0,200} = \SI{32,0}{\mL}$.
  \item \textbf{Protocole :} Mesurer \SI{18,0}{\mL} acide acétique \SI{0,200}{\mol\per\L} + \SI{32,0}{\mL} acétate de sodium \SI{0,200}{\mol\per\L} dans une fiole jaugée \SI{100}{\mL}, compléter à l'eau distillée, homogénéiser. Vérifier le pH au pH-mètre.
\end{enumerate}
\end{exoresolu}

\begin{experience}[Test de l'effet tampon]
\textbf{Protocole :}
\begin{enumerate}
  \item Préparer 2 béchers : (1) \SI{50}{\mL} eau distillée, (2) \SI{50}{\mL} tampon acétate pH 4,75 ($C_{\text{tot}} = \SI{0,1}{\mol\per\L}$).
  \item Mesurer pH initial.
  \item Ajouter \SI{1}{\mL} \ce{HCl} \SI{0,1}{\mol\per\L} dans chaque bécher. Mesurer nouveau pH.
  \item Ajouter \SI{1}{\mL} \ce{NaOH} \SI{0,1}{\mol\per\L} dans chaque bécher (neufs). Mesurer nouveau pH.
\end{enumerate}
\textbf{Résultats attendus :}
\begin{itemize}
  \item Eau : pH passe de 7 $\rightarrow$ $\approx$ 3 (acide) ou $\approx$ 11 (base). Variation $\approx 4$ à $8$ unités !
  \item Tampon : pH passe de 4,75 $\rightarrow$ $\approx$ 4,6 (acide) ou $\approx$ 4,9 (base). Variation $\approx 0,1$ à $0,2$ unité.
\end{itemize}
\emph{Conclusion :} Le tampon résiste efficacement.
\end{experience}

\begin{savaistu}[Tampons dans la vie et l'industrie au Cameroun]
\begin{itemize}
  \item \textbf{Sang humain :} Tampon bicarbonate \ce{H2CO3}/\ce{HCO3-} ($pK_a \approx 6,1$, pH 7,4) + hémoglobine. Maintient le pH sanguin entre 7,35 et 7,45. Vital pour l'activité enzymatique.
  \item \textbf{Industrie agro-alimentaire (brasseries, jus) :} Tampons citrate/acide citrique pour stabiliser le pH des boissons (bière, jus de bissap pasteurisé) et éviter la précipitation ou le goût métallique.
  \item \textbf{Cosmétiques (savonneries artisanales, huiles) :} Tampons phosphate ou citrate pour fixer le pH des crèmes, shampoings (pH 5,5 « pH physiologique ») et savons (pH $\approx$ 9-10 tamponné pour douceur).
  \item \textbf{CIMENCAM / SONARA :} Traitement des eaux usées : ajustement pH par tampons avant rejet ou traitement biologique (bactéries sensibles au pH).
  \item \textbf{Pharmacopée :} Préparation de solutions injectables, collyres (tampons phosphate, borate) pour éviter l'irritation (pH larmes $\approx$ 7,4).
\end{itemize}
\end{savaistu}

\begin{exercice}[Tampon phosphate]
On souhaite préparer \SI{500}{\mL} d'une solution tampon phosphate de pH = 7,20. $pK_{a2}(\ce{H2PO4-}) = 7,20$. On dispose de \ce{NaH2PO4} (acide) et \ce{Na2HPO4} (base) solides. $M(\ce{NaH2PO4}) = \SI{120}{\gram\per\mol}$, $M(\ce{Na2HPO4}) = \SI{142}{\gram\per\mol}$. Concentration totale souhaitée : \SI{0,050}{\mol\per\L}.
Calculer les masses de sel à peser.
\end{exercice}

\begin{corrige}[Exercice : Tampon phosphate]
pH = pKa = 7,20 $\Rightarrow$ $[\ce{A-}] = [\ce{AH}]$.
$C_{\text{tot}} = 0,050 = [\ce{AH}] + [\ce{A-}] = 2[\ce{AH}] \Rightarrow [\ce{AH}] = [\ce{A-}] = \SI{0,025}{\mol\per\L}$.
Quantités pour \SI{0,500}{\L} : $n = 0,025 \times 0,500 = \SI{0,0125}{\mol}$ chaque.
$m(\ce{NaH2PO4}) = 0,0125 \times 120 = \mathbf{\SI{1,50}{\gram}}$.
$m(\ce{Na2HPO4}) = 0,0125 \times 142 = \mathbf{\SI{1,78}{\gram}}$.
Peser ces masses, dissoudre dans \SI{\approx 400}{\mL} eau, transférer en fiole \SI{500}{\mL}, compléter, homogénéiser.
\end{corrige}

\coursec{Bilan synthétique et checklist Bac}

\begin{bilan}[L'essentiel de la leçon : Dosages acido-basiques et solutions tampons]

\textbf{1. Réactions de dosage et équivalence}
\begin{itemize}
  \item \cle{Dosage} : détermination de concentration par réaction \textbf{rapide, totale, unique}.
  \item \cle{Équivalence} : mélanges en proportions stœchiométriques exactes. Relation : $C_A V_A = C_B V_{BE}$ (adapter coefficients).
  \item \textbf{4 Équations-bilan types} (à connaître par cœur) :
  \begin{align*}
    \text{Fort/Fort} &: \ce{H3O+ + HO- -> 2 H2O} \\
    \text{Faible/Forte} &: \ce{AH + HO- -> A- + H2O} \\
    \text{Fort/Faible} &: \ce{B + H3O+ -> BH+ + H2O} \\
    \text{Faible/Faible} &: \ce{AH + B <=> A- + BH+} \quad (K = 10^{\Delta pK_a})
  \end{align*}
\end{itemize}

\textbf{2. Dosage pH-métrique}
\begin{itemize}
  \item \textbf{Fort/Fort} : pH$_E = 7$, saut immense (3-11), courbe symétrique si $C_A=C_B$.
  \item \textbf{Faible/Forte} : Zone tampon avant $E$, \textbf{pH$_{1/2E} = pK_a$}, saut réduit (pH basique), \textbf{pH$_E > 7$} (hydrolyse \ce{A-}).
  \item \textbf{Forte/Faible} : Symétrique, \textbf{pH$_{1/2E} = pK_a(\ce{BH+})$}, \textbf{pH$_E < 7$} (hydrolyse \ce{BH+}).
  \item \textbf{Repérage $E$} : Méthode des tangentes (2 tangentes parallèles + médiatrice) OU maximum de la dérivée $dpH/dV$.
\end{itemize}

\textbf{3. Dosage colorimétrique}
\begin{itemize}
  \item Indicateur = couple \ce{HIn}/\ce{In-} coloré. Zone de virage = $pK_a \pm 1$.
  \item \textbf{Règle d'or :} Zone de virage doit \textbf{contenir pH$_E$}.
  \item \textbf{B.T.B.} (6-8) pour Fort/Fort ; \textbf{Phénolphtaléine} (8,3-10) pour Faible/Forte ; \textbf{Hélianthine} (0,5-2,5) pour Forte/Faible.
\end{itemize}

\textbf{4. Solutions tampons}
\begin{itemize}
  \item Mélange \ce{AH}/\ce{A-} (concentrations voisines). Résiste à l'ajout d'acide/base/dilution.
  \item $\mathrm{pH} = pK_a + \log\frac{[\ce{A-}]}{[\ce{AH}]}$ (Henderson-Hasselbalch).
  \item Capacité tampon maxi si pH $= pK_a$ (i.e. $[\ce{AH}] = [\ce{A-}]$) et concentration totale forte.
  \item \textbf{Préparation pH $= pK_a$ :} Volumes égaux même concentration (ou dosage à $1/2E$).
  \item \textbf{Préparation pH imposé :} Calculer $R = 10^{\mathrm{pH}-pK_a}$, en déduire volumes/masses.
\end{itemize}

\textbf{Exemples concrets :} Sang (\ce{H2CO3}/\ce{HCO3-}), Bissap (anthocyanes), Cosmétiques, Eaux industrielles (CIMENCAM, SONARA).
\end{bilan}

\begin{aretenir}[Checklist finale avant le Baccalauréat TI]
\begin{itemize}
  \item[$\square$] Je sais écrire l'équation-bilan des \textbf{4 types} de dosage (fort/fort, faible/forte, forte/faible, faible/faible) avec flèches correctes ($\ce{->}$ ou $\ce{<=>}$).
  \item[$\square$] Je définis l'équivalence en une phrase : \emph{« Les réactifs sont mélangés dans les proportions stœchiométriques, ils sont tous deux entièrement consommés »}.
  \item[$\square$] J'applique $C_A V_A = C_B V_{BE}$ (avec coefficients) \textbf{sans erreur d'unités} (L ou mL cohérents).
  \item[$\square$] Je construis un \textbf{tableau d'avancement} pour justifier la composition du milieu à l'équivalence et en déduire le pH.
  \item[$\square$] Je sais que : pH$_E = 7$ (F/F) ; pH$_E > 7$ (acide faible dosé) ; pH$_E < 7$ (base faible dosée).
  \item[$\square$] J'exploite la \textbf{demi-équivalence} : $V_B = V_{BE}/2 \Rightarrow \mathrm{pH} = pK_a$ (acide faible) ou $pK_a(\ce{BH+})$ (base faible).
  \item[$\square$] Je détermine $V_{BE}$ graphiquement : \textbf{méthode des tangentes} (schéma clair) et \textbf{méthode de la dérivée} (maximum).
  \item[$\square$] Je choisis l'\textbf{indicateur coloré} dont la zone de virage ($pK_a \pm 1$) \textbf{contient pH$_E$}. Je cite B.T.B., Phénolphtaléine, Hélianthine avec leurs zones.
  \item[$\square$] Je maîtrise la formule $\mathrm{pH} = pK_a + \log\frac{[\text{base}]}{[\text{acide}]}$ pour les tampons (calcul pH, calcul rapport, préparation).
  \item[$\square$] Je sais préparer un tampon de pH $= pK_a$ (volumes égaux même $C$ ou demi-équivalence) et un tampon de pH quelconque.
  \item[$\square$] Je définis une solution tampon, cite ses propriétés (stabilité pH, capacité tampon) et donne des exemples (sang, industrie, cosmétique).
  \item[$\square$] \textbf{Vérification systématique :} Mon pH final est-il cohérent ? (Acide $<7$, Basique $>7$, Neutre $=7$). Mes concentrations sont-elles plausibles ?
\end{itemize}
\end{aretenir}

\findecours

\end{document}
